% !TeX program = xelatex
% 中国科学技术大学《热学B》总复习资料
% 主文件：只负责文档结构；章节正文统一放入 chapters/ch.tex

\documentclass[UTF8,fontset=fandol,zihao=-4,openany,oneside]{ctexbook}

% preamble.tex
% 版式、字体、颜色、目录、页眉页脚与环境样式。
% 本文件尽量不放课程内容；课程语义命令放在 macros.tex。

\usepackage{iftex}

\ifPDFTeX
  \PackageError{USTC Thermal B Review}{请使用 XeLaTeX 或 LuaLaTeX 编译本文档，不要使用 pdfLaTeX。}{}
\fi

% 页面与排版
\usepackage[
  a4paper,
  margin=2.45cm,
  headheight=15pt,
  headsep=0.7cm,
  footskip=1.0cm
]{geometry}

\usepackage{setspace}
\onehalfspacing
\setlength{\parindent}{2em}
\setlength{\parskip}{0.25em}
\linespread{1.08}

% 数学与单位
\usepackage{amsmath,amssymb,mathtools,bm,mathrsfs}
\usepackage{siunitx}

\sisetup{
  detect-all,
  per-mode=symbol,
  range-phrase=--,
  range-units=single
}

\allowdisplaybreaks[2]
\numberwithin{equation}{chapter}

% 表格、列表、图形
\usepackage{booktabs,array,tabularx,longtable,multirow}
\usepackage{graphicx}
\usepackage{float}
\usepackage{enumitem}

\setlist[itemize]{
  topsep=0.35em,
  itemsep=0.15em,
  parsep=0em,
  leftmargin=2em
}

\setlist[enumerate]{
  topsep=0.35em,
  itemsep=0.15em,
  parsep=0em,
  leftmargin=2.4em
}

% 常用表格列类型
\newcolumntype{Y}{>{\raggedright\arraybackslash}X}
\newcolumntype{C}{>{\centering\arraybackslash}X}

% 颜色与强调框
\usepackage[svgnames,dvipsnames,table]{xcolor}

\definecolor{USTCBlue}{HTML}{005BAC}
\definecolor{USTCDark}{HTML}{1F2933}
\definecolor{USTCLightBlue}{HTML}{EAF4FF}
\definecolor{USTCGold}{HTML}{B08A00}
\definecolor{ReviewGreen}{HTML}{EEF8F0}
\definecolor{WarnRed}{HTML}{B91C1C}
\definecolor{WarnBack}{HTML}{FFF1F2}
\definecolor{FormulaBack}{HTML}{F8FAFC}

\usepackage[most]{tcolorbox}

\tcbset{
  enhanced,
  breakable,
  sharp corners,
  boxrule=0.6pt,
  before skip=0.8em,
  after skip=0.8em,
  left=1em,
  right=1em,
  top=0.75em,
  bottom=0.75em,
  fonttitle=\bfseries
}

\newtcolorbox{KeyBox}[1][]{
  colback=USTCLightBlue,
  colframe=USTCBlue,
  coltitle=white,
  title={#1}
}

\newtcolorbox{FormulaBox}[1][]{
  colback=FormulaBack,
  colframe=USTCDark,
  coltitle=white,
  title={#1}
}

\newtcolorbox{WarnBox}[1][]{
  colback=WarnBack,
  colframe=WarnRed,
  coltitle=white,
  title={#1}
}

\newtcolorbox{MethodBox}[1][]{
  colback=ReviewGreen,
  colframe=ForestGreen,
  coltitle=white,
  title={#1}
}

\newtcolorbox{SourceBox}[1][]{
  colback=white,
  colframe=Gray,
  coltitle=white,
  title={#1}
}

% 超链接与 PDF 元数据
\usepackage{hyperref}
\usepackage{bookmark}

\hypersetup{
  colorlinks=true,
  linkcolor=USTCBlue,
  citecolor=USTCBlue,
  urlcolor=USTCBlue,
  pdftitle={中国科学技术大学 热学B 复习资料},
  pdfauthor={Generated with ChatGPT},
  pdfsubject={USTC Thermal Physics B Review},
  pdfkeywords={热学B, 热力学, 麦克斯韦分布, 玻尔兹曼分布, 输运过程, 相变}
}

% 页眉页脚
\usepackage{fancyhdr}

\pagestyle{fancy}
\fancyhf{}
\fancyhead[L]{\small 中国科学技术大学《热学B》复习}
\fancyhead[R]{\small \leftmark}
\fancyfoot[C]{\small \thepage}

\renewcommand{\headrulewidth}{0.4pt}
\renewcommand{\footrulewidth}{0pt}

\fancypagestyle{plain}{
  \fancyhf{}
  \fancyfoot[C]{\small \thepage}
  \renewcommand{\headrulewidth}{0pt}
}

% CTeX 标题样式
\ctexset{
  chapter={
    name={第,章},
    number=\chinese{chapter},
    format=\Huge\bfseries\color{USTCDark},
    aftername=\quad,
    beforeskip=0pt,
    afterskip=1.2em
  },
  section={
    format=\Large\bfseries\color{USTCDark},
    afterskip=0.6em
  },
  subsection={
    format=\large\bfseries\color{USTCDark},
    afterskip=0.4em
  }
}
% macros.tex
% 课程语义命令、常用公式命令与复习资料专用环境。

% ---------- 文档信息 ----------
\newcommand{\USTC}{中国科学技术大学}
\newcommand{\CourseName}{热学B}
\newcommand{\DocTitle}{\USTC《\CourseName》复习资料}
\newcommand{\DocVersion}{2026.1}

% 作者信息
\newcommand{\DocAuthor}{寒初晓}
\newcommand{\DocAuthorInfo}{中国科学技术大学 \quad 近代物理系}

% ---------- 封面 ----------
\newcommand{\makecoursecover}{%
  \begin{titlepage}
    \thispagestyle{empty}

    \vspace*{1.9cm}

    \begin{flushleft}
      {\color{USTCBlue}\rule{0.20\textwidth}{2.2pt}\par}

      \vspace{1.35cm}

      {\zihao{-4}\color{Gray}\textsf{THERMAL PHYSICS B}\par}

      \vspace{0.55cm}

      {\zihao{0}\bfseries 中国科学技术大学\par}

      \vspace{0.35cm}

      {\zihao{1}\bfseries 《热学B》复习资料\par}
    \end{flushleft}

    \begin{flushleft}
      \renewcommand{\arraystretch}{1.55}
      \begin{tabular}{@{}p{0.20\textwidth}p{0.70\textwidth}@{}}
        \textcolor{Gray}{版本} &
        \DocVersion \\
      \end{tabular}
    \end{flushleft}

    \vfill

    \begin{flushleft}
      {\color{USTCBlue}\rule{\textwidth}{0.7pt}\par}

      \vspace{0.45cm}

      {\zihao{5}\color{Gray}
      \DocAuthor
      \par}

      \vspace{0.15cm}

      {\zihao{5}\color{Gray}
      \DocAuthorInfo
      \par}

      \vspace{0.15cm}

      {\zihao{5}\color{Gray}
      \today
      \par}
    \end{flushleft}

  \end{titlepage}
}

% ---------- 微分、偏导与过程量 ----------
\newcommand{\dd}{\mathop{}\!\mathrm{d}}
\newcommand{\dbar}{\mathop{}\!\bar{\mathrm{d}}}

\newcommand{\pd}[2]{\frac{\partial #1}{\partial #2}}
\newcommand{\pdc}[3]{\left(\frac{\partial #1}{\partial #2}\right)_{#3}}

\newcommand{\od}[2]{\frac{\dd #1}{\dd #2}}
\newcommand{\evalat}[2]{\left.#1\right|_{#2}}

% ---------- 常用热学量 ----------
\newcommand{\kB}{k_{\mathrm B}}
\newcommand{\NA}{N_{\mathrm A}}

\newcommand{\Cv}{C_V}
\newcommand{\Cp}{C_p}
\newcommand{\Cvm}{C_{V,m}}
\newcommand{\Cpm}{C_{p,m}}

\newcommand{\gam}{\gamma}

\newcommand{\avg}[1]{\left\langle #1 \right\rangle}

\newcommand{\vrms}{v_{\mathrm{rms}}}
\newcommand{\vmp}{v_{\mathrm p}}
\newcommand{\vmean}{\bar v}

\newcommand{\eps}{\varepsilon}

% ---------- 复习资料语义环境 ----------
\newenvironment{ReviewList}
  {\begin{enumerate}[label=\textbf{\arabic*.},leftmargin=2.5em]}
  {\end{enumerate}}

\newenvironment{ConditionList}
  {\begin{itemize}[label=\textcolor{USTCBlue}{\ensuremath{\blacktriangleright}},leftmargin=2.2em]}
  {\end{itemize}}

\newenvironment{PitfallList}
  {\begin{itemize}[label=\textcolor{WarnRed}{\ensuremath{\times}},leftmargin=2.2em]}
  {\end{itemize}}

\newcommand{\concept}[1]{\textbf{\color{USTCBlue}#1}}
\newcommand{\keyword}[1]{\textbf{#1}}
\newcommand{\condition}[1]{\textit{适用条件：#1}}
\newcommand{\exam}[1]{\textcolor{WarnRed}{\textbf{考点：}#1}}
\newcommand{\sourcepdf}[1]{\par\smallskip\noindent\textcolor{Gray}{\small 来源参考：#1}}

\newcommand{\notA}{\textbf{注意：本资料按《热学B》范围整理，不按《热学A》额外拔高。}}

% ---------- 课程核心公式快捷命令 ----------

% 热力学基础
\newcommand{\IdealGas}{pV=\nu RT=Nk_{\mathrm B}T}
\newcommand{\VanDerWaals}{\left(p+\dfrac{\nu^2 a}{V^2}\right)(V-\nu b)=\nu RT}

% 第一章、第二章常用
\newcommand{\FirstLaw}{\dd U=\dbar Q+\dbar W}
\newcommand{\QuasiWork}{\dbar W=-p\dd V}
\newcommand{\HeatCapacityDef}{C_L=\left(\dfrac{\dbar Q}{\dd T}\right)_L}
\newcommand{\EnthalpyDef}{H=U+pV}

% 理想气体过程
\newcommand{\AdiabaticIdeal}{
  pV^{\gamma}=\mathrm{const},\quad
  TV^{\gamma-1}=\mathrm{const},\quad
  p^{1-\gamma}T^{\gamma}=\mathrm{const}
}

% 第二定律与熵
\newcommand{\EntropyDef}{\dd S=\dfrac{\dbar Q_{\mathrm{rev}}}{T}}
\newcommand{\ClausiusIneq}{\oint\dfrac{\dbar Q}{T}\le 0}
\newcommand{\BoltzmannEntropy}{S=\kB\ln W}

% 麦克斯韦--玻尔兹曼分布
\newcommand{\PressureKinetic}{p=\dfrac{1}{3}nm\avg{v^2}=\dfrac{2}{3}n\avg{\varepsilon_k}}
\newcommand{\MeanKinetic}{\avg{\varepsilon_k}=\dfrac{3}{2}\kB T}

% 输运过程
\newcommand{\FourierLaw}{j_Q=-\kappa\od{T}{z}}

% 相变与液体表面性质
\newcommand{\Clapeyron}{\od{p}{T}=\dfrac{l}{T(v_2-v_1)}}
\newcommand{\CapillaryRise}{h=\dfrac{2\sigma\cos\theta}{\rho g r}}

% ---------- 公式卡片快捷结构 ----------
\newcommand{\FormulaCard}[4]{%
  \begin{FormulaBox}[#1]
  \begin{align*}
    #2
  \end{align*}
  \condition{#3}
  \par\smallskip
  \textbf{复习提示：}#4
  \end{FormulaBox}
}

\newcommand{\Pitfall}[2]{%
  \begin{WarnBox}[#1]
  #2
  \end{WarnBox}
}

\newcommand{\Method}[2]{%
  \begin{MethodBox}[#1]
  #2
  \end{MethodBox}
}

% ---------- 占位命令：方便后续批量写章节 ----------
\newcommand{\chaptergoal}[1]{%
  \begin{KeyBox}[本章目标]
  #1
  \end{KeyBox}
}

\newcommand{\exammap}[1]{%
  \begin{MethodBox}[考点地图]
  #1
  \end{MethodBox}
}

\newcommand{\formulatable}[1]{%
  \begin{FormulaBox}[公式速查]
  #1
  \end{FormulaBox}
}

% ---------- 第二章：热力学第一定律常用命令 ----------

% 功的两种符号习惯：本资料默认 W 为外界对系统做功
\newcommand{\WorkOnSystem}{\dbar W=-p\dd V}
\newcommand{\WorkBySystem}{\dbar W_{\mathrm{by}}=p\dd V}

% 第一定律与焓
\newcommand{\FirstLawFull}{\dd U=\dbar Q+\dbar W=\dbar Q-p\dd V}
\newcommand{\Mayer}{C_{p,m}-C_{V,m}=R}

% 理想气体内能、焓
\newcommand{\IdealDU}{\dd U=\nu C_{V,m}\dd T}
\newcommand{\IdealDH}{\dd H=\nu C_{p,m}\dd T}

% 理想气体等温、绝热、多方过程
\newcommand{\IsoThermalWorkOn}{W=-\nu RT\ln\frac{V_2}{V_1}}
\newcommand{\AdiabaticWorkOn}{W=\Delta U=\nu C_{V,m}(T_2-T_1)}
\newcommand{\PolyProcess}{pV^n=\mathrm{const}}
\newcommand{\PolyHeatCapacity}{C_{n,m}=C_{V,m}\frac{\gamma-n}{1-n}}

% 循环、热机、制冷机
\newcommand{\EngineEff}{\eta=\frac{W'}{Q_1}=1-\frac{Q_2}{Q_1}}
\newcommand{\CarnotEff}{\eta_{\mathrm C}=1-\frac{T_2}{T_1}}
\newcommand{\COP}{\varepsilon=\frac{Q_2}{W}=\frac{Q_2}{Q_1-Q_2}}
\newcommand{\CarnotCOP}{\varepsilon_{\mathrm C}=\frac{T_2}{T_1-T_2}}

% 焦耳--汤姆孙系数
\newcommand{\JTcoef}{\alpha_i=\left(\frac{\partial T}{\partial p}\right)_H}

% ---------- 第三章：热力学第二定律与熵常用命令 ----------

% 第二定律、克劳修斯关系与熵
\newcommand{\ClausiusEq}{\oint_{\mathrm{rev}}\frac{\dbar Q}{T}=0}
\newcommand{\ClausiusGeneral}{\oint\frac{\dbar Q}{T}\le 0}
\newcommand{\EntropyIntegral}{\Delta S=S_b-S_a=\int_a^b\frac{\dbar Q_{\mathrm{rev}}}{T}}
\newcommand{\EntropyIneq}{\dd S\ge \frac{\dbar Q}{T}}
\newcommand{\IsolatedEntropy}{\Delta S_{\mathrm{isolated}}\ge 0}

% 理想气体熵变
\newcommand{\IdealEntropyTV}{\Delta S=\nu C_{V,m}\ln\frac{T_2}{T_1}+\nu R\ln\frac{V_2}{V_1}}
\newcommand{\IdealEntropyTp}{\Delta S=\nu C_{p,m}\ln\frac{T_2}{T_1}-\nu R\ln\frac{p_2}{p_1}}
\newcommand{\IdealEntropypV}{\Delta S=\nu C_{V,m}\ln\frac{p_2}{p_1}+\nu C_{p,m}\ln\frac{V_2}{V_1}}

% 卡诺热机、卡诺制冷机
\newcommand{\CarnotHeatRatio}{\frac{Q_2}{Q_1}=\frac{T_2}{T_1}}
\newcommand{\TSarea}{W'=\oint T\dd S}

% ---------- 第四章：麦克斯韦--玻尔兹曼分布律常用命令 ----------

% 麦克斯韦速度与速率分布
\newcommand{\MaxwellComponent}{
  f(v_x)=\left(\frac{m}{2\pi k_{\mathrm B}T}\right)^{1/2}
  \exp\left(-\frac{mv_x^2}{2k_{\mathrm B}T}\right)
}

\newcommand{\MaxwellVelocity}{
  f(v_x,v_y,v_z)=
  \left(\frac{m}{2\pi k_{\mathrm B}T}\right)^{3/2}
  \exp\left[-\frac{m(v_x^2+v_y^2+v_z^2)}{2k_{\mathrm B}T}\right]
}

\newcommand{\MaxwellSpeed}{
  f(v)=4\pi
  \left(\frac{m}{2\pi k_{\mathrm B}T}\right)^{3/2}
  v^2
  \exp\left(-\frac{mv^2}{2k_{\mathrm B}T}\right)
}

\newcommand{\MostProbSpeed}{
  v_{\mathrm p}=\sqrt{\frac{2k_{\mathrm B}T}{m}}
  =\sqrt{\frac{2RT}{\mu}}
}

\newcommand{\MeanSpeed}{
  \bar v=\sqrt{\frac{8k_{\mathrm B}T}{\pi m}}
  =\sqrt{\frac{8RT}{\pi\mu}}
}

\newcommand{\RmsSpeed}{
  v_{\mathrm{rms}}=\sqrt{\overline{v^2}}
  =\sqrt{\frac{3k_{\mathrm B}T}{m}}
  =\sqrt{\frac{3RT}{\mu}}
}

% 碰壁数与泻流
\newcommand{\CollisionRate}{
  G=\frac{1}{4}n\bar v
}

\newcommand{\EffusionRate}{
  \frac{\Delta N}{\Delta t}=\frac{1}{4}n\bar v A
}

% 玻尔兹曼分布
\newcommand{\BoltzmannDist}{
  n(\bm r)=n_0\exp\left[-\frac{\varepsilon_p(\bm r)}{k_{\mathrm B}T}\right]
}

\newcommand{\BarometricFormula}{
  n(z)=n_0\exp\left(-\frac{mgz}{k_{\mathrm B}T}\right),
  \qquad
  p(z)=p_0\exp\left(-\frac{\mu gz}{RT}\right)
}

% 能量均分
\newcommand{\Equipartition}{
  \overline{\varepsilon_i}=\frac{1}{2}k_{\mathrm B}T
}

\newcommand{\MolecularEnergy}{
  \bar\varepsilon=\frac{i}{2}k_{\mathrm B}T
}

\newcommand{\IdealCvFromFreedom}{
  C_{V,m}=\frac{i}{2}R
}

% ---------- 第五章：气体中的输运过程常用命令 ----------

% 宏观输运规律
\newcommand{\FourierTransport}{
  \frac{\Delta Q}{\Delta t}
  =
  -\kappa\frac{\dd T}{\dd z}\Delta S
}

\newcommand{\NewtonViscosity}{
  f
  =
  -\eta\frac{\dd u}{\dd z}\Delta S
}

\newcommand{\FickLaw}{
  \frac{\Delta N}{\Delta t}
  =
  -D\frac{\dd n}{\dd z}\Delta S
}

% 输运系数的微观表达
\newcommand{\ThermalConductivity}{
  \kappa=\frac13\lambda\rho c_V\bar v
}

\newcommand{\ViscosityCoeff}{
  \eta=\frac13\lambda\rho\bar v
}

\newcommand{\DiffusionCoeff}{
  D=\frac13\lambda\bar v
}

% 平均自由程与碰撞频率
\newcommand{\MeanFreePath}{
  \lambda=\frac{1}{\sqrt{2}\pi d^2 n}
}

\newcommand{\CollisionFreq}{
  Z=\frac{\bar v}{\lambda}
  =
  \sqrt{2}\pi d^2n\bar v
}

% 输运系数与压强、温度的近似关系
\newcommand{\KappaPressureIndep}{
  \kappa\sim \frac13\lambda\rho c_V\bar v
}

\newcommand{\EtaPressureIndep}{
  \eta\sim \frac13\lambda\rho\bar v
}

\newcommand{\DPressureRelation}{
  D\sim \frac13\lambda\bar v
}

% ---------- 第六章：固、液体性质简介与相变常用命令 ----------

% 固体热容量
\newcommand{\DulongPetit}{
  C_{V,m}\approx 3R
}

% 表面张力
\newcommand{\SurfaceForce}{
  f=\sigma l
}

\newcommand{\SurfaceWork}{
  \Delta A=\sigma\Delta S
}

% 弯曲液面附加压强
\newcommand{\LaplaceSphere}{
  \Delta p=\frac{2\sigma}{R}
}

\newcommand{\LaplaceCylinder}{
  \Delta p=\frac{\sigma}{R}
}

% 毛细现象
\newcommand{\CapillaryFormula}{
  h=\frac{2\sigma\cos\theta}{\rho gr}
}

% 克拉珀龙方程
\newcommand{\ClapeyronEq}{
  \frac{\dd p}{\dd T}=\frac{l}{T(v_2-v_1)}
}

% 克劳修斯--克拉珀龙近似
\newcommand{\ClausiusClapeyron}{
  \frac{\dd p}{\dd T}=\frac{lp}{RT^2}
}

\newcommand{\ClausiusClapeyronInt}{
  \ln p=-\frac{l}{RT}+C
}

% ---------- 总结部分常用命令 ----------

\newcommand{\finalnote}[1]{%
  \begin{KeyBox}[最后提醒]
  #1
  \end{KeyBox}
}

\begin{document}

\frontmatter
\makecoursecover

\tableofcontents

\mainmatter

\chapter{热力学基础知识与温度}

\chaptergoal{
本章是整门《热学B》的概念地基。复习重点不在复杂计算，而在准确区分系统、外界、状态、过程、平衡态、准静态过程、温度、温标与状态方程。后续第一定律、第二定律、熵、气体动理论和相变，都会反复使用本章语言。
}

\exammap{
\begin{ReviewList}
  \item \textbf{会定义：}热力学系统、外界、孤立系统、封闭系统、开放系统、绝热系统。
  \item \textbf{会判断：}一个状态是不是平衡态，一个过程能不能近似看作准静态过程。
  \item \textbf{会使用：}热力学第零定律解释温度概念，并区分经验温标、理想气体温标、热力学温标。
  \item \textbf{会列方程：}理想气体状态方程、混合理想气体状态方程、范德瓦耳斯方程、简单固液体状态方程。
  \item \textbf{会处理：}利用 \(\alpha,\beta\) 推出固体或液体在小范围内的状态变化。
\end{ReviewList}
}

\section{热力学系统、状态与过程}

\subsection{热力学系统与外界}

\concept{热力学系统}是指具有明确边界、被我们选作研究对象的宏观物体或物体系，简称系统。系统边界以外所有能与系统发生作用的物体，统称为\concept{外界}或环境。

\begin{KeyBox}[系统与外界的三类相互作用]
\begin{tabularx}{\textwidth}{@{}>{\bfseries}lXl@{}}
相互作用类型 & 物理含义 & 常见外界模型\\
\midrule
热相互作用 & 通过热传导、热辐射等方式交换能量 & 热源、热库\\
机械相互作用 & 通过广义力做功交换能量 & 功源、功库\\
质量相互作用 & 通过交换物质来交换能量与粒子数 & 粒子源、粒子库\\
\end{tabularx}
\end{KeyBox}

根据系统与外界之间的相互作用情况，常见系统类型如下：

\begin{FormulaBox}[系统分类]
\begin{tabularx}{\textwidth}{@{}>{\bfseries}lXX@{}}
系统类型 & 定义 & 复习提醒\\
\midrule
孤立系统 & 与外界无任何相互作用 & 不交换能量，也不交换物质。\\
封闭系统 & 与外界无质量相互作用 & 可交换能量，但不交换物质。\\
开放系统 & 与外界有质量相互作用 & 可交换物质，通常也可交换能量。\\
绝热系统 & 与外界无热相互作用 & 不交换热量，但可能有机械功。\\
\end{tabularx}
\end{FormulaBox}

\begin{WarnBox}[易错点：封闭系统不等于孤立系统]
封闭系统只是不交换物质，仍然可以通过做功或传热与外界交换能量。孤立系统才是既不交换物质，也不交换能量。绝热系统只排除热相互作用，不排除机械相互作用。
\end{WarnBox}

\subsection{状态、状态参量与态函数}

\concept{系统状态}由系统的宏观性质确定。对简单气体系统，常用状态参量包括压强 \(p\)、体积 \(V\)、温度 \(T\)。这些量也称为热力学参量或热力学坐标。

\begin{KeyBox}[状态参量的独立性]
状态参量是描述平衡态所需的几个相互独立的宏观量。系统越复杂，所需状态参量越多。
\begin{itemize}
  \item 化学纯气体：通常只需两个独立参量，例如 \(p,V\)，第三个由状态方程确定。
  \item 混合气体：除 \(p,V,T\) 关系外，还需描述组分或摩尔数的化学参量。
  \item 电磁场中的系统：还可能需要 \(E,B,P,M\) 等电磁参量。
\end{itemize}
\end{KeyBox}

若系统的某一宏观量只由当前状态决定，而与到达该状态的路径无关，则称其为\concept{态函数}。设在一个 \(p\)--\(V\) 系统中，某态函数为
\[
  S=S(p,V),
\]
则系统经历任意循环回到原状态时，
\[
  \oint \dd S=0.
\]

\begin{WarnBox}[易错点：态函数与过程量]
态函数的变化只看初、末态；过程量依赖路径。第一章先建立这个意识，第二章会专门区分内能 \(U\) 与热量、功。热量和功不能简单写成某个状态函数的全微分。
\end{WarnBox}

\subsection{热力学平衡态}

\concept{热力学平衡态}是指系统的宏观性质不随时间变化，且具有确定值的状态。平衡态不是“微观上完全静止”，而是宏观量稳定的动态平衡。

\begin{FormulaBox}[平衡态的必要条件]
\begin{tabularx}{\textwidth}{@{}>{\bfseries}lX@{}}
力学平衡 & 系统内部压强处处相同；不存在有限压强差导致的宏观运动。\\
热平衡 & 系统各部分冷热程度一致；不存在有限温差导致的热流。\\
质量平衡 & 包括化学平衡与相平衡。化学反应正逆速率达到平衡；多相系统中各相物质保持不变。\\
\end{tabularx}
\end{FormulaBox}

\concept{非平衡态}是不同时满足上述条件的状态。自然界中严格平衡态是理想化概念；实际系统常处于非平衡态，或只是近似处于平衡态。

\begin{MethodBox}[判断平衡态的思路]
\begin{ReviewList}
  \item 看宏观量是否随时间变化。
  \item 看系统内部是否仍存在有限的压强差、温度差、浓度差或相变趋势。
  \item 注意“稳定导热”这类情况：温度分布虽然不随时间变，但若存在持续热流，系统整体一般不是热力学平衡态。
\end{ReviewList}
\end{MethodBox}

\subsection{热力学过程与准静态过程}

\concept{热力学过程}是指系统状态发生变化时所经历的过程。系统从一个平衡态变化到另一个平衡态需要一定时间，这个时间尺度与系统内部恢复平衡的\concept{弛豫时间}有关。

\concept{准静态过程}是由一系列无限接近平衡态的状态组成的理想过程。换句话说，系统在变化过程中始终几乎处于平衡态。

\begin{KeyBox}[准静态过程的物理图像]
若气缸上方放一个大砝码，突然拿走砝码，活塞快速上升，系统经历明显非平衡过程；若把砝码换成一堆细沙，每次只拿走一粒沙子，系统每一步只偏离平衡态极小，则整个过程可近似看成准静态过程。
\end{KeyBox}

准静态过程是一种理想化过程。实际中，若过程进行得足够缓慢，并且系统内部弛豫足够快，就可以把它近似处理为准静态过程。

\begin{FormulaBox}[准静态过程与热力学坐标图]
准静态过程可以在热力学坐标图中用一条曲线表示，例如在 \(p\)--\(V\) 图中：
\[
  p=p(V).
\]
常见过程包括：
\[
  \text{等压过程：}p=\mathrm{const},\qquad
  \text{等容过程：}V=\mathrm{const},\qquad
  \text{等温过程：}T=\mathrm{const}.
\]
\end{FormulaBox}

\begin{WarnBox}[易错点：不是所有过程都能画成 \(p\)--\(V\) 曲线]
只有当系统在过程中可以用确定的平衡态参量描述时，才可以在热力学坐标系中画出过程曲线。非准静态过程经过的是连续变化的非平衡态，系统内部压强、温度等可能没有全局确定值，因此不能严格画成一条热力学过程曲线。
\end{WarnBox}

\section{热力学第零定律与温度}

\subsection{热力学第零定律}

\concept{热力学第零定律}可以表述为：若系统 \(A\) 和系统 \(B\) 分别与系统 \(C\) 处于热平衡，则在不受外界影响的情况下，系统 \(A\) 与系统 \(B\) 也彼此处于热平衡。

\begin{KeyBox}[第零定律的核心意义]
\begin{ReviewList}
  \item 它说明互为热平衡的物体之间存在一个共同特征，这个共同特征就是温度。
  \item 它给出了比较温度是否相同的方法：用同一个标准物体分别与两个系统接触，即温度计思想。
  \item 它只能判断是否热平衡，不能判断热量自发传递方向；方向性问题属于第二定律。
\end{ReviewList}
\end{KeyBox}

温度是对物体冷热程度的量度。更严格地说，温度是表征热平衡状态的状态函数。

\subsection{温度作为共同状态函数}

设系统 \(A,B,C\) 的平衡态各由两个状态参量表示：
\[
  A:(x_A,y_A),\qquad B:(x_B,y_B),\qquad C:(x_C,y_C).
\]
当 \(A\) 与 \(C\) 热平衡、\(B\) 与 \(C\) 热平衡时，有
\[
  f_{AC}(x_A,y_A;x_C,y_C)=0,\qquad
  f_{BC}(x_B,y_B;x_C,y_C)=0.
\]
由第零定律，\(A\) 与 \(B\) 也处于热平衡：
\[
  f_{AB}(x_A,y_A;x_B,y_B)=0.
\]
因此可以引入一个共同的状态函数：
\[
  \psi_A(x_A,y_A)=\psi_B(x_B,y_B)=\psi_C(x_C,y_C).
\]
这个共同状态函数就称为温度，记为 \(T\)。

\begin{WarnBox}[易错点：温度不是热量]
温度描述系统冷热程度，是状态参量；热量是由于温差而传递的能量，是过程量。不能说“物体含有很多热量”，更准确的说法是“物体温度较高”或“在某过程中吸收了热量”。
\end{WarnBox}

\subsection{温标}

\concept{温标}是具体给出温度数值的表示方法。温度概念来自第零定律，但要进行数值测量，必须规定温标。

\subsubsection{经验温标}

经验温标利用某种测温物质的某个物理属性随冷热程度单调变化的性质来定义温度。若测温属性为 \(x\)，常假定
\[
  T(x)=a x.
\]

\begin{FormulaBox}[经验温标的三个要素]
\begin{tabularx}{\textwidth}{@{}>{\bfseries}lX@{}}
测温物质与测温属性 & 例如水银体积、气体压强、金属电阻等。\\
固定点 & 例如冰的正常熔点、水的正常沸点。\\
分度方法 & 规定固定点之间如何等分，例如摄氏温标把 \(0^\circ\mathrm C\) 到 \(100^\circ\mathrm C\) 等分为100份。\\
\end{tabularx}
\end{FormulaBox}

若冰点时温度为 \(\theta_i\)、测量值为 \(x_i\)，汽点时温度为 \(\theta_s\)、测量值为 \(x_s\)，假设线性关系，则某一测量值 \(x\) 对应温度为
\[
  \theta
  =
  \theta_i+
  \frac{x-x_i}{x_s-x_i}(\theta_s-\theta_i).
\]

\begin{WarnBox}[经验温标的局限]
不同测温物质或不同测温属性建立的经验温标不会严格一致。这也是为什么需要理想气体温标和热力学温标。
\end{WarnBox}

\subsubsection{理想气体温标}

\concept{理想气体温标}以气体为测温物质，利用理想气体状态方程中定容时压强与温度成正比、或定压时体积与温度成正比的关系来定义温度。

\begin{FormulaBox}[气体温度计]
定容气体温度计：
\[
  T(p)=273.16\,\frac{p}{p_0}.
\]
其中 \(p_0\) 为气体在水的三相点温度时的压强。

定压气体温度计：
\[
  T(V)=273.16\,\frac{V}{V_0}.
\]
其中 \(V_0\) 为气体在水的三相点温度时的体积。
\end{FormulaBox}

这里 \(273.16\,\mathrm K\) 是水的三相点温度。气体温度计适用范围较宽，低温气体温度计常以氦气为测温物质。

\subsubsection{摄氏温标、华氏温标与国际实用温标}

摄氏温标 \(t\) 以冰的正常熔点为 \(0^\circ\mathrm C\)，水的正常沸点为 \(100^\circ\mathrm C\)。华氏温标 \(t_F\) 与摄氏温标的关系为
\[
  t_F=\frac{9}{5}t+32.
\]

国际实用温标中，热力学温度 \(T\) 与摄氏温度 \(t\) 的关系为
\[
  T=\left(\frac{t}{^\circ\mathrm C}+273.15\right)\mathrm K.
\]
也常写作
\[
  T/\mathrm K=t/^\circ\mathrm C+273.15.
\]

\begin{FormulaBox}[常见温度点]
\begin{tabularx}{\textwidth}{@{}lccc@{}}
\toprule
物理状态 & \(T/\mathrm K\) & \(t/^\circ\mathrm C\) & \(t_F/^\circ\mathrm F\)\\
\midrule
水的汽点 & \(373.15\) & \(100.00\) & \(212.00\)\\
水的三相点 & \(273.16\) & \(0.01\) & \(32.02\)\\
水的冰点 & \(273.15\) & \(0.00\) & \(32.00\)\\
绝对零度 & \(0.00\) & \(-273.15\) & \(-459.67\)\\
\bottomrule
\end{tabularx}
\end{FormulaBox}

\subsubsection{热力学温标}

由于经验温标依赖测温物质和测温属性，需要引入一种不依赖具体测温物质的温标，即\concept{热力学温标}，也称绝对温标。国际上规定热力学温标为基本温标，一切温度测量最终以热力学温标为准。

在理想气体温标可适用的范围内，热力学温标与理想气体温标一致。

\begin{WarnBox}[热学B复习边界]
实用温度计，如液体温度计、电阻温度计、热电偶温度计、光学高温计、低温温度计，主要作为了解内容。复习重点是温标思想、温标之间的关系，以及水的三相点 \(273.16\,\mathrm K\)。
\end{WarnBox}

\section{状态方程}

\subsection{状态方程的概念}

\concept{状态方程}是处于平衡态的热力学系统，其状态参量之间满足的函数关系。对化学纯气体，可写为
\[
  T=T(p,V),
\]
或
\[
  f(T,p,V)=0.
\]

状态方程可以由实验给出，也可以由微观理论推导。热学B中重点掌握三类状态方程：理想气体、范德瓦耳斯气体、简单固体和液体。

\subsection{理想气体状态方程}

\concept{理想气体}是压强趋于零极限状态下的气体。它是一种理想模型，实际不存在完全理想的气体；但当温度不太低、压强不太大时，许多真实气体可以近似看作理想气体。

\begin{KeyBox}[理想气体模型]
\begin{ReviewList}
  \item 分子本身的体积相对分子间距离可以忽略。
  \item 除碰撞瞬间外，分子间相互作用力可以忽略。
  \item 分子之间、分子与器壁之间的碰撞视为弹性碰撞。
  \item 从能量角度看，忽略分子间势能，理想气体内能只与温度有关。
\end{ReviewList}
\end{KeyBox}

\subsubsection{气体实验定律}

\begin{FormulaBox}[三条气体实验定律]
\begin{tabularx}{\textwidth}{@{}>{\bfseries}lXX@{}}
玻意耳--马略特定律 & \(T=\mathrm{const}\) & \(pV=C\)\\
查理定律 & \(V=\mathrm{const}\) & \(p/T=C\)\\
盖--吕萨克定律 & \(p=\mathrm{const}\) & \(V/T=C\)\\
\end{tabularx}
\end{FormulaBox}

由这些实验定律可得，对一定质量理想气体，
\[
  \frac{pV}{T}=C.
\]
由阿伏伽德罗定律，常数与摩尔数 \(\nu\) 成正比，得到理想气体状态方程：
\[
  \boxed{pV=\nu RT}.
\]

若用分子数 \(N=\nu N_A\)，且 \(R=N_A k\)，则
\[
  \boxed{pV=NkT}.
\]

\begin{FormulaBox}[常数]
\[
  R=8.31\,\mathrm{J\cdot mol^{-1}\cdot K^{-1}},
  \qquad
  N_A=6.02\times 10^{23}\,\mathrm{mol^{-1}},
\]
\[
  k=\frac{R}{N_A}
  \approx 1.38\times 10^{-23}\,\mathrm{J\cdot K^{-1}}.
\]
\end{FormulaBox}

\begin{WarnBox}[易错点：\(\nu\)、\(N\)、\(N_A\)]
\(\nu\) 是摩尔数，\(N\) 是分子数，\(N_A\) 是阿伏伽德罗常数。公式
\[
  pV=\nu RT
\]
与
\[
  pV=NkT
\]
等价，但不能把 \(\nu\) 和 \(N\) 混用。
\end{WarnBox}

\subsection{混合理想气体状态方程}

对于混合理想气体，遵循\concept{道尔顿分压定律}：
\[
  p=\sum_i p_i.
\]
第 \(i\) 种气体满足
\[
  p_iV=\nu_i RT.
\]
各组分相加得
\[
  \boxed{pV=\nu RT},
  \qquad
  \nu=\sum_i\nu_i.
\]

若混合气体总质量为 \(M\)，平均摩尔质量为 \(\bar m\)，则
\[
  \nu=\frac{M}{\bar m},
  \qquad
  pV=\frac{M}{\bar m}RT.
\]

\begin{MethodBox}[混合理想气体题的做法]
\begin{ReviewList}
  \item 若题目问总压强、总体积或总摩尔数，直接对混合气体整体使用 \(pV=\nu RT\)。
  \item 若题目问某组分，例如氧气质量或某组分分压，要先确定该组分所占的摩尔分数、质量分数或分压比例。
  \item 若温度和体积不变，压强比通常等于物质的量比。
\end{ReviewList}
\end{MethodBox}

\subsection{理想气体状态方程的典型应用}

\subsubsection{抽气过程}

容积为 \(V\) 的密封容器，抽气机抽气速率为 \(u\)。若温度恒定，气压由 \(p_0\) 降到 \(p\)，所需时间为
\[
  \boxed{t=\frac{V}{u}\ln\frac{p_0}{p}}.
\]

推导思路如下：在 \(t\) 到 \(t+\dd t\) 内，压强由 \(p\) 变为 \(p+\dd p\)，排出气体体积为 \(u\dd t\)。等温下有
\[
  pV=(p+\dd p)(V+u\dd t).
\]
忽略高阶小量：
\[
  V\dd p+p\,u\dd t=0,
\]
即
\[
  \frac{\dd p}{p}=-\frac{u}{V}\dd t.
\]
积分得到上式。

\begin{WarnBox}[抽气题常见错法]
不要把 \(p\) 当常量。抽气过程中容器内压强随时间指数下降，微分方程的核心是
\[
  \frac{\dd p}{p}=-\frac{u}{V}\dd t.
\]
\end{WarnBox}

\subsubsection{高空吸氧质量变化}

若空气温度及组分比例不随高度变化，人在相同肺容积下每次吸入氧气质量与空气压强成正比：
\[
  \frac{m}{m_0}=\frac{p}{p_0}.
\]
例如当
\[
  p=5.0\times 10^4\,\mathrm{Pa},
  \qquad
  p_0=1.01\times 10^5\,\mathrm{Pa},
\]
标准状况下一次吸入氧气 \(m_0=1.0\,\mathrm g\)，则
\[
  m=\frac{p}{p_0}m_0
  =
  \frac{5.0\times 10^4}{1.01\times 10^5}\times 1.0\,\mathrm g
  \approx 0.495\,\mathrm g.
\]

\section{实际气体状态方程}

\subsection{范德瓦耳斯方程}

理想气体模型忽略了两件事：
\begin{ReviewList}
  \item 分子具有固有体积。
  \item 分子间存在除碰撞外的相互作用力。
\end{ReviewList}

范德瓦耳斯方程正是对这两点进行修正。

\subsubsection{体积修正}

若 \(\nu\) mol 气体占据体积 \(V\)，考虑分子固有体积后，分子可自由活动的有效体积不再是 \(V\)，而近似为
\[
  V-\nu b,
\]
其中 \(b\) 是与气体种类有关的常数。仅考虑体积修正时，
\[
  p(V-\nu b)=\nu RT.
\]

\subsubsection{压强修正}

实际气体分子间存在吸引力。靠近器壁的分子受到指向气体内部的吸引力，撞击器壁时给器壁的冲量小于理想气体情形，因此实际测得压强 \(p\) 小于气体内部对应的动理压强。设压强修正量为
\[
  \Delta p=\frac{\nu^2 a}{V^2},
\]
则
\[
  p+\frac{\nu^2a}{V^2}
\]
对应被修正后的内部动理压强。

因此 \(\nu\) mol 真实气体的范德瓦耳斯方程为
\[
  \boxed{
  \left(p+\frac{\nu^2a}{V^2}\right)(V-\nu b)=\nu RT
  }.
\]

对 \(1\) mol 气体，
\[
  \boxed{
  \left(p+\frac{a}{V_m^2}\right)(V_m-b)=RT
  }.
\]
这里 \(V_m\) 表示摩尔体积。

\begin{FormulaBox}[范德瓦耳斯参数的物理意义]
\begin{tabularx}{\textwidth}{@{}>{\bfseries}lX@{}}
\(a\) & 描述分子间吸引作用。吸引力越明显，压强修正越大。\\
\(b\) & 描述分子固有体积造成的有效体积修正。\\
\end{tabularx}
\end{FormulaBox}

\begin{WarnBox}[易错点：内压强不是气体内部压强]
授课中“内压强”指由于分子吸引力造成的压强修正 \(\Delta p\)，不是简单意义上的“气体内部的压强”。公式里要写
\[
  p+\Delta p=p+\frac{\nu^2a}{V^2}.
\]
\end{WarnBox}

\subsection{范德瓦耳斯方程的适用性}

范德瓦耳斯方程比理想气体方程更接近真实气体，但仍是近似方程。一般在压强不是很高、温度不是太低的情况下，对许多真实气体是较好的近似。

\begin{KeyBox}[为什么范德瓦耳斯方程重要]
它的价值不只是“比理想气体更准”，更重要的是物理图像清楚：通过 \(a\) 和 \(b\) 两个修正量，同时体现分子吸引与分子固有体积，并能初步解释气液转变和临界现象。
\end{KeyBox}

\subsection{昂内斯方程与维里展开}

另一类真实气体状态方程是昂内斯方程，也称维里展开。按体积展开可写为
\[
  pV=A+\frac{B_V}{V}+\frac{C_V}{V^2}+\cdots.
\]
按压强展开可写为
\[
  pV=A+Bp+Cp^2+Dp^3+\cdots.
\]
其中 \(A,B,C,\ldots\) 或 \(A,B_V,C_V,\ldots\) 是温度的函数，并与气体种类有关，称为维里系数。

理想气体是一级近似：
\[
  A=\nu RT.
\]

范德瓦耳斯方程也可展开为维里形式。当 \(b/V\ll 1\) 时，
\[
  pV
  =
  \nu RT+
  \frac{\nu^2RTb-\nu^2a}{V}
  +\cdots.
\]
因此前两项维里系数为
\[
  A_V=\nu RT,\qquad
  B_V=\nu^2RTb-\nu^2a.
\]

\begin{WarnBox}[维里系数的物理符号]
\(B_V\) 中 \(+\nu^2RTb\) 来自排斥或体积效应，贡献为正；\(-\nu^2a\) 来自吸引效应，贡献为负。这个符号关系经常被用来判断真实气体相对于理想气体的偏离方向。
\end{WarnBox}

\section{简单固体与液体的状态方程}

\subsection{等压体膨胀系数与等温压缩系数}

简单固体与液体的状态方程也可写为
\[
  f(p,V,T)=0.
\]
若选择 \(p,T\) 为独立状态参量，则
\[
  V=V(p,T).
\]
其全微分为
\[
  \dd V
  =
  \pdc{V}{T}{p}\dd T+
  \pdc{V}{p}{T}\dd p.
\]

定义\concept{等压体膨胀系数}
\[
  \boxed{
  \alpha
  =
  \frac{1}{V}\pdc{V}{T}{p}
  }
\]
和\concept{等温压缩系数}
\[
  \boxed{
  \beta
  =
  -\frac{1}{V}\pdc{V}{p}{T}
  }.
\]

因此
\[
  \boxed{
  \dd V=\alpha V\dd T-\beta V\dd p
  }.
\]

\begin{KeyBox}[\(\alpha,\beta\) 的物理意义]
\begin{itemize}
  \item \(\alpha\)：压强不变时，温度每变化 \(1\,\mathrm K\) 引起的相对体积变化，反映物体热胀冷缩的容易程度。
  \item \(\beta\)：温度不变时，压强每变化一个单位引起的相对体积变化的负值，反映物体被压缩的容易程度。
\end{itemize}
\end{KeyBox}

\subsection{小范围近似状态方程}

若在一定温度、压强范围内，\(\alpha,\beta\) 都可近似看作常数，初态为 \((p_0,V_0,T_0)\)，则由
\[
  \frac{\dd V}{V}=\alpha\dd T-\beta\dd p
\]
在小变化近似下可得
\[
  \boxed{
  V
  =
  V_0\left[
  1+\alpha(T-T_0)-\beta(p-p_0)
  \right]
  }.
\]

\begin{WarnBox}[易错点：\(\beta\) 前面有负号]
压强增大时，大多数物体体积减小，所以
\[
  \pdc{V}{p}{T}<0.
\]
为了让压缩系数 \(\beta\) 通常为正，定义中必须加负号：
\[
  \beta=-\frac{1}{V}\pdc{V}{p}{T}.
\]
\end{WarnBox}

\subsection{保持体积不变时的压强变化}

由
\[
  \dd V=\alpha V\dd T-\beta V\dd p
\]
若保持体积不变，即 \(\dd V=0\)，则
\[
  \boxed{
  \dd p=\frac{\alpha}{\beta}\dd T
  }.
\]
由于固体和液体通常 \(\alpha/\beta\) 很大，所以哪怕温度变化不大，若要保持体积不变，也可能需要很大的压强变化。

\subsection{三变量循环关系}

若系统满足
\[
  p=p(V,T),
\]
则
\[
  \dd p=
  \pdc{p}{V}{T}\dd V+
  \pdc{p}{T}{V}\dd T.
\]
在 \(p=\mathrm{const}\) 时，
\[
  \pdc{p}{V}{T}\dd V+
  \pdc{p}{T}{V}\dd T=0.
\]
因此得到三变量循环关系：
\[
  \boxed{
  \pdc{p}{V}{T}
  \pdc{V}{T}{p}
  \pdc{T}{p}{V}
  =
  -1
  }.
\]

这个关系不是热学特有公式，而是三个变量之间存在函数关系时的全微分恒等式。热学中常用它在不同偏导之间转换。

\subsection{典型题：水银定容升温}

\begin{MethodBox}[题型模板：定容升温求压强]
已知某液体或固体 \(\alpha,\beta\) 近似为常数，保持体积不变，使温度由 \(T_1\) 升至 \(T_2\)。由
\[
  \dd p=\frac{\alpha}{\beta}\dd T
\]
积分得
\[
  \boxed{
  p_2-p_1=\frac{\alpha}{\beta}(T_2-T_1)
  }.
\]
\end{MethodBox}

例如一团水银在 \(1\) 标准大气压下，温度为 \(0^\circ\mathrm C\)，保持体积不变升温至 \(10^\circ\mathrm C\)。若
\[
  \alpha=1.81\times 10^{-6}\,\mathrm K^{-1},
  \qquad
  \beta=3.82\times 10^{-11}\,\mathrm{Pa}^{-1},
\]
则
\[
  p_2
  =
  p_1+\frac{\alpha}{\beta}(T_2-T_1).
\]
代入
\[
  p_1=1.01\times 10^5\,\mathrm{Pa},
  \qquad
  T_2-T_1=10\,\mathrm K,
\]
得
\[
  p_2
  \approx
  1.01\times 10^5+
  \frac{1.81\times 10^{-6}}{3.82\times 10^{-11}}\times 10
  \approx
  4.75\times 10^5\,\mathrm{Pa}.
\]

\section{第一章公式速查}

\formulatable{
\begin{tabularx}{\textwidth}{@{}>{\bfseries}lXl@{}}
\toprule
内容 & 公式 & 条件\\
\midrule
态函数循环积分 & \(\displaystyle \oint \dd S=0\) & \(S\) 为态函数\\[0.4em]
经验温标线性关系 &
\(\displaystyle \theta=\theta_i+\frac{x-x_i}{x_s-x_i}(\theta_s-\theta_i)\)
& 经验线性假设\\[0.8em]
定容气体温度计 &
\(\displaystyle T=273.16\frac{p}{p_0}\)
& \(V=\mathrm{const}\)\\[0.8em]
定压气体温度计 &
\(\displaystyle T=273.16\frac{V}{V_0}\)
& \(p=\mathrm{const}\)\\[0.8em]
摄氏温度与热力学温度 &
\(\displaystyle T/\mathrm K=t/^\circ\mathrm C+273.15\)
& 国际实用温标\\[0.8em]
摄氏与华氏 &
\(\displaystyle t_F=\frac95 t+32\)
& 温标换算\\[0.8em]
理想气体 &
\(\displaystyle pV=\nu RT=NkT\)
& 理想气体近似\\[0.8em]
混合理想气体 &
\(\displaystyle p=\sum_i p_i,\quad pV=\left(\sum_i\nu_i\right)RT\)
& 道尔顿分压定律\\[0.8em]
范德瓦耳斯方程 &
\(\displaystyle \left(p+\frac{\nu^2a}{V^2}\right)(V-\nu b)=\nu RT\)
& 真实气体近似\\[1em]
固液体微分状态方程 &
\(\displaystyle \dd V=\alpha V\dd T-\beta V\dd p\)
& 简单固体、液体\\[0.8em]
固液体小范围状态方程 &
\(\displaystyle V=V_0[1+\alpha(T-T_0)-\beta(p-p_0)]\)
& \(\alpha,\beta\) 近似常数\\[0.8em]
三变量循环关系 &
\(\displaystyle \pdc{p}{V}{T}\pdc{V}{T}{p}\pdc{T}{p}{V}=-1\)
& 三变量有函数关系\\
\bottomrule
\end{tabularx}
}

\section{第一章易错点总表}

\begin{WarnBox}[本章高频错误]
\begin{PitfallList}
  \item 把封闭系统误认为孤立系统。封闭系统不交换物质，但可以交换能量。
  \item 把绝热系统误认为不做功系统。绝热只是不传热，仍可通过机械方式交换能量。
  \item 认为宏观量不随时间变化就一定是热力学平衡态。若系统内部存在持续热流或浓度流，通常仍不是平衡态。
  \item 把非准静态过程画成 \(p\)--\(V\) 图上的确定曲线。只有准静态过程才可以严格用热力学坐标曲线表示。
  \item 把温度与热量混为一谈。温度是状态参量，热量是过程量。
  \item 忘记水的三相点是 \(273.16\,\mathrm K\)，而水的冰点是 \(273.15\,\mathrm K\)。
  \item 在 \(pV=\nu RT=NkT\) 中混淆摩尔数 \(\nu\) 与分子数 \(N\)。
  \item 写范德瓦耳斯方程时漏掉平方：压强修正是 \(\nu^2a/V^2\)，不是 \(\nu a/V\)。
  \item 写等温压缩系数时漏掉负号：\(\beta=-\frac{1}{V}\pdc{V}{p}{T}\)。
  \item 使用固液体近似状态方程时忘记前提：\(\alpha,\beta\) 只能在小范围内近似为常数。
\end{PitfallList}
\end{WarnBox}

\section{第一章复习路线}

\begin{MethodBox}[建议背诵顺序]
\begin{ReviewList}
  \item 先背系统分类：孤立、封闭、开放、绝热。
  \item 再背平衡态条件：力学平衡、热平衡、质量平衡。
  \item 再理解准静态过程：无限接近平衡态，才能画过程曲线。
  \item 再背第零定律：热平衡的传递性给出温度概念。
  \item 再背温标关系：\(T=t+273.15\)，水三相点 \(273.16\,\mathrm K\)。
  \item 最后整理状态方程：理想气体、混合理想气体、范德瓦耳斯气体、固液体。
\end{ReviewList}
\end{MethodBox}

\begin{KeyBox}[第一章一句话总结]
第一章的核心不是计算技巧，而是建立热学语言：选定系统，确认状态是否平衡，判断过程是否准静态，用温度描述热平衡，用状态方程连接 \(p,V,T\)。后面所有热力学计算都建立在这些前提之上。
\end{KeyBox}
\chapter{热力学第一定律}

\chaptergoal{
本章的核心是把“能量守恒”写成可计算的热学语言。需要清楚区分功、热量、内能、焓、热容这些量的性质：功和热量是过程量，内能和焓是态函数。计算题的主线是：先判断过程，再写第一定律，最后用状态方程与热容关系求 \(Q,W,\Delta U,\Delta H\)。
}

\exammap{
\begin{ReviewList}
  \item \textbf{会算功：}准静态体积功 \(\dbar W=-p\dd V\)，并能从 \(p\)--\(V\) 图判断正负。
  \item \textbf{会区分：}功和热量是过程量；内能、焓是态函数。
  \item \textbf{会用第一定律：}\(\dd U=\dbar Q+\dbar W\)，本资料采用“外界对系统做功”为正的符号约定。
  \item \textbf{会用热容：}掌握 \(C_V,C_p\) 与 \(\Delta U,\Delta H\) 的关系。
  \item \textbf{会处理理想气体：}等容、等压、等温、绝热、多方过程。
  \item \textbf{会判断实际气体节流：}节流过程为等焓过程，焦耳--汤姆孙系数描述温度随压强变化。
  \item \textbf{会算循环：}热机效率、制冷机制冷系数，以及卡诺、奥托、狄塞尔循环的基本结果。
\end{ReviewList}
}

\section{功和热量}

\subsection{改变系统状态的两种方式}

改变热力学系统状态的方式主要有两类：
\begin{ReviewList}
  \item \concept{做功}：通过机械相互作用改变系统能量。
  \item \concept{传热}：通过热相互作用改变系统能量。
\end{ReviewList}

两者都描述能量传递，但物理机制不同。做功通常伴随宏观有序运动，例如活塞移动；传热来自分子无规则热运动能量在温差驱动下的传递。

\begin{WarnBox}[本章符号约定]
本资料采用授课 PDF 中的约定：\(W\) 表示\textbf{外界对系统所做的功}。
因此：
\[
  \dbar W=-p\dd V.
\]
若系统膨胀，\(\dd V>0\)，则外界对系统做负功；若系统被压缩，\(\dd V<0\)，则外界对系统做正功。
\end{WarnBox}

\subsection{准静态过程中的体积功}

考虑气体被活塞封闭，外界压强为 \(p'\)。忽略摩擦时，外界对气体所做元功为
\[
  \dbar W=-p'\dd V.
\]

若过程不是准静态过程，系统内部压强可能不均匀，此时一般有
\[
  p'\ne p,
\]
计算外力对系统做功只能使用外压强 \(p'\)。

若过程为准静态过程，系统始终近似处于平衡态，则
\[
  p'=p,
\]
所以准静态体积功为
\[
  \boxed{\dbar W=-p\dd V}.
\]

对于有限准静态过程，
\[
  \boxed{
  W=-\int_{V_1}^{V_2}p\dd V
  }.
\]

\begin{KeyBox}[\(p\)--\(V\) 图中的功]
在 \(p\)--\(V\) 图上，系统对外做功 \(W_{\mathrm{by}}\) 等于过程曲线下方面积：
\[
  W_{\mathrm{by}}=\int_{V_1}^{V_2}p\dd V.
\]
本资料中外界对系统做功为
\[
  W=-W_{\mathrm{by}}=-\int_{V_1}^{V_2}p\dd V.
\]
因此看图时一定先确认题目使用的是“系统对外做功”还是“外界对系统做功”。
\end{KeyBox}

\subsection{循环过程中的功}

对于一闭合准静态循环，
\[
  \oint \dd U=0
\]
因为内能是态函数。但功不是态函数，一般有
\[
  \oint \dbar W\ne 0.
\]

在 \(p\)--\(V\) 图中：
\begin{itemize}
  \item 顺时针循环：系统对外输出净功，外界对系统做功为负。
  \item 逆时针循环：外界对系统输入净功，外界对系统做功为正。
\end{itemize}

\begin{WarnBox}[易错点：功不是状态函数]
同一初态和末态之间，沿不同路径所作的功一般不同。例如理想气体从 \(C\) 到 \(D\)，走等温路径、先等压后等容路径、先等容后等压路径，功的结果不同。因此不能写 \(W=W_2-W_1\)，也不能写 \(\dd W\) 表示某个态函数的全微分。
\end{WarnBox}

\subsection{几种基本过程的体积功}

\subsubsection{等容过程}

等容过程满足
\[
  V=\mathrm{const},
  \qquad
  \dd V=0.
\]
因此
\[
  \boxed{W=0}.
\]

\subsubsection{等压过程}

等压过程满足
\[
  p=\mathrm{const}.
\]
外界对系统做功为
\[
  \boxed{
  W=-p(V_2-V_1)
  }.
\]
若气体膨胀，\(V_2>V_1\)，则 \(W<0\)。

\subsubsection{理想气体等温过程}

对理想气体等温准静态过程，
\[
  pV=\nu RT,
  \qquad
  T=\mathrm{const}.
\]
因此
\[
  W
  =
  -\int_{V_1}^{V_2}\frac{\nu RT}{V}\dd V
  =
  \boxed{
  -\nu RT\ln\frac{V_2}{V_1}
  }.
\]
也可写为
\[
  W=\nu RT\ln\frac{p_2}{p_1}.
\]

\begin{WarnBox}[等温功的符号]
若等温膨胀，\(V_2>V_1\)，则
\[
  \ln\frac{V_2}{V_1}>0,\qquad W<0.
\]
这表示外界对气体做负功，等价于气体对外做正功。
\end{WarnBox}

\subsection{固体、液体中的体积功}

体积功公式
\[
  W=-\int p\dd V
\]
不仅适用于气体，也适用于简单固体和液体。但固体、液体体积变化通常很小，计算时常借助
\[
  \dd V=\alpha V\dd T-\beta V\dd p.
\]

若质量为 \(M\)、密度为 \(\rho\)、等温压缩系数 \(\beta\) 近似为常数的简单固体，在等温准静态过程中压强从 \(p_i\) 增至 \(p_f\)，则
\[
  \dd T=0,\qquad
  \dd V=-\beta V\dd p\approx-\beta\frac{M}{\rho}\dd p.
\]
于是
\[
  W=-\int p\dd V
  =
  \int_{p_i}^{p_f}p\frac{\beta M}{\rho}\dd p
  =
  \boxed{
  \frac{\beta M}{2\rho}(p_f^2-p_i^2)
  }.
\]

\section{热量与热容}

\subsection{热量的概念}

\concept{热量}是由于温差而在系统和外界之间传递的能量。热量不是系统“含有”的东西，而是一个过程中传递能量的量度。

\begin{KeyBox}[功与热量的相同点和区别]
\begin{tabularx}{\textwidth}{@{}>{\bfseries}lXX@{}}
 & 功 & 热量\\
\midrule
共同点 & \multicolumn{2}{X}{都是能量传递方式，单位均为焦耳 \(\mathrm J\)，都是过程量。}\\
来源 & 机械相互作用 & 热相互作用\\
微观图像 & 分子宏观有序运动对应的能量转移 & 分子无规则热运动能量的传递\\
数学记号 & \(\dbar W\) & \(\dbar Q\)\\
\end{tabularx}
\end{KeyBox}

因为热量是过程量，微小热量不写作 \(\dd Q\)，而写作
\[
  \dbar Q.
\]

\subsection{热容}

若系统在某一无限小过程中吸收热量 \(\dbar Q\)，温度变化为 \(\dd T\)，则该过程的热容定义为
\[
  \boxed{
  C_L=\left(\frac{\dbar Q}{\dd T}\right)_L
  }.
\]
这里下标 \(L\) 表示具体过程。

\begin{FormulaBox}[常用热容]
\[
  C_V=\left(\frac{\dbar Q}{\dd T}\right)_V,
  \qquad
  C_p=\left(\frac{\dbar Q}{\dd T}\right)_p.
\]
摩尔热容：
\[
  C_{V,m}=\frac{C_V}{\nu},
  \qquad
  C_{p,m}=\frac{C_p}{\nu}.
\]
比热容：
\[
  c=\frac{C}{m}.
\]
\end{FormulaBox}

在某一指定过程 \(L\) 中，如果热容近似为常数，则
\[
  Q=\int C_L\dd T=C_L(T_2-T_1).
\]

\begin{WarnBox}[易错点：热容依赖过程]
同一系统的热容不是唯一的。定容热容 \(C_V\) 与定压热容 \(C_p\) 对应不同过程，不能随意替换。更一般地，不同路径 \(L\) 会有不同的过程热容 \(C_L\)。
\end{WarnBox}

\section{内能与热力学第一定律}

\subsection{内能}

\concept{内能} \(U\) 是系统内部所有微观能量的总和，包括分子的热运动能、分子间相互作用势能，以及分子、原子内部能量等。

热力学中通常不关心内能绝对值，而关心两个平衡态之间的内能变化：
\[
  \Delta U=U_2-U_1.
\]
内能是态函数，因此
\[
  \oint\dd U=0.
\]

\begin{WarnBox}[易错点：内能是态函数，不是过程量]
功和热量依赖路径，内能只依赖状态。两个状态之间不管经历什么过程，只要初末态相同，\(\Delta U\) 就相同。
\end{WarnBox}

\subsection{热力学第一定律}

热力学第一定律表达能量守恒：
\[
  \boxed{
  \dd U=\dbar Q+\dbar W
  }.
\]
在只有准静态体积功时，
\[
  \dbar W=-p\dd V,
\]
所以
\[
  \boxed{
  \dd U=\dbar Q-p\dd V
  }.
\]
也可写作
\[
  \boxed{
  \dbar Q=\dd U+p\dd V
  }.
\]

\begin{KeyBox}[第一定律的物理意义]
系统内能的增加，等于系统从外界吸收的热量，加上外界对系统所做的功。第一定律本质上是能量守恒定律在热现象中的表达。
\end{KeyBox}

\begin{WarnBox}[第一类永动机]
第一类永动机指不消耗任何能量而能持续对外做功的机器。它违反热力学第一定律，因此不可能制造。
\end{WarnBox}

\subsection{焓}

定义\concept{焓}
\[
  \boxed{
  H=U+pV
  }.
\]
焓是态函数，因为 \(U,p,V\) 都是状态量。

对简单可压缩系统，
\[
  \dd H=\dd U+p\dd V+V\dd p.
\]
由第一定律
\[
  \dbar Q=\dd U+p\dd V,
\]
可得
\[
  \boxed{
  \dbar Q=\dd H-V\dd p
  }.
\]

因此在等压过程中，
\[
  \dd p=0,
  \qquad
  \boxed{
  \dbar Q_p=\dd H
  }.
\]

\begin{KeyBox}[为什么要引入焓]
定压过程在热学题中很常见。引入焓后，定压吸热可以直接写成
\[
  Q_p=\Delta H.
\]
这比每次都从 \(\Delta U+p\Delta V\) 展开更清楚。
\end{KeyBox}

\subsection{热容与内能、焓}

定容过程中，
\[
  \dd V=0.
\]
由第一定律
\[
  \dbar Q_V=\dd U.
\]
所以
\[
  \boxed{
  C_V=\left(\frac{\partial U}{\partial T}\right)_V
  }.
\]

定压过程中，
\[
  \dbar Q_p=\dd H.
\]
所以
\[
  \boxed{
  C_p=\left(\frac{\partial H}{\partial T}\right)_p
  }.
\]

若热容近似为常数：
\[
  \Delta U=\int_{T_1}^{T_2}C_V\dd T,
  \qquad
  \Delta H=\int_{T_1}^{T_2}C_p\dd T.
\]

\begin{FormulaBox}[热容计算的最基本结论]
\[
  \Delta U=C_V(T_2-T_1)
  \quad\text{用于定容过程，或用于理想气体任意过程。}
\]
\[
  \Delta H=C_p(T_2-T_1)
  \quad\text{用于定压过程，或用于理想气体任意过程。}
\]
第二句话成立的前提是理想气体 \(H=H(T)\)。
\end{FormulaBox}

\section{热力学第一定律对理想气体的应用}

\subsection{理想气体的内能和焓}

理想气体的内能只与温度有关：
\[
  \boxed{
  U=U(T)
  }.
\]
理想气体的焓也只与温度有关：
\[
  \boxed{
  H=H(T)
  }.
\]

因此对 \(\nu\) mol 理想气体，
\[
  \boxed{
  \dd U=\nu C_{V,m}\dd T
  },
  \qquad
  \boxed{
  \dd H=\nu C_{p,m}\dd T
  }.
\]

若摩尔热容近似为常数：
\[
  \Delta U=\nu C_{V,m}(T_2-T_1),
  \qquad
  \Delta H=\nu C_{p,m}(T_2-T_1).
\]

\subsection{迈耶公式}

对理想气体，
\[
  H=U+pV=U+\nu RT.
\]
两边对 \(T\) 求导：
\[
  \nu C_{p,m}
  =
  \nu C_{V,m}+\nu R.
\]
故
\[
  \boxed{
  C_{p,m}-C_{V,m}=R
  }.
\]

定义热容比
\[
  \boxed{
  \gamma=\frac{C_{p,m}}{C_{V,m}}
  }.
\]

由迈耶公式可得
\[
  C_{V,m}=\frac{R}{\gamma-1},
  \qquad
  C_{p,m}=\frac{\gamma R}{\gamma-1}.
\]

\begin{WarnBox}[易错点：迈耶公式只适用于理想气体]
\[
  C_{p,m}-C_{V,m}=R
\]
是理想气体结论。真实气体、液体和固体不能直接套用。
\end{WarnBox}

\subsection{理想气体等容过程}

等容过程满足
\[
  V=\mathrm{const},
  \qquad
  W=0.
\]
由第一定律：
\[
  Q=\Delta U.
\]
因此
\[
  \boxed{
  Q=\nu C_{V,m}(T_2-T_1)
  },
  \qquad
  \boxed{
  W=0
  }.
\]

\subsection{理想气体等压过程}

等压过程满足
\[
  p=\mathrm{const}.
\]
外界对系统做功：
\[
  W=-p(V_2-V_1).
\]
由理想气体状态方程，
\[
  p(V_2-V_1)=\nu R(T_2-T_1),
\]
所以
\[
  \boxed{
  W=-\nu R(T_2-T_1)
  }.
\]

吸热量：
\[
  Q=\Delta H.
\]
因此
\[
  \boxed{
  Q=\nu C_{p,m}(T_2-T_1)
  }.
\]

内能变化：
\[
  \boxed{
  \Delta U=\nu C_{V,m}(T_2-T_1)
  }.
\]

\subsection{理想气体等温过程}

等温过程满足
\[
  T=\mathrm{const}.
\]
理想气体内能只与温度有关，所以
\[
  \boxed{\Delta U=0}.
\]

由第一定律
\[
  \Delta U=Q+W=0,
\]
故
\[
  \boxed{
  Q=-W
  }.
\]

又
\[
  W=-\nu RT\ln\frac{V_2}{V_1},
\]
所以
\[
  \boxed{
  Q=\nu RT\ln\frac{V_2}{V_1}
  }.
\]

\begin{WarnBox}[等温过程不是绝热过程]
理想气体等温膨胀时 \(\Delta U=0\)，但气体对外做功，需要从外界吸热补偿。等温不代表不吸热。
\end{WarnBox}

\subsection{理想气体绝热过程}

绝热过程满足
\[
  \boxed{Q=0}.
\]
第一定律变为
\[
  \dd U=\dbar W.
\]
对理想气体：
\[
  \nu C_{V,m}\dd T=-p\dd V.
\]

利用
\[
  pV=\nu RT
\]
可推得可逆绝热过程方程：
\[
  \boxed{
  pV^\gamma=C_1
  },
  \qquad
  \boxed{
  TV^{\gamma-1}=C_2
  },
  \qquad
  \boxed{
  p^{1-\gamma}T^\gamma=C_3
  }.
\]

绝热过程中
\[
  W=\Delta U,
\]
因此
\[
  \boxed{
  W=\nu C_{V,m}(T_2-T_1)
  }.
\]

也可写为
\[
  W=\frac{p_2V_2-p_1V_1}{\gamma-1}.
\]

\begin{WarnBox}[绝热线与等温线]
在同一点的 \(p\)--\(V\) 图上，绝热线比等温线更陡。因为绝热膨胀时气体不吸热，内能下降，温度降低，压强下降得比等温膨胀更快。
\end{WarnBox}

\subsection{理想气体多方过程}

实际气体过程中常可近似满足
\[
  \boxed{
  pV^n=C
  }.
\]
这种过程称为\concept{多方过程}，其中 \(n\) 称为多方指数。

由理想气体状态方程可得等价形式：
\[
  \boxed{
  TV^{n-1}=C_2
  },
  \qquad
  \boxed{
  p^{1-n}T^n=C_3
  }.
\]

几种特殊情形：
\[
  n=0 \Rightarrow p=\mathrm{const}\quad\text{等压过程},
\]
\[
  n=1 \Rightarrow pV=\mathrm{const}\quad\text{等温过程},
\]
\[
  n=\gamma \Rightarrow pV^\gamma=\mathrm{const}\quad\text{绝热过程},
\]
\[
  n\to\infty \Rightarrow V=\mathrm{const}\quad\text{等容过程}.
\]

\begin{FormulaBox}[多方过程的功]
当 \(n\ne 1\) 时，
\[
  W=-\int_{V_1}^{V_2}p\dd V
  =
  \boxed{
  \frac{p_2V_2-p_1V_1}{n-1}
  }.
\]
利用 \(pV=\nu RT\)，也可写为
\[
  \boxed{
  W=\frac{\nu R(T_2-T_1)}{n-1}
  }.
\]
当 \(n=1\) 时，退化为等温过程：
\[
  W=-\nu RT\ln\frac{V_2}{V_1}.
\]
\end{FormulaBox}

多方过程摩尔热容定义为
\[
  C_{n,m}=\frac{\dbar Q}{\nu\dd T}.
\]
由第一定律可得
\[
  \boxed{
  C_{n,m}
  =
  C_{V,m}\frac{\gamma-n}{1-n}
  }.
\]

\begin{WarnBox}[多方热容的特殊值]
\[
  n=0:\quad C_{n,m}=C_{p,m}.
\]
\[
  n=\gamma:\quad C_{n,m}=0.
\]
\[
  n=1:\quad C_{n,m}\to\infty.
\]
等温过程温度不变，有限热量对应 \(\dd T=0\)，所以过程热容形式上发散。
\end{WarnBox}

\section{实际气体：绝热节流与焦耳--汤姆孙效应}

\subsection{节流过程}

\concept{节流过程}是气体通过多孔塞、细孔或节流阀，由高压侧流向低压侧的过程。典型条件是：
\begin{itemize}
  \item 管壁绝热；
  \item 过程前后达到稳定流动状态；
  \item 气体从高压流向低压。
\end{itemize}

节流过程不是准静态过程，不能用一条实线过程曲线表示。它通常在热力学坐标图中用虚线连接初末态。

\subsection{节流过程是等焓过程}

设初态为 \(1\)，末态为 \(2\)。绝热节流中
\[
  Q=0.
\]
外界推动气体进入节流区、气体推出外界所涉及的流动功合并后，可由第一定律得到
\[
  U_1+p_1V_1=U_2+p_2V_2.
\]
即
\[
  \boxed{
  H_1=H_2
  }.
\]

所以节流过程是\concept{等焓过程}。

\begin{WarnBox}[易错点：节流过程不是沿等焓线进行]
节流过程的初末态焓相等，但中间经过的是非平衡过程，不能说系统沿某条等焓线准静态变化。更准确地说：节流前后两个平衡态满足 \(H_1=H_2\)。
\end{WarnBox}

\subsection{焦耳--汤姆孙效应}

\concept{焦耳--汤姆孙效应}是指气体在节流过程中温度随压强变化的现象。

定义焦耳--汤姆孙系数：
\[
  \boxed{
  \alpha_i
  =
  \left(\frac{\partial T}{\partial p}\right)_H
  }.
\]

因为节流过程中压强降低，
\[
  \dd p<0.
\]
若
\[
  \alpha_i>0,
\]
则
\[
  \dd T=\alpha_i\dd p<0,
\]
气体节流后降温。

若
\[
  \alpha_i<0,
\]
则
\[
  \dd T>0,
\]
气体节流后升温。

若
\[
  \alpha_i=0,
\]
节流前后温度不变，对应反转点。

\begin{KeyBox}[焦耳--汤姆孙效应的复习重点]
热学B不需要把真实气体节流曲线推得太深。重点是：
\begin{ReviewList}
  \item 节流过程绝热、非准静态。
  \item 初末态满足 \(H_1=H_2\)。
  \item \(\alpha_i=(\partial T/\partial p)_H\) 判断节流后升温还是降温。
  \item 节流制冷依赖 \(\alpha_i>0\) 的区域。
\end{ReviewList}
\end{KeyBox}

\section{循环过程、热机与制冷机}

\subsection{循环过程}

系统经过一系列过程后回到初态，称为循环过程。由于内能为态函数，
\[
  \Delta U_{\mathrm{cycle}}=0.
\]
由第一定律：
\[
  0=Q_{\mathrm{cycle}}+W_{\mathrm{cycle}}.
\]

若采用系统对外做功 \(W'\) 为正，则
\[
  W'=Q_{\mathrm{net}}.
\]

\begin{WarnBox}[符号切换]
前面 \(W\) 表示外界对系统做功。讨论热机时，习惯用 \(W'\) 表示系统对外输出的功。两者关系为
\[
  W'=-W.
\]
\end{WarnBox}

\subsection{热机效率}

热机从高温热源吸收热量 \(Q_1\)，向低温热源放出热量 \(Q_2\)，对外输出功
\[
  W'=Q_1-Q_2.
\]
热机效率定义为
\[
  \boxed{
  \eta
  =
  \frac{W'}{Q_1}
  =
  1-\frac{Q_2}{Q_1}
  }.
\]

\begin{KeyBox}[正循环与热机]
热机必须至少与两个热源交换热量。若只有一个热源并把吸收热量全部转化为功，会涉及第二定律问题；第一定律本身只给能量数量关系。
\end{KeyBox}

\subsection{卡诺循环效率}

卡诺循环由两个等温过程和两个绝热过程组成。设高温热源温度为 \(T_1\)，低温热源温度为 \(T_2\)，则卡诺热机效率为
\[
  \boxed{
  \eta_{\mathrm C}=1-\frac{T_2}{T_1}
  }.
\]

这里 \(T_1,T_2\) 必须使用热力学温度，单位为 \(\mathrm K\)。

\begin{WarnBox}[易错点：不能用摄氏温度直接代入卡诺效率]
卡诺效率公式中的温度必须是绝对温度：
\[
  \eta_{\mathrm C}=1-\frac{T_2}{T_1}.
\]
若题目给的是摄氏温度，必须先转为开尔文温度。
\end{WarnBox}

\subsection{奥托循环效率}

奥托循环是理想汽油机循环模型，由两个绝热过程和两个等容过程组成。其热效率为
\[
  \boxed{
  \eta
  =
  1-\left(\frac{V_2}{V_1}\right)^{\gamma-1}
  }.
\]

若定义压缩比
\[
  r=\frac{V_1}{V_2},
\]
则
\[
  \boxed{
  \eta=1-\frac{1}{r^{\gamma-1}}
  }.
\]

\begin{KeyBox}[奥托循环记忆法]
奥托循环的吸热和放热都发生在等容过程，因此
\[
  Q_1=\nu C_{V,m}(T_3-T_2),
  \qquad
  Q_2=\nu C_{V,m}(T_4-T_1).
\]
效率从
\[
  \eta=1-\frac{Q_2}{Q_1}
\]
出发，再利用两段绝热关系化简。
\end{KeyBox}

\subsection{狄塞尔循环效率}

狄塞尔循环是理想柴油机循环模型，通常包括：
\begin{itemize}
  \item 绝热压缩；
  \item 等压吸热；
  \item 绝热膨胀；
  \item 等容放热。
\end{itemize}

热效率可写为
\[
  \eta
  =
  1-\frac{Q_2}{Q_1}.
\]
其中
\[
  Q_1=\nu C_{p,m}(T_3-T_2),
  \qquad
  Q_2=\nu C_{V,m}(T_4-T_1).
\]
所以
\[
  \boxed{
  \eta
  =
  1-\frac{1}{\gamma}\frac{T_4-T_1}{T_3-T_2}
  }.
\]

按授课知识要点中体积比形式，也可写为
\[
  \boxed{
  \eta
  =
  1-
  \frac{(V_3/V_2)^\gamma-1}
  {\gamma (V_1/V_2)^{\gamma-1}(V_3/V_2-1)}
  }.
\]

\begin{WarnBox}[奥托与狄塞尔不要混]
奥托循环：等容吸热、等容放热，用 \(C_V\)。狄塞尔循环：等压吸热、等容放热，吸热用 \(C_p\)，放热用 \(C_V\)。
\end{WarnBox}

\subsection{制冷机与制冷系数}

制冷机是逆循环装置。它从低温热源吸收热量 \(Q_2\)，向高温热源放出热量 \(Q_1\)，外界需要对其输入功
\[
  W=Q_1-Q_2.
\]

制冷系数定义为
\[
  \boxed{
  \varepsilon
  =
  \frac{Q_2}{W}
  =
  \frac{Q_2}{Q_1-Q_2}
  }.
\]

对逆卡诺循环，
\[
  \boxed{
  \varepsilon_{\mathrm C}
  =
  \frac{T_2}{T_1-T_2}
  }.
\]

\begin{KeyBox}[热机效率与制冷系数的区别]
热机效率关心“输入热量中有多少变成对外功”：
\[
  \eta=\frac{W'}{Q_1}.
\]
制冷系数关心“输入单位功能从低温处搬走多少热量”：
\[
  \varepsilon=\frac{Q_2}{W}.
\]
制冷系数可以大于 \(1\)，这不违反能量守恒。
\end{KeyBox}

\section{第二章题型模板}

\subsection{题型一：给定过程求 \(Q,W,\Delta U\)}

\begin{MethodBox}[通用流程]
\begin{ReviewList}
  \item 明确过程类型：等容、等压、等温、绝热、多方、循环。
  \item 先求功：
  \[
    W=-\int p\dd V.
  \]
  \item 对理想气体先求温度变化：
  \[
    \Delta U=\nu C_{V,m}(T_2-T_1).
  \]
  \item 最后由第一定律：
  \[
    Q=\Delta U-W.
  \]
  注意这里 \(W\) 是外界对系统做功。
\end{ReviewList}
\end{MethodBox}

\subsection{题型二：理想气体过程速查}

\begin{FormulaBox}[理想气体常见过程]
\begin{tabularx}{\textwidth}{@{}>{\bfseries}lXXX@{}}
\toprule
过程 & \(W\)：外界对系统做功 & \(\Delta U\) & \(Q\)\\
\midrule
等容 & \(0\) & \(\nu C_{V,m}\Delta T\) & \(\nu C_{V,m}\Delta T\)\\[0.5em]
等压 & \(-\nu R\Delta T\) & \(\nu C_{V,m}\Delta T\) & \(\nu C_{p,m}\Delta T\)\\[0.5em]
等温 & \(-\nu RT\ln(V_2/V_1)\) & \(0\) & \(\nu RT\ln(V_2/V_1)\)\\[0.5em]
绝热 & \(\nu C_{V,m}\Delta T\) & \(\nu C_{V,m}\Delta T\) & \(0\)\\[0.5em]
多方 \(n\ne1\) & \(\dfrac{\nu R\Delta T}{n-1}\) & \(\nu C_{V,m}\Delta T\) & \(\nu C_{n,m}\Delta T\)\\
\bottomrule
\end{tabularx}
\end{FormulaBox}

\subsection{题型三：循环效率}

\begin{MethodBox}[热机效率计算路线]
\begin{ReviewList}
  \item 找出吸热过程，计算总吸热 \(Q_1\)。
  \item 找出放热过程，计算放热量的绝对值 \(Q_2\)。
  \item 用
  \[
    \eta=1-\frac{Q_2}{Q_1}.
  \]
  \item 若是卡诺循环，直接用
  \[
    \eta=1-\frac{T_2}{T_1}.
  \]
\end{ReviewList}
\end{MethodBox}

\section{第二章公式速查}

\formulatable{
\begin{tabularx}{\textwidth}{@{}>{\bfseries}lXl@{}}
\toprule
内容 & 公式 & 条件\\
\midrule
准静态体积功 &
\(\displaystyle \dbar W=-p\dd V,\quad W=-\int_{V_1}^{V_2}p\dd V\)
& 外界对系统做功\\[0.8em]

热容定义 &
\(\displaystyle C_L=\left(\frac{\dbar Q}{\dd T}\right)_L\)
& 指定过程 \(L\)\\[0.8em]

第一定律 &
\(\displaystyle \dd U=\dbar Q+\dbar W\)
& 普遍能量守恒\\[0.8em]

只有体积功时 &
\(\displaystyle \dd U=\dbar Q-p\dd V\)
& 准静态体积功\\[0.8em]

焓 &
\(\displaystyle H=U+pV,\quad \dbar Q_p=\dd H\)
& 等压过程\\[0.8em]

定容热容 &
\(\displaystyle C_V=\left(\frac{\partial U}{\partial T}\right)_V\)
& 一般系统\\[0.8em]

定压热容 &
\(\displaystyle C_p=\left(\frac{\partial H}{\partial T}\right)_p\)
& 一般系统\\[0.8em]

理想气体内能 &
\(\displaystyle U=U(T),\quad \dd U=\nu C_{V,m}\dd T\)
& 理想气体\\[0.8em]

理想气体焓 &
\(\displaystyle H=H(T),\quad \dd H=\nu C_{p,m}\dd T\)
& 理想气体\\[0.8em]

迈耶公式 &
\(\displaystyle C_{p,m}-C_{V,m}=R\)
& 理想气体\\[0.8em]

绝热过程 &
\(\displaystyle pV^\gamma=C_1,\quad TV^{\gamma-1}=C_2,\quad p^{1-\gamma}T^\gamma=C_3\)
& 理想气体可逆绝热\\[1em]

多方过程 &
\(\displaystyle pV^n=C,\quad C_{n,m}=C_{V,m}\frac{\gamma-n}{1-n}\)
& 理想气体\\[1em]

节流过程 &
\(\displaystyle H_1=H_2,\quad \alpha_i=\left(\frac{\partial T}{\partial p}\right)_H\)
& 绝热节流\\[1em]

热机效率 &
\(\displaystyle \eta=\frac{W'}{Q_1}=1-\frac{Q_2}{Q_1}\)
& 正循环\\[0.8em]

卡诺效率 &
\(\displaystyle \eta_{\mathrm C}=1-\frac{T_2}{T_1}\)
& 可逆卡诺热机\\[0.8em]

制冷系数 &
\(\displaystyle \varepsilon=\frac{Q_2}{W}=\frac{Q_2}{Q_1-Q_2}\)
& 逆循环\\[0.8em]

卡诺制冷系数 &
\(\displaystyle \varepsilon_{\mathrm C}=\frac{T_2}{T_1-T_2}\)
& 逆卡诺循环\\
\bottomrule
\end{tabularx}
}

\section{第二章易错点总表}

\begin{WarnBox}[本章高频错误]
\begin{PitfallList}
  \item 混淆“外界对系统做功”和“系统对外做功”的符号。本资料默认 \(W\) 为外界对系统做功。
  \item 把功和热量当成状态函数。功、热量都是过程量，不能写成态函数差值。
  \item 非准静态过程直接用 \(W=-\int p\dd V\)。非准静态过程一般不能用系统内部压强 \(p\) 积分。
  \item 忘记等容过程 \(W=0\)，等温理想气体过程 \(\Delta U=0\)，绝热过程 \(Q=0\)。
  \item 把理想气体的 \(U=U(T)\) 用到真实气体上。真实气体内能一般与 \(T,V\) 都有关。
  \item 把 \(C_p-C_V=R\) 用到非理想气体、液体或固体上。
  \item 绝热方程 \(pV^\gamma=\mathrm{const}\) 只适用于理想气体可逆绝热过程。
  \item 多方过程中 \(n=1\) 是等温过程，不能直接套 \(W=\nu R\Delta T/(n-1)\)。
  \item 节流过程只保证初末态等焓，不是准静态地沿等焓线变化。
  \item 卡诺效率和卡诺制冷系数中的温度必须使用开尔文温度。
  \item 制冷系数可以大于 \(1\)，这不是错误，也不违反第一定律。
\end{PitfallList}
\end{WarnBox}

\section{第二章复习路线}

\begin{MethodBox}[建议背诵顺序]
\begin{ReviewList}
  \item 先确定符号约定：外界对系统做功 \(W=-\int p\dd V\)。
  \item 再背第一定律：\(\dd U=\dbar Q+\dbar W\)。
  \item 再背态函数与过程量：\(U,H\) 是态函数，\(Q,W\) 是过程量。
  \item 再背热容定义：\(C_V=(\partial U/\partial T)_V\)，\(C_p=(\partial H/\partial T)_p\)。
  \item 再背理想气体结论：\(U=U(T)\)，\(H=H(T)\)，\(C_{p,m}-C_{V,m}=R\)。
  \item 再整理五类过程：等容、等压、等温、绝热、多方。
  \item 最后看循环：\(\eta=1-Q_2/Q_1\)，\(\varepsilon=Q_2/(Q_1-Q_2)\)。
\end{ReviewList}
\end{MethodBox}

\begin{KeyBox}[第二章一句话总结]
第二章的核心计算逻辑是：先判断过程，再算功和内能变化，最后用第一定律补出热量。所有复杂题本质上都在反复使用
\[
  \dd U=\dbar Q-p\dd V
\]
以及理想气体的
\[
  \Delta U=\nu C_{V,m}\Delta T.
\]
\end{KeyBox}
\chapter{热力学第二定律与熵}

\chaptergoal{
本章解决第一定律无法回答的问题：满足能量守恒的过程，是否一定能够自发发生？热力学第二定律给出自然过程的方向性；熵则把这种方向性写成态函数语言。复习时要抓住一条主线：不可逆性 \(\rightarrow\) 第二定律经典表述 \(\rightarrow\) 卡诺定理 \(\rightarrow\) 克劳修斯不等式 \(\rightarrow\) 熵 \(\rightarrow\) 熵增加原理。
}

\exammap{
\begin{ReviewList}
  \item \textbf{会判断：}可逆过程与不可逆过程，知道实际宏观热过程通常不可逆。
  \item \textbf{会表述：}热力学第二定律的开尔文表述和克劳修斯表述，并理解二者等价。
  \item \textbf{会使用：}卡诺定理，知道可逆热机效率只由两个热源温度决定。
  \item \textbf{会理解：}热力学温标不依赖具体测温物质，并与理想气体温标一致。
  \item \textbf{会推导：}从克劳修斯等式引入熵，从克劳修斯不等式得到熵增加原理。
  \item \textbf{会计算：}理想气体、热源、相变、自由膨胀等常见过程的熵变。
  \item \textbf{会读图：}\(T\)--\(S\) 图中可逆过程吸热量对应曲线下面积。
\end{ReviewList}
}

\section{为什么需要热力学第二定律}

热力学第一定律说明一个过程若要发生，必须满足能量守恒。但经验表明，满足能量守恒的过程并不一定会自发发生。

\begin{KeyBox}[自然过程的方向性]
\begin{itemize}
  \item 热量会自发地从高温物体传到低温物体，但反向过程不会自发发生。
  \item 气体可以自由膨胀充满容器，但不会自发收缩回原来一侧。
  \item 机械功可以完全转化为热，例如摩擦生热；但热量不能在不引起其它变化的情况下完全转化为有用功。
  \item 两种气体可以自发扩散混合，但不会自发反扩散分离。
\end{itemize}
这些过程都不违反第一定律，但自发方向是确定的。
\end{KeyBox}

热力学第二定律研究的正是自然过程的\concept{方向性}与\concept{不可逆性}。

\section{可逆过程与不可逆过程}

\subsection{定义}

\concept{可逆过程}：一个热力学系统由某状态出发，经过一个过程达到另一状态；若存在另一过程或某种方法，可以使系统和外界都完全恢复到原来的状态，则原过程称为可逆过程。

\concept{不可逆过程}：若在不引起其它变化的条件下，不能使系统和外界都完全复原，则该过程称为不可逆过程。

\begin{FormulaBox}[可逆过程的条件]
一个过程要可逆，至少需要满足：
\begin{ReviewList}
  \item 必须是准静态过程，即过程经过一系列无限接近平衡态的状态。
  \item 过程中不能包含摩擦、磁滞、电阻、有限温差传热、有限压强差膨胀等耗散效应。
\end{ReviewList}
\end{FormulaBox}

\begin{WarnBox}[易错点：准静态不一定可逆]
可逆过程一定是准静态过程，但准静态过程不一定可逆。若准静态过程中仍存在摩擦、电阻、磁滞等耗散效应，该过程仍然不可逆。
\end{WarnBox}

\subsection{四类不可逆因素}

\begin{KeyBox}[不可逆性的常见来源]
\begin{tabularx}{\textwidth}{@{}>{\bfseries}lX@{}}
耗散不可逆 & 功自发转化为热，例如摩擦、电阻发热、黏滞耗散。\\
力学不可逆 & 系统内部各部分之间存在非无穷小压强差。\\
热学不可逆 & 系统内部或系统与外界之间存在非无穷小温度差。\\
化学不可逆 & 系统内部存在非无穷小浓度差或化学组成差异，例如扩散、反应。\\
\end{tabularx}
\end{KeyBox}

只要过程中包含上述任一不可逆因素，就不是可逆过程。

\begin{WarnBox}[热学B判断题常见陷阱]
“缓慢”不自动等于“可逆”。缓慢只可能保证近似准静态；是否可逆还要看有没有耗散因素。
\end{WarnBox}

\section{热力学第二定律的经典表述}

\subsection{开尔文表述}

\concept{开尔文表述}：不可能从单一热源吸取热量，使之完全变为有用功而不引起其它变化。

这里“单一热源”指温度处处相同并保持恒定的热源；“其它变化”指除了从单一热源吸热并全部变为功以外的任何其它影响。

\begin{KeyBox}[开尔文表述的含义]
开尔文表述并不是说热不能变成功。热机可以把一部分热量转化为功。它说的是：不能只从一个热源吸热，并把这部分热量全部变成有用功，同时不产生其它影响。
\end{KeyBox}

由开尔文表述可知，\concept{第二类永动机}不可能制造。第二类永动机指只从单一热源吸热，并将其全部转化为有用功而不引起其它变化的机器。

\subsection{克劳修斯表述}

\concept{克劳修斯表述}：不可能把热量从低温物体传到高温物体而不引起其它变化。

\begin{KeyBox}[克劳修斯表述的含义]
克劳修斯表述并不是说热量绝对不能从低温物体传到高温物体。冰箱、空调都能把热量从低温处搬到高温处，但必须消耗外界功。因此，禁止的是“无代价地从低温向高温传热”。
\end{KeyBox}

克劳修斯表述说明热传导过程是不可逆过程。

\subsection{两种表述的等价性}

开尔文表述和克劳修斯表述是等价的。其证明通常采用反证法：

\begin{MethodBox}[等价性证明思路]
\begin{ReviewList}
  \item 若开尔文表述不成立，就存在一种装置，可以从单一高温热源吸热并全部变成功。
  \item 用这部分功驱动一台制冷机，就可以把热量从低温热源搬到高温热源，而整体不引起其它变化。
  \item 这将违反克劳修斯表述。
  \item 反过来，若克劳修斯表述不成立，即热量可以自动从低温热源传到高温热源，则可与普通热机组合，构造出从单一热源吸热并全部变成功的装置，违反开尔文表述。
\end{ReviewList}
\end{MethodBox}

\begin{WarnBox}[易错点：第二定律不是第一定律的推论]
第一定律只限制能量数值守恒，第二定律限制过程方向。热量从低温物体自发传给高温物体不违反能量守恒，但违反第二定律。
\end{WarnBox}

\section{卡诺定理与热力学温标}

\subsection{卡诺循环回顾}

卡诺循环由四个可逆过程组成：
\begin{ReviewList}
  \item 高温热源 \(T_1\) 下的可逆等温膨胀，吸热 \(Q_1\)。
  \item 可逆绝热膨胀，温度由 \(T_1\) 降至 \(T_2\)。
  \item 低温热源 \(T_2\) 下的可逆等温压缩，放热 \(Q_2\)。
  \item 可逆绝热压缩，温度由 \(T_2\) 升回 \(T_1\)。
\end{ReviewList}

热机效率为
\[
  \eta=\frac{W'}{Q_1}=1-\frac{Q_2}{Q_1}.
\]

\subsection{卡诺定理}

\concept{卡诺定理}包括两个结论：

\begin{FormulaBox}[卡诺定理]
\begin{ReviewList}
  \item 工作在相同高温热源和低温热源之间的一切可逆热机，其效率相等，与工作物质无关。
  \item 工作在相同两个热源之间的一切不可逆热机，其效率都不可能大于可逆热机效率。
\end{ReviewList}
即
\[
  \eta_{\mathrm{ir}}\le \eta_{\mathrm{rev}}.
\]
\end{FormulaBox}

\begin{KeyBox}[卡诺定理的物理意义]
卡诺定理把热机效率的上限与工作物质、循环细节分离开来。只要两个热源温度给定，最高效率就由可逆热机决定。
\end{KeyBox}

\subsection{热力学温标}

由于所有可逆热机在相同两个热源之间效率相同，所以可用可逆热机来定义不依赖测温物质的温标，即\concept{热力学温标}。

对可逆卡诺热机，定义
\[
  \frac{Q_2}{Q_1}=\frac{T_2}{T_1}.
\]
于是卡诺效率为
\[
  \boxed{
  \eta_{\mathrm C}=1-\frac{T_2}{T_1}
  }.
\]

\begin{WarnBox}[温度必须用开尔文]
卡诺效率中的 \(T_1,T_2\) 必须是热力学温度，单位为 \(\mathrm K\)。不能直接代入摄氏温度。
\end{WarnBox}

\subsection{卡诺制冷机}

卡诺循环逆向运行就是卡诺制冷机。它从低温热源吸热 \(Q_2\)，向高温热源放热 \(Q_1\)，外界输入功
\[
  W=Q_1-Q_2.
\]
制冷系数为
\[
  \varepsilon=\frac{Q_2}{W}.
\]
对逆卡诺循环：
\[
  \boxed{
  \varepsilon_{\mathrm C}=\frac{T_2}{T_1-T_2}
  }.
\]

\begin{KeyBox}[卡诺热机与卡诺制冷机]
\[
  \eta_{\mathrm C}=1-\frac{T_2}{T_1},
  \qquad
  \varepsilon_{\mathrm C}=\frac{T_2}{T_1-T_2}.
\]
二者都来自同一条热量比关系：
\[
  \frac{Q_2}{Q_1}=\frac{T_2}{T_1}.
\]
\end{KeyBox}

\section{克劳修斯等式与不等式}

\subsection{克劳修斯等式}

对任意可逆循环，可以把循环分解为许多微小卡诺循环。由卡诺循环的热量与温度关系，可得
\[
  \boxed{
  \oint_{\mathrm{rev}}\frac{\dbar Q}{T}=0
  }.
\]
这称为\concept{克劳修斯等式}。

\begin{KeyBox}[克劳修斯等式的意义]
对可逆循环，\(\dbar Q/T\) 的循环积分为零。这说明 \(\dbar Q_{\mathrm{rev}}/T\) 具有态函数全微分的性质，从而可以引入新的态函数：熵。
\end{KeyBox}

\subsection{克劳修斯不等式}

对不可逆循环，有
\[
  \boxed{
  \oint_{\mathrm{ir}}\frac{\dbar Q}{T}<0
  }.
\]
综合可逆和不可逆循环：
\[
  \boxed{
  \oint\frac{\dbar Q}{T}\le 0
  }.
\]
其中等号对应可逆循环，小于号对应不可逆循环。

\begin{WarnBox}[克劳修斯不等式中的 \(T\)]
\[
  \oint\frac{\dbar Q}{T}\le0
\]
中的 \(T\) 是发生热交换处外界热源的温度；对于可逆过程，它也等于系统边界处温度。不可逆有限温差传热时，系统内部未必有统一温度，不能随意把系统某个温度代进去。
\end{WarnBox}

\section{熵}

\subsection{熵的定义}

由克劳修斯等式，任意两个平衡态 \(a,b\) 之间，可逆路径积分
\[
  \int_a^b\frac{\dbar Q_{\mathrm{rev}}}{T}
\]
只与初末态有关，与可逆路径无关。因此定义态函数\concept{熵} \(S\)：
\[
  \boxed{
  S_b-S_a
  =
  \int_a^b\frac{\dbar Q_{\mathrm{rev}}}{T}
  }.
\]
微分形式为
\[
  \boxed{
  \dd S=\frac{\dbar Q_{\mathrm{rev}}}{T}
  }.
\]

\begin{WarnBox}[易错点：熵变计算必须用可逆路径]
即使实际过程不可逆，熵变仍然只由初末态决定。计算不可逆过程的熵变时，不能直接把实际过程中的 \(\dbar Q/T\) 积分；应在相同初末态之间设计一条可逆路径，然后计算
\[
  \Delta S=\int\frac{\dbar Q_{\mathrm{rev}}}{T}.
\]
\end{WarnBox}

\subsection{热力学第二定律的数学表达式}

对任意过程，有
\[
  \boxed{
  \dd S\ge \frac{\dbar Q}{T}
  }.
\]
其中等号对应可逆过程，不等号对应不可逆过程。

若系统绝热，
\[
  \dbar Q=0,
\]
则
\[
  \boxed{
  \dd S\ge0.
  }
\]

\section{熵增加原理}

\subsection{孤立系统的熵增加}

孤立系统与外界既不交换热量，也不交换功和物质。对孤立系统，
\[
  \dbar Q=0.
\]
由第二定律的数学表达式：
\[
  \dd S\ge0.
\]
因此
\[
  \boxed{
  \Delta S_{\mathrm{isolated}}\ge0.
  }
\]

这就是\concept{熵增加原理}：孤立系统中，自然发生的过程总是向熵不减的方向进行。

\begin{FormulaBox}[熵增加原理]
\[
  \Delta S_{\mathrm{isolated}}>0
  \quad\text{不可逆过程},
\]
\[
  \Delta S_{\mathrm{isolated}}=0
  \quad\text{可逆过程}.
\]
孤立系统中不可能出现
\[
  \Delta S_{\mathrm{isolated}}<0.
\]
\end{FormulaBox}

\subsection{系统与外界的总熵}

若研究对象不是孤立系统，可以把系统与相关外界合在一起组成孤立系统，则
\[
  \Delta S_{\mathrm{total}}
  =
  \Delta S_{\mathrm{system}}+\Delta S_{\mathrm{environment}}\ge0.
\]

\begin{KeyBox}[判断过程能否自发发生]
某一过程若要自发发生，需要满足
\[
  \Delta S_{\mathrm{total}}>0.
\]
若
\[
  \Delta S_{\mathrm{total}}=0,
\]
则过程可逆。若
\[
  \Delta S_{\mathrm{total}}<0,
\]
则该过程不可能自发发生。
\end{KeyBox}

\begin{WarnBox}[易错点：系统熵可以减少]
熵增加原理适用于孤立系统或“系统+外界”的总熵。单独一个非孤立系统的熵可以减少，例如制冷机内部低温物体降温时熵减少，但外界熵增加更多，总熵仍增加。
\end{WarnBox}

\section{熵变的计算}

\subsection{计算熵变的通用方法}

\begin{MethodBox}[熵变计算三步法]
\begin{ReviewList}
  \item 明确初态和末态，因为熵是态函数。
  \item 判断实际过程是否可逆。若可逆，直接用实际路径积分：
  \[
    \Delta S=\int\frac{\dbar Q_{\mathrm{rev}}}{T}.
  \]
  \item 若实际过程不可逆，在同一初末态之间设计一条方便计算的可逆路径，再积分。
\end{ReviewList}
\end{MethodBox}

\subsection{等温可逆过程}

若系统在温度 \(T\) 下经历等温可逆过程，吸热量为 \(Q_{\mathrm{rev}}\)，则
\[
  \boxed{
  \Delta S=\frac{Q_{\mathrm{rev}}}{T}
  }.
\]

典型例子包括等温相变、理想气体等温可逆膨胀等。

\subsection{理想气体熵变}

对 \(\nu\) mol 理想气体，
\[
  \dbar Q_{\mathrm{rev}}=\dd U+p\dd V.
\]
又
\[
  \dd U=\nu C_{V,m}\dd T,
  \qquad
  p=\frac{\nu RT}{V}.
\]
因此
\[
  \dd S
  =
  \frac{\dbar Q_{\mathrm{rev}}}{T}
  =
  \nu C_{V,m}\frac{\dd T}{T}
  +
  \nu R\frac{\dd V}{V}.
\]
积分得
\[
  \boxed{
  \Delta S
  =
  \nu C_{V,m}\ln\frac{T_2}{T_1}
  +
  \nu R\ln\frac{V_2}{V_1}
  }.
\]

也可用 \(T,p\) 表示：
\[
  \boxed{
  \Delta S
  =
  \nu C_{p,m}\ln\frac{T_2}{T_1}
  -
  \nu R\ln\frac{p_2}{p_1}
  }.
\]

还可用 \(p,V\) 表示：
\[
  \boxed{
  \Delta S
  =
  \nu C_{V,m}\ln\frac{p_2}{p_1}
  +
  \nu C_{p,m}\ln\frac{V_2}{V_1}
  }.
\]

\begin{WarnBox}[理想气体熵变公式的适用条件]
上述公式默认摩尔热容可近似看作常数。若热容随温度变化，应保留积分形式：
\[
  \Delta S
  =
  \nu\int_{T_1}^{T_2}\frac{C_{V,m}(T)}{T}\dd T
  +
  \nu R\ln\frac{V_2}{V_1}.
\]
\end{WarnBox}

\subsection{理想气体等温膨胀}

理想气体等温过程 \(T_1=T_2=T\)，因此
\[
  \Delta S
  =
  \nu R\ln\frac{V_2}{V_1}.
\]

若 \(V_2>V_1\)，则
\[
  \Delta S>0.
\]

\begin{KeyBox}[自由膨胀的熵变]
理想气体自由膨胀是不可逆绝热过程：
\[
  Q=0,\qquad W=0,\qquad \Delta U=0,\qquad T_2=T_1.
\]
但熵变不是零。因为熵是态函数，可在同一初末态之间设计等温可逆膨胀路径：
\[
  \boxed{
  \Delta S=\nu R\ln\frac{V_2}{V_1}>0
  }.
\]
\end{KeyBox}

\subsection{可逆绝热过程}

可逆绝热过程满足
\[
  \dbar Q_{\mathrm{rev}}=0.
\]
因此
\[
  \boxed{
  \Delta S=0
  }.
\]
所以可逆绝热过程也称为\concept{等熵过程}。

\begin{WarnBox}[绝热过程不一定等熵]
可逆绝热过程才是等熵过程。不可逆绝热过程虽然 \(Q=0\)，但熵可以增加。例如自由膨胀是绝热不可逆过程，\(\Delta S>0\)。
\end{WarnBox}

\subsection{热源的熵变}

若一个温度恒定的大热源在温度 \(T\) 下吸收热量 \(Q\)，则热源熵变为
\[
  \boxed{
  \Delta S_{\mathrm{res}}=\frac{Q}{T}
  }.
\]
若热源放出热量 \(Q\)，则
\[
  \Delta S_{\mathrm{res}}=-\frac{Q}{T}.
\]

\begin{MethodBox}[有限温差传热的总熵变]
高温热源 \(T_1\) 向低温热源 \(T_2\) 传热 \(Q\)，其中 \(T_1>T_2\)。则
\[
  \Delta S_1=-\frac{Q}{T_1},
  \qquad
  \Delta S_2=\frac{Q}{T_2}.
\]
总熵变
\[
  \Delta S_{\mathrm{total}}
  =
  Q\left(\frac{1}{T_2}-\frac{1}{T_1}\right)>0.
\]
这说明有限温差传热不可逆。
\end{MethodBox}

\subsection{相变过程的熵变}

若物质在温度 \(T\) 下发生可逆等温相变，相变潜热为 \(L\)，则
\[
  \boxed{
  \Delta S=\frac{L}{T}
  }.
\]
这里 \(L\) 是系统吸收的热量。熔化、汽化通常吸热，熵增加；凝固、液化放热，熵减少。

\section{\texorpdfstring{\(T\)--\(S\)}{T-S} 图}

\subsection{\texorpdfstring{\(T\)--\(S\)}{T-S} 图的基本意义}

在 \(T\)--\(S\) 图中，横轴为熵 \(S\)，纵轴为温度 \(T\)。对可逆过程，
\[
  \dbar Q_{\mathrm{rev}}=T\dd S.
\]
所以可逆过程中吸收的热量为
\[
  \boxed{
  Q_{\mathrm{rev}}=\int T\dd S
  }.
\]

\begin{KeyBox}[\(T\)--\(S\) 图中的面积]
在 \(T\)--\(S\) 图上，可逆过程曲线下方面积表示该过程的吸热量：
\[
  Q_{\mathrm{rev}}=\int T\dd S.
\]
可逆循环包围的面积表示循环净功：
\[
  W'=\oint T\dd S.
\]
\end{KeyBox}

\subsection{常见过程在 \texorpdfstring{\(T\)--\(S\)}{T-S} 图中的形状}

\begin{FormulaBox}[常见图像]
\begin{tabularx}{\textwidth}{@{}>{\bfseries}lX@{}}
可逆绝热过程 & \(\Delta S=0\)，在 \(T\)--\(S\) 图中为竖直线。\\
可逆等温过程 & \(T=\mathrm{const}\)，在 \(T\)--\(S\) 图中为水平线。\\
卡诺循环 & 两条等温线和两条绝热线构成矩形。\\
\end{tabularx}
\end{FormulaBox}

对卡诺循环，
\[
  Q_1=T_1(S_2-S_1),
  \qquad
  Q_2=T_2(S_2-S_1).
\]
故
\[
  \frac{Q_2}{Q_1}=\frac{T_2}{T_1},
\]
从而
\[
  \eta_{\mathrm C}=1-\frac{T_2}{T_1}.
\]

\section{热力学第二定律的统计意义}

热力学第二定律描述的是大量粒子系统的统计规律，而不是单个分子的力学规律。

\subsection{宏观不可逆性的微观原因}

以气体自由膨胀为例。隔板抽去后，气体分子从容器一侧扩散到整个容器。微观上，每个分子仍服从力学规律；但从统计角度看，所有分子重新自发聚集到原来一侧的概率极小。

\begin{KeyBox}[统计意义]
自然过程总是从概率小的宏观态向概率大的宏观态发展。平衡态对应的微观状态数最多，因此孤立系统趋向平衡态。
\end{KeyBox}

玻尔兹曼关系给出熵与微观状态数的联系：
\[
  \boxed{
  S=k_{\mathrm B}\ln W
  }.
\]
其中 \(W\) 表示对应某宏观态的微观状态数。

\begin{WarnBox}[第二定律的统计性质]
第二定律不是说熵减少在微观上绝对不可能，而是说对含有巨大数量粒子的宏观系统，熵自发减少的概率小到可以忽略。因此实际宏观过程表现为不可逆。
\end{WarnBox}

\section{第三章题型模板}

\subsection{题型一：判断过程方向}

\begin{MethodBox}[用总熵判断自发方向]
\begin{ReviewList}
  \item 明确系统和外界。
  \item 计算系统熵变 \(\Delta S_{\mathrm{sys}}\)。
  \item 计算外界熵变 \(\Delta S_{\mathrm{env}}\)。
  \item 相加：
  \[
    \Delta S_{\mathrm{total}}
    =
    \Delta S_{\mathrm{sys}}+\Delta S_{\mathrm{env}}.
  \]
  \item 判断：
  \[
    \Delta S_{\mathrm{total}}>0
    \Rightarrow
    \text{不可逆自发过程};
  \]
  \[
    \Delta S_{\mathrm{total}}=0
    \Rightarrow
    \text{可逆过程};
  \]
  \[
    \Delta S_{\mathrm{total}}<0
    \Rightarrow
    \text{不能自发发生}.
  \]
\end{ReviewList}
\end{MethodBox}

\subsection{题型二：不可逆过程熵变}

\begin{MethodBox}[不可逆过程熵变计算]
不可逆过程不能直接对实际过程使用
\[
  \int\frac{\dbar Q}{T}.
\]
正确做法是：
\begin{ReviewList}
  \item 固定初态和末态。
  \item 设计一条方便计算的可逆路径。
  \item 沿可逆路径计算
  \[
    \Delta S=\int\frac{\dbar Q_{\mathrm{rev}}}{T}.
  \]
\end{ReviewList}
典型例子：理想气体自由膨胀，用等温可逆膨胀路径计算熵变。
\end{MethodBox}

\subsection{题型三：热机或制冷机极限}

\begin{MethodBox}[卡诺类题目]
\begin{ReviewList}
  \item 若问最高热机效率，直接用
  \[
    \eta_{\max}=1-\frac{T_2}{T_1}.
  \]
  \item 若问最高制冷系数，直接用
  \[
    \varepsilon_{\max}=\frac{T_2}{T_1-T_2}.
  \]
  \item 所有温度先转成开尔文。
  \item 若实际机器效率超过卡诺值，说明题目条件不可能。
\end{ReviewList}
\end{MethodBox}

\section{第三章公式速查}

\formulatable{
\begin{tabularx}{\textwidth}{@{}>{\bfseries}lXl@{}}
\toprule
内容 & 公式 & 条件\\
\midrule

热机效率 &
\(\displaystyle \eta=\frac{W'}{Q_1}=1-\frac{Q_2}{Q_1}\)
& 正循环\\[0.8em]

卡诺热量比 &
\(\displaystyle \frac{Q_2}{Q_1}=\frac{T_2}{T_1}\)
& 可逆卡诺循环\\[0.8em]

卡诺效率 &
\(\displaystyle \eta_{\mathrm C}=1-\frac{T_2}{T_1}\)
& 可逆热机\\[0.8em]

制冷系数 &
\(\displaystyle \varepsilon=\frac{Q_2}{W}=\frac{Q_2}{Q_1-Q_2}\)
& 逆循环\\[0.8em]

卡诺制冷系数 &
\(\displaystyle \varepsilon_{\mathrm C}=\frac{T_2}{T_1-T_2}\)
& 逆卡诺循环\\[0.8em]

克劳修斯等式 &
\(\displaystyle \oint_{\mathrm{rev}}\frac{\dbar Q}{T}=0\)
& 可逆循环\\[0.8em]

克劳修斯不等式 &
\(\displaystyle \oint\frac{\dbar Q}{T}\le0\)
& 任意循环\\[0.8em]

熵定义 &
\(\displaystyle \dd S=\frac{\dbar Q_{\mathrm{rev}}}{T}\)
& 可逆路径\\[0.8em]

熵变 &
\(\displaystyle \Delta S=\int_a^b\frac{\dbar Q_{\mathrm{rev}}}{T}\)
& 任意过程的初末态\\[0.8em]

第二定律数学表达 &
\(\displaystyle \dd S\ge\frac{\dbar Q}{T}\)
& 等号可逆\\[0.8em]

孤立系统熵增加 &
\(\displaystyle \Delta S_{\mathrm{isolated}}\ge0\)
& 孤立系统\\[0.8em]

理想气体熵变 \(T,V\) 式 &
\(\displaystyle \Delta S=\nu C_{V,m}\ln\frac{T_2}{T_1}+\nu R\ln\frac{V_2}{V_1}\)
& \(C_{V,m}\) 近似常数\\[1em]

理想气体熵变 \(T,p\) 式 &
\(\displaystyle \Delta S=\nu C_{p,m}\ln\frac{T_2}{T_1}-\nu R\ln\frac{p_2}{p_1}\)
& \(C_{p,m}\) 近似常数\\[1em]

等温相变熵变 &
\(\displaystyle \Delta S=\frac{L}{T}\)
& 可逆等温相变\\[0.8em]

\(T\)--\(S\) 图吸热量 &
\(\displaystyle Q_{\mathrm{rev}}=\int T\dd S\)
& 可逆过程\\

\bottomrule
\end{tabularx}
}

\section{第三章易错点总表}

\begin{WarnBox}[本章高频错误]
\begin{PitfallList}
  \item 认为满足第一定律就一定能发生。第一定律管能量守恒，第二定律管过程方向。
  \item 把“准静态过程”直接等同于“可逆过程”。有耗散因素时，准静态也不可逆。
  \item 误解开尔文表述，以为它禁止热变功。它禁止的是从单一热源吸热并全部变成功而不引起其它变化。
  \item 误解克劳修斯表述，以为热不能从低温传到高温。冰箱可以做到，但必须消耗外界功。
  \item 卡诺效率中直接代入摄氏温度。必须使用开尔文温度。
  \item 计算不可逆过程熵变时，直接沿不可逆实际路径积分 \(\int \dbar Q/T\)。正确做法是构造可逆路径。
  \item 认为绝热过程一定等熵。只有可逆绝热过程才等熵。
  \item 认为系统熵不能减少。熵增加原理约束的是孤立系统总熵。
  \item 混淆 \(Q/T\) 和 \(\Delta S\)。只有可逆等温过程中，才可直接写 \(\Delta S=Q_{\mathrm{rev}}/T\)。
  \item 忽略外界熵变。判断自发方向要看系统加外界的总熵变。
  \item 在克劳修斯不等式中随便使用系统内部温度。不可逆过程中系统未必有统一温度。
\end{PitfallList}
\end{WarnBox}

\section{第三章复习路线}

\begin{MethodBox}[建议背诵顺序]
\begin{ReviewList}
  \item 先理解自然过程方向性：热传导、自由膨胀、摩擦生热、扩散。
  \item 再背可逆过程条件：准静态且无耗散。
  \item 再背第二定律两种表述：开尔文表述、克劳修斯表述。
  \item 再背卡诺定理：相同热源之间，可逆热机效率最大且相等。
  \item 再背克劳修斯关系：
  \[
    \oint_{\mathrm{rev}}\frac{\dbar Q}{T}=0,
    \qquad
    \oint\frac{\dbar Q}{T}\le0.
  \]
  \item 再背熵定义：
  \[
    \dd S=\frac{\dbar Q_{\mathrm{rev}}}{T}.
  \]
  \item 再背熵增加原理：
  \[
    \Delta S_{\mathrm{isolated}}\ge0.
  \]
  \item 最后练理想气体熵变、自由膨胀、有限温差传热和卡诺效率。
\end{ReviewList}
\end{MethodBox}

\begin{KeyBox}[第三章一句话总结]
第二定律说明自然过程有方向，熵把这种方向写成态函数形式。实际计算时，最重要的判断是：若过程不可逆，熵变仍然要沿同一初末态之间的可逆路径计算；若判断过程是否自发，要看系统与外界组成的孤立整体总熵是否增加。
\end{KeyBox}
\chapter{麦克斯韦--玻尔兹曼分布律}

\chaptergoal{
本章从微观统计角度解释理想气体的宏观性质。复习核心不是记住一堆孤立公式，而是掌握三条线：第一，分子无规则运动如何给出压强和温度；第二，麦克斯韦分布如何描述分子速度和速率；第三，玻尔兹曼分布和能量均分定理如何把外力场、自由度和热容量联系起来。
}

\exammap{
\begin{ReviewList}
  \item \textbf{会推导图像：}理想气体压强公式 \(p=\frac13 nm\overline{v^2}\)，温度统计解释 \(\overline{\varepsilon_k}=\frac32 k_{\mathrm B}T\)。
  \item \textbf{会区分：}速度分布与速率分布；\(v_x\) 可以正负，速率 \(v\ge0\)。
  \item \textbf{会使用：}麦克斯韦速率分布函数 \(f(v)\)，以及最概然速率、平均速率、方均根速率。
  \item \textbf{会判断：}温度升高或分子质量变大时，分布曲线如何变化。
  \item \textbf{会应用：}碰壁数、泻流现象、麦克斯韦分布实验验证的基本物理图像。
  \item \textbf{会使用：}玻尔兹曼分布律解释重力场中的大气密度和压强随高度变化。
  \item \textbf{会计算：}能量均分定理给出的理想气体摩尔热容，并知道经典理论的局限。
\end{ReviewList}
}

\section{气体分子动理论与理想气体压强公式}

\subsection{理想气体的微观模型}

气体分子动理论从分子热运动出发，研究大量分子组成系统的热学性质。理想气体的微观模型包括：

\begin{KeyBox}[理想气体微观模型]
\begin{ReviewList}
  \item 分子本身大小远小于分子间平均距离，可忽略分子固有体积。
  \item 除碰撞瞬间外，分子间相互作用力可忽略。
  \item 分子间、分子与器壁之间的碰撞视为弹性碰撞。
  \item 大量分子的个别运动无规则，但整体遵守稳定的统计规律。
\end{ReviewList}
\end{KeyBox}

\begin{WarnBox}[微观无规则不等于宏观无规律]
单个分子运动具有偶然性，无法准确预言某个分子的瞬时速度。但大量分子的统计平均值具有确定规律，这正是气体动理论能够解释压强、温度和热容的原因。
\end{WarnBox}

\subsection{理想气体压强公式}

设单位体积内分子数为 \(n=N/V\)，每个分子质量为 \(m\)，分子速度为
\[
  \bm v=(v_x,v_y,v_z),
  \qquad
  v^2=v_x^2+v_y^2+v_z^2.
\]

由于气体处于平衡态，各方向等价，有
\[
  \overline{v_x^2}
  =
  \overline{v_y^2}
  =
  \overline{v_z^2}
  =
  \frac13\overline{v^2}.
\]

气体压强来自分子碰撞器壁时对单位面积、单位时间的冲量。统计平均后得到
\[
  p=nm\overline{v_x^2}.
\]
于是
\[
  \boxed{
  p=\frac13 nm\overline{v^2}
  }.
\]

如果记分子平均平动动能为
\[
  \overline{\varepsilon_k}
  =
  \frac12m\overline{v^2},
\]
则
\[
  \boxed{
  p=\frac23 n\overline{\varepsilon_k}
  }.
\]

\begin{KeyBox}[压强的微观意义]
气体压强不是由静止“重量”直接造成，而是大量分子持续撞击器壁产生的平均效果。分子数密度越大、分子平均平动动能越大，压强越大。
\end{KeyBox}

\subsection{温度的统计解释}

理想气体状态方程可写为
\[
  pV=Nk_{\mathrm B}T.
\]
因为 \(n=N/V\)，所以
\[
  p=nk_{\mathrm B}T.
\]

与压强公式
\[
  p=\frac23n\overline{\varepsilon_k}
\]
比较，得到
\[
  \boxed{
  \overline{\varepsilon_k}=\frac32 k_{\mathrm B}T
  }.
\]

也即
\[
  \boxed{
  T=\frac{2}{3k_{\mathrm B}}\overline{\varepsilon_k}
  }.
\]

\begin{KeyBox}[温度的统计意义]
温度反映大量分子无规则热运动平均平动动能的大小。温度越高，分子平均平动动能越大。
\end{KeyBox}

由
\[
  \frac12m\overline{v^2}=\frac32 k_{\mathrm B}T
\]
可得方均根速率
\[
  \boxed{
  v_{\mathrm{rms}}
  =
  \sqrt{\overline{v^2}}
  =
  \sqrt{\frac{3k_{\mathrm B}T}{m}}
  =
  \sqrt{\frac{3RT}{\mu}}
  }.
\]
其中 \(\mu\) 为摩尔质量。

\begin{WarnBox}[易错点：温度不是分子平均速率]
温度直接对应的是平均平动动能 \(\overline{\varepsilon_k}\)，也就是 \(\overline{v^2}\)，不是 \(\bar v\)。所以温度与 \(v_{\mathrm{rms}}\) 的平方成正比。
\end{WarnBox}

\section{统计规律与分布函数}

\subsection{统计规律}

\concept{统计规律}是大量偶然事件整体表现出来的确定规律。它有三个特点：
\begin{ReviewList}
  \item 描述大量粒子的综合性质，不描述单个粒子的个别行为。
  \item 在给定宏观条件下具有确定性。
  \item 伴随涨落；粒子数越多，相对涨落越小。
\end{ReviewList}

\subsection{概率与平均值}

若某事件出现次数为 \(N_i\)，总实验次数为 \(N\)，则概率为
\[
  P_i=\frac{N_i}{N}
  \qquad
  (N\to\infty).
\]
归一化条件为
\[
  \sum_i P_i=1.
\]

若物理量 \(M\) 的可能取值为 \(M_i\)，对应概率为 \(P_i\)，则统计平均值为
\[
  \boxed{
  \overline{M}=\sum_i M_iP_i
  }.
\]

若变量连续，概率密度为 \(f(x)\)，则
\[
  \int f(x)\dd x=1,
  \qquad
  \overline{M}=\int M(x)f(x)\dd x.
\]

\begin{KeyBox}[分布函数的读法]
若 \(f(v)\) 是速率分布函数，则
\[
  f(v)\dd v
\]
表示速率在 \(v\sim v+\dd v\) 范围内的分子数占总分子数的比例，而不是速率恰好等于 \(v\) 的概率。
\end{KeyBox}

\section{麦克斯韦速度分布律}

\subsection{速度分量分布}

平衡态理想气体中，一个速度分量 \(v_x\) 的分布为
\[
  \boxed{
  f(v_x)=
  \left(\frac{m}{2\pi k_{\mathrm B}T}\right)^{1/2}
  \exp\left(-\frac{mv_x^2}{2k_{\mathrm B}T}\right)
  }.
\]

它满足归一化条件：
\[
  \int_{-\infty}^{+\infty}f(v_x)\dd v_x=1.
\]

由于 \(v_x\) 可正可负，且正负方向等价：
\[
  \overline{v_x}=0,
  \qquad
  \overline{v_x^2}=\frac{k_{\mathrm B}T}{m}.
\]

\subsection{三维速度分布}

三维速度分布函数为
\[
  \boxed{
  f(v_x,v_y,v_z)=
  \left(\frac{m}{2\pi k_{\mathrm B}T}\right)^{3/2}
  \exp\left[
  -\frac{m(v_x^2+v_y^2+v_z^2)}{2k_{\mathrm B}T}
  \right]
  }.
\]

归一化条件为
\[
  \int_{-\infty}^{+\infty}
  \int_{-\infty}^{+\infty}
  \int_{-\infty}^{+\infty}
  f(v_x,v_y,v_z)
  \dd v_x\dd v_y\dd v_z
  =
  1.
\]

速度在体积元
\[
  \dd v_x\dd v_y\dd v_z
\]
内的分子数为
\[
  \dd N=Nf(v_x,v_y,v_z)\dd v_x\dd v_y\dd v_z.
\]

\begin{WarnBox}[速度分布与速率分布不要混]
速度 \(\bm v\) 是矢量，有方向；速率 \(v=|\bm v|\) 是标量，只有非负值。三维速度分布的变量是 \(v_x,v_y,v_z\)，速率分布的变量是 \(v\)。
\end{WarnBox}

\section{麦克斯韦速率分布律}

\subsection{速率分布函数}

速率在 \(v\sim v+\dd v\) 之间的分子数占总分子数的比例为
\[
  f(v)\dd v.
\]
麦克斯韦速率分布函数为
\[
  \boxed{
  f(v)=
  4\pi
  \left(\frac{m}{2\pi k_{\mathrm B}T}\right)^{3/2}
  v^2
  \exp\left(-\frac{mv^2}{2k_{\mathrm B}T}\right)
  }.
\]

对应分子数为
\[
  \boxed{
  \dd N=Nf(v)\dd v
  }.
\]

归一化条件为
\[
  \int_0^{+\infty}f(v)\dd v=1.
\]

\begin{KeyBox}[速率分布的两个因素]
\[
  f(v)\propto v^2\exp\left(-\frac{mv^2}{2k_{\mathrm B}T}\right).
\]
其中 \(v^2\) 来自速度空间中半径为 \(v\) 的球壳面积，倾向于让较大速率占比增大；指数因子来自能量概率，倾向于抑制高速分子。二者共同决定分布曲线存在一个峰值。
\end{KeyBox}

\subsection{三种特征速率}

\subsubsection{最概然速率}

\concept{最概然速率} \(v_{\mathrm p}\) 是使 \(f(v)\) 取最大值的速率：
\[
  \boxed{
  v_{\mathrm p}
  =
  \sqrt{\frac{2k_{\mathrm B}T}{m}}
  =
  \sqrt{\frac{2RT}{\mu}}
  }.
\]

\subsubsection{平均速率}

\concept{平均速率}定义为
\[
  \bar v=\int_0^\infty vf(v)\dd v.
\]
结果为
\[
  \boxed{
  \bar v
  =
  \sqrt{\frac{8k_{\mathrm B}T}{\pi m}}
  =
  \sqrt{\frac{8RT}{\pi\mu}}
  }.
\]

\subsubsection{方均根速率}

\concept{方均根速率}为
\[
  v_{\mathrm{rms}}=\sqrt{\overline{v^2}}.
\]
结果为
\[
  \boxed{
  v_{\mathrm{rms}}
  =
  \sqrt{\frac{3k_{\mathrm B}T}{m}}
  =
  \sqrt{\frac{3RT}{\mu}}
  }.
\]

三者大小关系：
\[
  \boxed{
  v_{\mathrm p}<\bar v<v_{\mathrm{rms}}
  }.
\]

\begin{WarnBox}[易错点：三个速率含义不同]
\begin{itemize}
  \item \(v_{\mathrm p}\)：分布函数最大处的速率，不是平均值。
  \item \(\bar v\)：速率的一次平均，用于碰壁数、平均运动距离等。
  \item \(v_{\mathrm{rms}}\)：与平均平动动能直接相关。
\end{itemize}
不能在公式中随意互换。
\end{WarnBox}

\subsection{温度和摩尔质量对分布曲线的影响}

麦克斯韦速率分布随 \(T\) 和 \(m\) 变化：

\begin{KeyBox}[曲线变化规律]
\begin{itemize}
  \item 温度升高：\(v_{\mathrm p},\bar v,v_{\mathrm{rms}}\) 都增大，分布峰向右移，曲线变宽，峰值降低。
  \item 分子质量增大：三个特征速率都减小，分布峰向左移，曲线变窄，峰值升高。
  \item 曲线下面积恒为 \(1\)，因为总概率不变。
\end{itemize}
\end{KeyBox}

\section{麦克斯韦分布律的应用}

\subsection{由速度分布求理想气体状态方程}

麦克斯韦速度分布给出
\[
  \overline{v_x^2}=\frac{k_{\mathrm B}T}{m}.
\]
代入压强公式
\[
  p=nm\overline{v_x^2},
\]
得
\[
  p=nk_{\mathrm B}T.
\]
由于 \(n=N/V\)，所以
\[
  pV=Nk_{\mathrm B}T.
\]
这正是理想气体状态方程。

\begin{KeyBox}[这一步的意义]
第一章把 \(pV=Nk_{\mathrm B}T\) 当作状态方程使用；第四章从分子速度分布出发说明它的微观来源：压强来自分子碰壁，温度来自平均平动动能。
\end{KeyBox}

\subsection{碰壁数}

\concept{碰壁数} \(G\) 是单位时间内碰撞单位器壁面积的分子数。对平衡理想气体，
\[
  \boxed{
  G=\frac14 n\bar v
  }.
\]

若器壁面积为 \(A\)，单位时间内碰到该面积的分子数为
\[
  \Delta N=GA\Delta t
  =
  \frac14 n\bar v A\Delta t.
\]

\begin{WarnBox}[为什么是 \(\bar v\)，不是 \(v_{\mathrm{rms}}\)]
碰壁数统计的是单位时间内有多少分子穿过某个面积，直接涉及分子速度在法向上的一次平均，因此结果与平均速率 \(\bar v\) 相关，而不是方均根速率。
\end{WarnBox}

\subsection{泻流现象}

\concept{泻流}是指气体通过小孔从容器中逸出的现象。当小孔足够小，孔附近不破坏容器内平衡分布时，单位时间通过面积 \(A\) 泻出的分子数为
\[
  \boxed{
  \frac{\Delta N}{\Delta t}
  =
  \frac14 n\bar v A
  }.
\]

由于
\[
  \bar v\propto \frac{1}{\sqrt{\mu}},
\]
在相同温度和数密度下，轻分子泻出更快。

若两种气体温度相同、孔两边条件相同，则泻流速率近似满足
\[
  \frac{\Gamma_1}{\Gamma_2}
  =
  \sqrt{\frac{\mu_2}{\mu_1}}.
\]

\begin{KeyBox}[泻流的物理图像]
泻流不是普通宏观气流。普通气流由压强差造成整体定向运动；泻流强调分子热运动使部分分子通过小孔逸出。高速分子更容易通过孔口，因此泻出的分子束速率分布会偏向高速。
\end{KeyBox}

\subsection{麦克斯韦分布律的实验验证}

麦克斯韦分布律可通过分子束实验验证。基本思想是：
\begin{ReviewList}
  \item 从高温原子炉或分子源中产生分子束。
  \item 利用狭缝准直，使分子束方向较为确定。
  \item 用速度选择装置筛选不同速率的分子。
  \item 改变选择条件并检测分子强度，从而得到速率分布。
\end{ReviewList}

密勒--库士实验使用旋转选择器验证了麦克斯韦分布律的正确性。授课 PDF 中也把该实验作为麦克斯韦分布律应用与验证的一部分。:contentReference[oaicite:2]{index=2}

\section{玻尔兹曼分布律}

\subsection{外力场中的分子数密度}

如果气体分子处在外力场中，不同空间位置具有不同势能 \(\varepsilon_p(\bm r)\)。平衡态下，分子数密度满足\concept{玻尔兹曼分布律}：
\[
  \boxed{
  n(\bm r)
  =
  n_0
  \exp\left[
  -\frac{\varepsilon_p(\bm r)}{k_{\mathrm B}T}
  \right]
  }.
\]

更一般地，在相空间中，分子分布与总能量有关：
\[
  \dd N
  \propto
  \exp\left[
  -\frac{\varepsilon_k+\varepsilon_p}{k_{\mathrm B}T}
  \right]
  \dd x\dd y\dd z\dd v_x\dd v_y\dd v_z.
\]

\begin{KeyBox}[玻尔兹曼分布的核心]
势能越高的位置，分子出现概率越小；温度越高，热运动越强，分子越不容易被外力场造成的势能差限制。
\end{KeyBox}

\subsection{重力场中的大气压强公式}

在重力场中，取高度 \(z=0\) 处势能为零，则
\[
  \varepsilon_p=mgz.
\]
于是
\[
  \boxed{
  n(z)=n_0\exp\left(-\frac{mgz}{k_{\mathrm B}T}\right)
  }.
\]

因为理想气体满足
\[
  p=nk_{\mathrm B}T,
\]
若温度近似不随高度变化，则
\[
  \boxed{
  p(z)=p_0\exp\left(-\frac{mgz}{k_{\mathrm B}T}\right)
  }.
\]

对一摩尔气体，\(m=\mu/N_A\)，且 \(R=N_Ak_{\mathrm B}\)，所以也可写为
\[
  \boxed{
  p(z)=p_0\exp\left(-\frac{\mu gz}{RT}\right)
  }.
\]

\begin{WarnBox}[大气压强公式的条件]
该公式默认大气近似等温，并把空气视为理想气体。真实大气温度随高度变化，所以实际大气压强随高度的关系会更复杂。
\end{WarnBox}

\subsection{离心力场中的分布}

在角速度为 \(\omega\) 的旋转系统中，分子受到离心作用。相对于转轴距离为 \(r\) 处，平衡分布可写为
\[
  \boxed{
  n(r)
  =
  n(0)
  \exp\left(
  \frac{m\omega^2r^2}{2k_{\mathrm B}T}
  \right)
  }.
\]

因此分子质量越大、角速度越大，向外侧聚集越明显。这是离心分离的基本物理图像。

\begin{KeyBox}[重力场与离心场的对比]
\begin{itemize}
  \item 重力场中，势能 \(mgz\) 随高度升高而增大，所以密度随高度指数下降。
  \item 离心场中，有效势能随 \(r\) 增大而降低，所以密度随半径增大而升高。
\end{itemize}
\end{KeyBox}

\section{能量均分定理}

\subsection{自由度}

\concept{自由度}是确定一个物体空间位置或运动状态所需的独立坐标数。对分子热运动，常见自由度包括：
\begin{itemize}
  \item 平动自由度 \(t\)；
  \item 转动自由度 \(r\)；
  \item 振动自由度 \(s\)。
\end{itemize}

对理想气体分子：
\[
  t=3.
\]
单原子分子通常只有平动自由度；双原子分子在常温下通常考虑 \(3\) 个平动自由度和 \(2\) 个转动自由度，振动自由度常因量子效应未被充分激发而暂不计入。

\subsection{能量按自由度均分定理}

\concept{能量均分定理}：在温度为 \(T\) 的平衡状态下，分子每一个以平方项形式出现在能量中的自由度，平均具有
\[
  \boxed{
  \frac12 k_{\mathrm B}T
  }
\]
的能量。

例如平动动能为
\[
  \varepsilon_k=
  \frac12mv_x^2+
  \frac12mv_y^2+
  \frac12mv_z^2.
\]
因此
\[
  \overline{\frac12mv_x^2}
  =
  \overline{\frac12mv_y^2}
  =
  \overline{\frac12mv_z^2}
  =
  \frac12k_{\mathrm B}T.
\]

总平动平均动能为
\[
  \overline{\varepsilon_k}
  =
  \frac32k_{\mathrm B}T.
\]

\begin{WarnBox}[能量均分定理是统计规律]
能量均分定理不是说每个分子在任一瞬间各自由度能量都相等。它描述的是大量分子在平衡态下的统计平均结果。
\end{WarnBox}

\subsection{理想气体内能与热容量}

若一个理想气体分子共有 \(i\) 个被激发的自由度，则单个分子平均能量为
\[
  \bar\varepsilon=\frac{i}{2}k_{\mathrm B}T.
\]
\(N\) 个分子的内能为
\[
  U=N\bar\varepsilon
  =
  \frac{i}{2}Nk_{\mathrm B}T.
\]
对 \(1\) mol 理想气体，
\[
  U_m=\frac{i}{2}RT.
\]

因此摩尔定容热容为
\[
  \boxed{
  C_{V,m}=\left(\frac{\partial U_m}{\partial T}\right)_V
  =
  \frac{i}{2}R
  }.
\]
由迈耶公式
\[
  C_{p,m}=C_{V,m}+R,
\]
得
\[
  \boxed{
  C_{p,m}=\frac{i+2}{2}R
  }.
\]
热容比为
\[
  \boxed{
  \gamma=\frac{C_{p,m}}{C_{V,m}}=\frac{i+2}{i}
  }.
\]

\subsection{常见理想气体热容量}

\begin{FormulaBox}[经典能量均分给出的热容量]
\begin{tabularx}{\textwidth}{@{}>{\bfseries}lcccX@{}}
\toprule
气体类型 & 有效自由度 \(i\) & \(C_{V,m}\) & \(C_{p,m}\) & \(\gamma\)\\
\midrule
单原子分子 & \(3\) & \(\frac32R\) & \(\frac52R\) & \(\frac53\)\\[0.4em]
双原子分子，常温 & \(5\) & \(\frac52R\) & \(\frac72R\) & \(\frac75\)\\[0.4em]
多原子非线性分子，忽略振动 & \(6\) & \(3R\) & \(4R\) & \(\frac43\)\\
\bottomrule
\end{tabularx}
\end{FormulaBox}

\subsection{经典理论解释热容量的困难}

经典能量均分定理预言：只要存在某种自由度，它就应该平均分得 \(\frac12k_{\mathrm B}T\) 的能量。但实验表明，分子的转动、振动自由度是否参与热容量，与温度有关。

例如双原子分子：
\begin{itemize}
  \item 低温下，转动和振动自由度可能未充分激发；
  \item 常温下，通常主要有 \(3\) 个平动自由度和 \(2\) 个转动自由度；
  \item 高温下，振动自由度逐渐被激发，热容量增大。
\end{itemize}

\begin{KeyBox}[经典理论的局限]
经典能量均分定理无法解释热容量随温度变化的细节。根本原因是微观能级量子化：某些自由度在低温下不能连续吸收任意小能量，因此不会按经典理论贡献热容量。
\end{KeyBox}

\section{第四章题型模板}

\subsection{题型一：由温度求分子速率}

\begin{MethodBox}[三种速率计算]
给定温度 \(T\) 与摩尔质量 \(\mu\)，直接使用：
\[
  v_{\mathrm p}=\sqrt{\frac{2RT}{\mu}},
  \qquad
  \bar v=\sqrt{\frac{8RT}{\pi\mu}},
  \qquad
  v_{\mathrm{rms}}=\sqrt{\frac{3RT}{\mu}}.
\]
注意 \(\mu\) 必须使用 \(\mathrm{kg/mol}\)，不能直接用 \(\mathrm{g/mol}\)。
\end{MethodBox}

\subsection{题型二：判断分布曲线变化}

\begin{MethodBox}[麦克斯韦曲线判断]
\begin{ReviewList}
  \item 温度升高：峰右移、变低、变宽。
  \item 温度降低：峰左移、变高、变窄。
  \item 摩尔质量增大：峰左移、变高、变窄。
  \item 摩尔质量减小：峰右移、变低、变宽。
  \item 不管怎么变，曲线下面积都等于 \(1\)。
\end{ReviewList}
\end{MethodBox}

\subsection{题型三：玻尔兹曼分布应用}

\begin{MethodBox}[重力场大气公式]
如果题目给出温度近似恒定，且要求高度 \(z\) 处压强：
\[
  p(z)=p_0\exp\left(-\frac{\mu gz}{RT}\right).
\]
若问数密度：
\[
  n(z)=n_0\exp\left(-\frac{mgz}{k_{\mathrm B}T}\right).
\]
单个分子质量 \(m\) 与摩尔质量 \(\mu\) 的关系是
\[
  m=\frac{\mu}{N_A}.
\]
\end{MethodBox}

\subsection{题型四：能量均分与热容量}

\begin{MethodBox}[自由度算热容]
\begin{ReviewList}
  \item 先判断有效自由度 \(i\)。
  \item 用
  \[
    C_{V,m}=\frac{i}{2}R.
  \]
  \item 再用迈耶公式
  \[
    C_{p,m}=C_{V,m}+R.
  \]
  \item 最后算
  \[
    \gamma=\frac{C_{p,m}}{C_{V,m}}.
  \]
\end{ReviewList}
\end{MethodBox}

\section{第四章公式速查}

\formulatable{
\begin{tabularx}{\textwidth}{@{}>{\bfseries}lXl@{}}
\toprule
内容 & 公式 & 条件\\
\midrule

理想气体压强公式 &
\(\displaystyle p=\frac13nm\overline{v^2}=\frac23n\overline{\varepsilon_k}\)
& 平衡理想气体\\[0.8em]

温度统计解释 &
\(\displaystyle \overline{\varepsilon_k}=\frac32k_{\mathrm B}T\)
& 平动平均动能\\[0.8em]

方均根速率 &
\(\displaystyle v_{\mathrm{rms}}=\sqrt{\frac{3k_{\mathrm B}T}{m}}=\sqrt{\frac{3RT}{\mu}}\)
& 理想气体\\[0.8em]

速度分量分布 &
\(\displaystyle f(v_x)=\left(\frac{m}{2\pi k_{\mathrm B}T}\right)^{1/2}\exp\left(-\frac{mv_x^2}{2k_{\mathrm B}T}\right)\)
& \(v_x\in(-\infty,\infty)\)\\[1em]

速率分布 &
\(\displaystyle f(v)=4\pi\left(\frac{m}{2\pi k_{\mathrm B}T}\right)^{3/2}v^2\exp\left(-\frac{mv^2}{2k_{\mathrm B}T}\right)\)
& \(v\ge0\)\\[1em]

最概然速率 &
\(\displaystyle v_{\mathrm p}=\sqrt{\frac{2k_{\mathrm B}T}{m}}=\sqrt{\frac{2RT}{\mu}}\)
& \(f(v)\) 最大\\[0.8em]

平均速率 &
\(\displaystyle \bar v=\sqrt{\frac{8k_{\mathrm B}T}{\pi m}}=\sqrt{\frac{8RT}{\pi\mu}}\)
& 一次平均\\[0.8em]

三速率大小 &
\(\displaystyle v_{\mathrm p}<\bar v<v_{\mathrm{rms}}\)
& 同一气体同一温度\\[0.8em]

碰壁数 &
\(\displaystyle G=\frac14n\bar v\)
& 单位面积单位时间\\[0.8em]

泻流分子数率 &
\(\displaystyle \frac{\Delta N}{\Delta t}=\frac14n\bar v A\)
& 小孔泻流\\[0.8em]

玻尔兹曼分布 &
\(\displaystyle n(\bm r)=n_0\exp\left[-\frac{\varepsilon_p(\bm r)}{k_{\mathrm B}T}\right]\)
& 外力场平衡态\\[1em]

重力场分布 &
\(\displaystyle n(z)=n_0e^{-mgz/k_{\mathrm B}T},\quad p(z)=p_0e^{-\mu gz/RT}\)
& 等温大气\\[1em]

能量均分 &
\(\displaystyle \overline{\varepsilon_i}=\frac12k_{\mathrm B}T\)
& 每个平方型自由度\\[0.8em]

摩尔定容热容 &
\(\displaystyle C_{V,m}=\frac{i}{2}R\)
& \(i\) 个有效自由度\\[0.8em]

摩尔定压热容 &
\(\displaystyle C_{p,m}=\frac{i+2}{2}R\)
& 理想气体\\[0.8em]

热容比 &
\(\displaystyle \gamma=\frac{i+2}{i}\)
& 理想气体\\

\bottomrule
\end{tabularx}
}

\section{第四章易错点总表}

\begin{WarnBox}[本章高频错误]
\begin{PitfallList}
  \item 把速度分布和速率分布混为一谈。速度是矢量，速率是标量。
  \item 忘记速率分布的 \(4\pi v^2\) 因子。它来自速度空间球壳面积。
  \item 认为最概然速率就是平均速率。实际上 \(v_{\mathrm p}<\bar v<v_{\mathrm{rms}}\)。
  \item 用 \(\bar v\) 计算平均动能。平均动能对应 \(\overline{v^2}\)，不是 \(\bar v^2\)。
  \item 温度升高时只记得曲线右移，忘记峰值降低、曲线变宽、面积不变。
  \item 计算速率时把摩尔质量 \(\mu\) 的单位用成 \(\mathrm{g/mol}\)。必须换成 \(\mathrm{kg/mol}\)。
  \item 把碰壁数中的 \(\bar v\) 写成 \(v_{\mathrm{rms}}\)。
  \item 玻尔兹曼分布中指数符号写反。势能越高，粒子数密度越小。
  \item 大气压强公式默认等温大气；真实大气不严格满足这个条件。
  \item 能量均分定理是统计平均规律，不适用于单个分子瞬时状态。
  \item 经典热容量公式忽略了量子效应，不能解释所有温区的实验结果。
\end{PitfallList}
\end{WarnBox}

\section{第四章复习路线}

\begin{MethodBox}[建议背诵顺序]
\begin{ReviewList}
  \item 先背压强公式：
  \[
    p=\frac13nm\overline{v^2}.
  \]
  \item 再背温度统计解释：
  \[
    \overline{\varepsilon_k}=\frac32k_{\mathrm B}T.
  \]
  \item 再区分速度分布和速率分布。
  \item 再背麦克斯韦速率分布和三个特征速率。
  \item 再看应用：碰壁数、泻流、实验验证。
  \item 再背玻尔兹曼分布：
  \[
    n\propto e^{-\varepsilon_p/k_{\mathrm B}T}.
  \]
  \item 最后用能量均分定理推热容量：
  \[
    C_{V,m}=\frac{i}{2}R,\qquad C_{p,m}=\frac{i+2}{2}R.
  \]
\end{ReviewList}
\end{MethodBox}

\begin{KeyBox}[第四章一句话总结]
第四章把宏观热学量翻译成微观统计语言：压强来自分子碰壁，温度来自平均平动动能，麦克斯韦分布描述速度概率，玻尔兹曼分布描述外力场中的空间概率，能量均分定理把自由度与热容量连接起来。
\end{KeyBox}
\chapter{气体中的输运过程}

\chaptergoal{
本章讨论气体处于非平衡态时，由分子热运动和分子间碰撞导致的宏观输运现象。复习核心是把三种过程统一起来看：温度不均匀导致热量输运，流速不均匀导致动量输运，组分数密度不均匀导致分子数输运。宏观上对应傅里叶定律、牛顿黏性定律、菲克定律；微观上对应平均自由程、平均速率和分子碰撞。
}

\exammap{
\begin{ReviewList}
  \item \textbf{会识别：}热传导、黏滞、扩散分别输运什么物理量。
  \item \textbf{会写宏观规律：}傅里叶热传导定律、牛顿黏性定律、菲克扩散定律。
  \item \textbf{会解释负号：}输运总是沿着使不均匀性减弱的方向进行，即沿负梯度方向。
  \item \textbf{会用微观图像：}分子热运动造成跨层转移，分子碰撞限制输运强弱。
  \item \textbf{会背输运系数：}\(\kappa=\frac13\lambda\rho c_V\bar v\)，\(\eta=\frac13\lambda\rho\bar v\)，\(D=\frac13\lambda\bar v\)。
  \item \textbf{会判断压强依赖：}普通稀薄气体中 \(\kappa,\eta\) 近似与压强无关，\(D\) 与压强近似成反比。
  \item \textbf{会辨析：}极稀薄气体中平均自由程可与容器尺寸相比，普通输运系数公式不再适用。
\end{ReviewList}
}

\section{输运现象的基本图像}

\subsection{什么是输运过程}

\concept{输运现象}是指系统处于非平衡态时，由于某种物理量空间分布不均匀，导致该物理量出现宏观流动的现象。气体中的典型输运过程包括：
\begin{ReviewList}
  \item \concept{热传导过程}：温度不均匀导致热量输运。
  \item \concept{黏滞过程}：宏观定向流速不均匀导致动量输运。
  \item \concept{扩散过程}：组分数密度不均匀导致分子数输运。
\end{ReviewList}

\begin{KeyBox}[输运过程的一句话理解]
输运过程的共同作用是消除不均匀性：温度差趋于消失，流速差趋于消失，浓度差趋于消失。
\end{KeyBox}

\subsection{三种输运现象的类比}

\begin{FormulaBox}[三种输运过程的统一结构]
\begin{tabularx}{\textwidth}{@{}>{\bfseries}lCCC@{}}
\toprule
类型 & 热传导 & 黏滞 & 扩散\\
\midrule
不均匀量 & 温度 \(T\) & 定向流速 \(u\) & 数密度 \(n\)\\
梯度 & \(\dd T/\dd z\) & \(\dd u/\dd z\) & \(\dd n/\dd z\)\\
输运量 & 热量 \(Q\) & 定向动量 \(p_x\) & 分子数 \(N\)\\
输运系数 & \(\kappa\) & \(\eta\) & \(D\)\\
宏观规律 & 傅里叶定律 & 牛顿黏性定律 & 菲克定律\\
\bottomrule
\end{tabularx}
\end{FormulaBox}

三个输运系数均为正值。公式中的负号表示输运方向与对应物理量梯度方向相反，即沿着使不均匀性减弱的方向输运。

\begin{WarnBox}[易错点：梯度方向和输运方向相反]
若 \(T\) 随 \(z\) 增大而增大，则 \(\dd T/\dd z>0\)，热量沿 \(-z\) 方向流动。负号不是数学装饰，而是物理方向。
\end{WarnBox}

\section{热传导的宏观规律}

\subsection{热传导现象}

\concept{热传导}是指系统内部不同部分之间，或系统与外界之间存在温度差时，热量从温度较高处传递到温度较低处的现象。

考虑气体夹在两个恒温热源之间，上方温度为 \(T_1\)，下方温度为 \(T_2\)，且 \(T_1>T_2\)。若沿 \(z\) 方向存在温度梯度，则通过垂直于 \(z\) 轴的面积 \(\Delta S\) 的热量遵循傅里叶定律。

\subsection{傅里叶热传导定律}

\begin{FormulaBox}[傅里叶定律]
\[
  \boxed{
  \frac{\Delta Q}{\Delta t}
  =
  -\kappa
  \left(\frac{\dd T}{\dd z}\right)_{z_0}
  \Delta S
  }
\]
或写成热流密度：
\[
  \boxed{
  j_Q
  =
  \frac{1}{\Delta S}\frac{\Delta Q}{\Delta t}
  =
  -\kappa\frac{\dd T}{\dd z}
  }.
\]
\end{FormulaBox}

其中 \(\kappa\) 称为\concept{热传导系数}，单位为
\[
  \mathrm{W\cdot m^{-1}\cdot K^{-1}}.
\]
\(\kappa\) 越大，物质导热能力越强。

\begin{KeyBox}[傅里叶定律的物理含义]
单位时间通过单位面积输运的热量，与温度梯度成正比；热量从高温处流向低温处。
\end{KeyBox}

\subsection{平板稳态导热}

若两块平行板之间距离为 \(L\)，温度分别为 \(T_1,T_2\)，并且温度梯度近似均匀，则
\[
  \frac{\dd T}{\dd z}\approx \frac{T_2-T_1}{L}.
\]
热流率大小为
\[
  \boxed{
  \frac{\Delta Q}{\Delta t}
  =
  \kappa\frac{T_1-T_2}{L}\Delta S
  }.
\]

若导热持续时间为 \(\Delta t\)，则传递热量为
\[
  \boxed{
  \Delta Q
  =
  \kappa\frac{T_1-T_2}{L}\Delta S\,\Delta t
  }.
\]

\begin{WarnBox}[热传导题常见错误]
如果题目问的是“热量”，要乘以时间 \(\Delta t\)；如果问的是“热流率”或“单位时间传热量”，不要再乘时间。
\end{WarnBox}

\section{黏滞现象的宏观规律}

\subsection{黏滞现象}

\concept{黏滞现象}是指流体内部不同流层流速不同时，相邻流层之间出现切向相互作用的现象。其结果是：快流层被减速，慢流层被加速，流速不均匀性逐渐减弱。

典型模型是两块平行板之间夹有气体。下板静止，上板以速度 \(u_1\) 沿 \(x\) 方向运动。气体流速 \(u\) 随垂直方向 \(z\) 改变，因此存在速度梯度
\[
  \frac{\dd u}{\dd z}.
\]

\subsection{牛顿黏性定律}

\begin{FormulaBox}[牛顿黏性定律]
若在 \(z=z_0\) 处取面积为 \(\Delta S\) 的流体层，则该层间的黏滞力为
\[
  \boxed{
  f
  =
  -\eta
  \left(\frac{\dd u}{\dd z}\right)_{z_0}
  \Delta S
  }.
\]
也可写为切应力：
\[
  \boxed{
  \tau
  =
  \frac{f}{\Delta S}
  =
  -\eta\frac{\dd u}{\dd z}
  }.
\]
\end{FormulaBox}

其中 \(\eta\) 称为\concept{黏滞系数}或动力黏度，单位为
\[
  \mathrm{Pa\cdot s}.
\]

\begin{KeyBox}[黏滞的物理含义]
黏滞力本质上对应定向动量的输运。分子从快流层跑到慢流层，会把较大的定向动量带过去；分子从慢流层跑到快流层，会把较小的定向动量带过去。结果是快层减速，慢层加速。
\end{KeyBox}

\subsection{平板间线性速度分布}

若两平行板间距为 \(L\)，上板速度为 \(u_1\)，下板静止，且速度分布近似线性，则
\[
  \frac{\dd u}{\dd z}\approx \frac{u_1}{L}.
\]
黏滞力大小为
\[
  \boxed{
  f
  =
  \eta\frac{u_1}{L}\Delta S
  }.
\]

\begin{WarnBox}[黏滞力的方向]
公式中的负号表示黏滞力总是阻碍相对运动。做题时若只求大小，可取
\[
  |f|=\eta\left|\frac{\dd u}{\dd z}\right|\Delta S.
\]
若要求方向，必须结合流速梯度判断。
\end{WarnBox}

\section{扩散现象的宏观规律}

\subsection{扩散现象}

\concept{扩散}是指由于组分分子数密度不均匀，分子从数密度较高处向数密度较低处迁移的现象。扩散的结果是组分分布趋于均匀。

设某种组分的数密度为 \(n(z)\)，沿 \(z\) 方向存在数密度梯度：
\[
  \frac{\dd n}{\dd z}.
\]

\subsection{菲克定律}

\begin{FormulaBox}[菲克定律]
通过面积 \(\Delta S\) 的分子数输运率为
\[
  \boxed{
  \frac{\Delta N}{\Delta t}
  =
  -D
  \left(\frac{\dd n}{\dd z}\right)_{z_0}
  \Delta S
  }.
\]
分子数流密度为
\[
  \boxed{
  j_N
  =
  \frac{1}{\Delta S}\frac{\Delta N}{\Delta t}
  =
  -D\frac{\dd n}{\dd z}
  }.
\]
\end{FormulaBox}

其中 \(D\) 称为\concept{扩散系数}，单位为
\[
  \mathrm{m^2/s}.
\]

\begin{KeyBox}[扩散的物理含义]
扩散并不要求存在宏观风速。即使气体整体没有定向流动，分子无规则热运动仍会使高数密度区域的分子更多地穿过某一截面，从而形成净扩散流。
\end{KeyBox}

\subsection{扩散与热传导、黏滞的共同点}

\begin{FormulaBox}[三种宏观规律的统一形式]
\[
  \text{输运流密度}
  =
  -
  \text{输运系数}
  \times
  \text{对应物理量的梯度}.
\]
具体地：
\[
  j_Q=-\kappa\frac{\dd T}{\dd z},
  \qquad
  \tau=-\eta\frac{\dd u}{\dd z},
  \qquad
  j_N=-D\frac{\dd n}{\dd z}.
\]
\end{FormulaBox}

\begin{WarnBox}[扩散不是对流]
扩散来自分子无规则热运动导致的组分迁移；对流是流体整体宏观运动导致的物质输运。二者可以同时存在，但概念不同。
\end{WarnBox}

\section{气体输运过程的微观解释}

\subsection{微观机制}

气体输运过程来自两个微观因素：
\begin{ReviewList}
  \item \concept{分子的热运动}：使分子能够从一个空间区域移动到另一个区域，从而携带能量、动量或分子数穿过某一截面。
  \item \concept{分子间碰撞}：不断改变分子运动方向和携带信息的距离，决定输运过程强弱。
\end{ReviewList}

若没有分子热运动，就没有跨层转移；若没有分子间碰撞，分子会直接飞到器壁，普通气体输运模型也会失效。

\begin{KeyBox}[平均自由程的作用]
平均自由程 \(\lambda\) 表示分子两次相邻碰撞之间运动路程的平均值。它决定分子大致从多远的地方把能量、动量或数密度信息带到当前截面。
\end{KeyBox}

\subsection{平均自由程与碰撞频率}

对分子直径为 \(d\)、数密度为 \(n\) 的理想气体，平均自由程近似为
\[
  \boxed{
  \lambda=\frac{1}{\sqrt{2}\pi d^2n}
  }.
\]

平均碰撞频率可写为
\[
  \boxed{
  Z=\frac{\bar v}{\lambda}
  =
  \sqrt{2}\pi d^2n\bar v
  }.
\]

由理想气体状态方程
\[
  p=nk_{\mathrm B}T
\]
可知，在温度一定时，
\[
  n\propto p,
  \qquad
  \lambda\propto \frac1p.
\]

\begin{WarnBox}[平均自由程的压强依赖]
温度一定时，压强越低，分子数密度越小，平均自由程越大。高真空下，\(\lambda\) 可能变得和容器尺寸同量级，这时普通输运公式要谨慎使用。
\end{WarnBox}

\subsection{统一的输运公式}

设某物理量 \(W\) 的体密度可写成
\[
  nW.
\]
其中 \(W\) 表示单个分子携带的某种物理量，例如能量、定向动量或组分标记。分子热运动使这些量跨过面积 \(\Delta S\) 输运。经过平均自由程近似处理，可以得到统一形式：
\[
  \boxed{
  \frac{\Delta W_{\mathrm{total}}}{\Delta t}
  =
  -\frac13\bar v\lambda
  \frac{\dd (nW)}{\dd z}
  \Delta S
  }.
\]

三种输运过程都可看作这个公式的特例。

\begin{KeyBox}[为什么有 \(\frac13\)]
因分子运动在三维空间中各向同性，只有一部分速度分量有效地穿过给定截面。严格平均后出现数量级因子 \(\frac13\)。热学B复习中重点是理解它来自三维无规则运动的方向平均，不必深究复杂积分。
\end{KeyBox}

\section{三种输运系数的微观表达}

\subsection{热传导系数}

热传导中，分子携带的是热运动能量。单位质量定容比热为 \(c_V\)，质量密度为 \(\rho\)。分子能量密度与温度有关，可写成
\[
  \rho c_VT.
\]
由统一输运公式得到
\[
  \frac{\Delta Q}{\Delta t}
  =
  -\frac13\lambda\bar v\rho c_V
  \frac{\dd T}{\dd z}\Delta S.
\]
与傅里叶定律比较：
\[
  \frac{\Delta Q}{\Delta t}
  =
  -\kappa\frac{\dd T}{\dd z}\Delta S.
\]
得到
\[
  \boxed{
  \kappa=\frac13\lambda\rho c_V\bar v
  }.
\]

\begin{KeyBox}[热传导系数的含义]
\(\lambda\) 越大，分子能把能量带得越远；\(\bar v\) 越大，分子带能量穿过截面越快；\(\rho c_V\) 越大，单位体积单位温差可携带的热量越多。
\end{KeyBox}

\subsection{黏滞系数}

黏滞过程中，分子携带的是沿流动方向的定向动量。设宏观流速为 \(u(z)\)，单位体积内定向动量密度为
\[
  \rho u.
\]
由统一输运公式得到动量输运率：
\[
  \frac{\Delta p_x}{\Delta t}
  =
  -\frac13\lambda\bar v\rho
  \frac{\dd u}{\dd z}\Delta S.
\]
动量输运率对应层间切向力 \(f\)。与牛顿黏性定律比较：
\[
  f
  =
  -\eta\frac{\dd u}{\dd z}\Delta S.
\]
得到
\[
  \boxed{
  \eta=\frac13\lambda\rho\bar v
  }.
\]

\begin{KeyBox}[黏滞系数的含义]
黏滞不是因为分子“粘住”不动，而是因为分子热运动在不同流层之间输运定向动量。气体温度升高时，分子平均速率增大，气体黏滞性通常增强。
\end{KeyBox}

\subsection{扩散系数}

扩散过程中，分子携带的是组分本身。由统一公式可得
\[
  \frac{\Delta N}{\Delta t}
  =
  -\frac13\lambda\bar v
  \frac{\dd n}{\dd z}\Delta S.
\]
与菲克定律比较：
\[
  \frac{\Delta N}{\Delta t}
  =
  -D\frac{\dd n}{\dd z}\Delta S.
\]
得到
\[
  \boxed{
  D=\frac13\lambda\bar v
  }.
\]

\begin{KeyBox}[扩散系数的含义]
扩散快慢由两个量控制：分子跑得多快，即 \(\bar v\)；分子在改变方向前能跑多远，即 \(\lambda\)。
\end{KeyBox}

\section{输运系数与气体状态的关系}

\subsection{普通稀薄气体中的压强依赖}

对普通稀薄气体，
\[
  \lambda=\frac{1}{\sqrt{2}\pi d^2n}.
\]
又
\[
  \rho=mn.
\]
因此
\[
  \lambda\rho
  =
  \frac{1}{\sqrt{2}\pi d^2n}mn
  =
  \frac{m}{\sqrt{2}\pi d^2}.
\]
可见 \(\lambda\rho\) 与数密度 \(n\) 无关。

所以：
\[
  \kappa=\frac13\lambda\rho c_V\bar v
\]
近似与 \(n\) 和 \(p\) 无关；
\[
  \eta=\frac13\lambda\rho\bar v
\]
也近似与 \(n\) 和 \(p\) 无关。

而扩散系数
\[
  D=\frac13\lambda\bar v
\]
由于
\[
  \lambda\propto\frac1n,
\]
所以在温度一定时
\[
  \boxed{
  D\propto\frac1n\propto\frac1p
  }.
\]

\begin{FormulaBox}[压强依赖总结]
在普通稀薄气体、温度一定时：
\[
  \boxed{
  \kappa \ \text{近似与}\ p\ \text{无关}
  }
\]
\[
  \boxed{
  \eta \ \text{近似与}\ p\ \text{无关}
  }
\]
\[
  \boxed{
  D\propto \frac1p
  }.
\]
\end{FormulaBox}

\begin{WarnBox}[看起来反直觉但很重要]
降低压强会减少分子数量，但同时会增大平均自由程。对热传导和黏滞，这两个因素近似抵消，所以 \(\kappa,\eta\) 对压强不敏感。对扩散，没有 \(\rho\) 因子抵消，所以 \(D\) 随压强降低而增大。
\end{WarnBox}

\subsection{温度依赖的定性判断}

对理想气体，
\[
  \bar v\propto \sqrt{T}.
\]
若分子直径近似不变，在压强不太低的普通稀薄气体中：
\begin{itemize}
  \item \(\kappa\) 随温度升高通常增大；
  \item \(\eta\) 随温度升高通常增大；
  \item \(D\) 随温度升高通常增大，并且在定压下还会受到 \(n\propto 1/T\) 的影响。
\end{itemize}

热学B复习中，定性趋势比精确温度幂律更重要。

\section{极稀薄气体中的输运过程}

\subsection{普通输运公式的适用条件}

前面得到的输运系数公式默认分子在气体内部频繁碰撞，即平均自由程 \(\lambda\) 远小于容器特征长度 \(L\)：
\[
  \lambda\ll L.
\]
此时分子携带的信息大致来自距离截面一个平均自由程的区域。

\subsection{极稀薄气体}

当气体压强很低时，
\[
  \lambda
\]
会显著增大。如果
\[
  \lambda\sim L
  \quad \text{甚至} \quad
  \lambda>L,
\]
则分子主要在器壁之间碰撞，分子间碰撞反而不再占主导。此时普通公式
\[
  \kappa=\frac13\lambda\rho c_V\bar v,
  \qquad
  \eta=\frac13\lambda\rho\bar v,
  \qquad
  D=\frac13\lambda\bar v
\]
不再直接适用。

\begin{KeyBox}[极稀薄气体的物理图像]
极稀薄气体中，分子从一个器壁飞到另一个器壁，中间很少与其它分子碰撞。热量和动量更多通过分子与器壁的碰撞交换来输运，而不是通过气体内部的分子间碰撞逐层传递。
\end{KeyBox}

\subsection{杜瓦瓶与真空绝热}

杜瓦瓶夹层抽成高真空后，气体分子数密度 \(n\) 变小。在极稀薄气体范围内，单位时间能从热壁飞到冷壁并输运能量的分子数减少，因此气体热传导显著降低。这是杜瓦瓶真空层隔热的重要原因之一。

\begin{WarnBox}[普通压强无关结论有范围]
“气体热导率近似与压强无关”只适用于普通稀薄气体范围，不适用于极稀薄气体。当平均自由程与容器尺寸可比时，热导率会随压强降低而降低。
\end{WarnBox}

\section{第五章题型模板}

\subsection{题型一：判断输运类型}

\begin{MethodBox}[三问判断法]
\begin{ReviewList}
  \item 不均匀的是温度 \(T\) 吗？若是，热传导，输运热量。
  \item 不均匀的是宏观流速 \(u\) 吗？若是，黏滞，输运动量。
  \item 不均匀的是组分数密度 \(n\) 吗？若是，扩散，输运分子数。
\end{ReviewList}
\end{MethodBox}

\subsection{题型二：用宏观规律求输运量}

\begin{MethodBox}[通用计算流程]
\begin{ReviewList}
  \item 先确定梯度：
  \[
    \frac{\dd T}{\dd z},\quad
    \frac{\dd u}{\dd z},\quad
    \frac{\dd n}{\dd z}.
  \]
  \item 再选对应定律：
  \[
    \frac{\Delta Q}{\Delta t}
    =
    -\kappa\frac{\dd T}{\dd z}\Delta S,
  \]
  \[
    f
    =
    -\eta\frac{\dd u}{\dd z}\Delta S,
  \]
  \[
    \frac{\Delta N}{\Delta t}
    =
    -D\frac{\dd n}{\dd z}\Delta S.
  \]
  \item 最后根据题目要求判断是否乘以时间 \(\Delta t\)。
\end{ReviewList}
\end{MethodBox}

\subsection{题型三：由微观量求输运系数}

\begin{MethodBox}[输运系数计算]
若已知平均自由程 \(\lambda\)、平均速率 \(\bar v\)、质量密度 \(\rho\)、定容比热 \(c_V\)，则
\[
  \kappa=\frac13\lambda\rho c_V\bar v.
\]
若求黏滞系数：
\[
  \eta=\frac13\lambda\rho\bar v.
\]
若求扩散系数：
\[
  D=\frac13\lambda\bar v.
\]
注意三者维度不同，不能混用。
\end{MethodBox}

\subsection{题型四：判断压强变化对输运系数的影响}

\begin{MethodBox}[压强依赖判断]
温度一定、普通稀薄气体中：
\[
  \lambda\propto\frac1p,
  \qquad
  \rho\propto p.
\]
因此
\[
  \lambda\rho\approx \mathrm{const}.
\]
所以：
\[
  \kappa,\eta\ \text{近似不随}\ p\ \text{变化},
  \qquad
  D\propto\frac1p.
\]
如果题目进入极稀薄气体条件 \(\lambda\sim L\)，上述结论失效。
\end{MethodBox}

\section{第五章公式速查}

\formulatable{
\begin{tabularx}{\textwidth}{@{}>{\bfseries}lXl@{}}
\toprule
内容 & 公式 & 条件\\
\midrule

傅里叶热传导定律 &
\(\displaystyle \frac{\Delta Q}{\Delta t}=-\kappa\frac{\dd T}{\dd z}\Delta S\)
& 温度梯度\\[0.9em]

热流密度 &
\(\displaystyle j_Q=-\kappa\frac{\dd T}{\dd z}\)
& 单位面积单位时间\\[0.9em]

牛顿黏性定律 &
\(\displaystyle f=-\eta\frac{\dd u}{\dd z}\Delta S\)
& 流速梯度\\[0.9em]

切应力 &
\(\displaystyle \tau=-\eta\frac{\dd u}{\dd z}\)
& 单位面积黏滞力\\[0.9em]

菲克定律 &
\(\displaystyle \frac{\Delta N}{\Delta t}=-D\frac{\dd n}{\dd z}\Delta S\)
& 数密度梯度\\[0.9em]

扩散流密度 &
\(\displaystyle j_N=-D\frac{\dd n}{\dd z}\)
& 单位面积单位时间\\[0.9em]

平均自由程 &
\(\displaystyle \lambda=\frac{1}{\sqrt{2}\pi d^2n}\)
& 硬球模型\\[0.9em]

平均碰撞频率 &
\(\displaystyle Z=\frac{\bar v}{\lambda}=\sqrt{2}\pi d^2n\bar v\)
& 平衡气体\\[0.9em]

热传导系数 &
\(\displaystyle \kappa=\frac13\lambda\rho c_V\bar v\)
& 普通稀薄气体\\[0.9em]

黏滞系数 &
\(\displaystyle \eta=\frac13\lambda\rho\bar v\)
& 普通稀薄气体\\[0.9em]

扩散系数 &
\(\displaystyle D=\frac13\lambda\bar v\)
& 普通稀薄气体\\[0.9em]

压强依赖 &
\(\displaystyle \kappa,\eta \approx \text{与 }p\text{ 无关},\quad D\propto\frac1p\)
& 温度一定、普通稀薄气体\\

\bottomrule
\end{tabularx}
}

\section{第五章易错点总表}

\begin{WarnBox}[本章高频错误]
\begin{PitfallList}
  \item 把热传导、黏滞、扩散分开死记，忘记它们共同结构都是“流密度 = 负的输运系数乘以梯度”。
  \item 忘记公式中的负号。负号表示输运方向使不均匀性减弱。
  \item 把热流率和热量混淆。热流率是 \(\Delta Q/\Delta t\)，热量还要乘以时间。
  \item 把黏滞力方向判断错。黏滞力总是阻碍不同流层之间的相对运动。
  \item 把扩散和对流混为一谈。扩散不需要流体整体定向运动。
  \item 把 \(\kappa,\eta,D\) 三个输运系数混用。它们对应不同输运量，单位也不同。
  \item 忘记 \(D=\frac13\lambda\bar v\) 中没有 \(\rho\) 和 \(c_V\)。
  \item 认为气体热导率在所有压强下都与压强无关。该结论只适用于普通稀薄气体，不适用于极稀薄气体。
  \item 忽略普通输运公式的条件：平均自由程必须远小于容器特征尺寸。
  \item 把分子间碰撞看成阻碍输运的唯一因素。实际上分子热运动提供空间转移，碰撞决定输运的有效尺度，两者共同产生普通气体输运规律。
\end{PitfallList}
\end{WarnBox}

\section{第五章复习路线}

\begin{MethodBox}[建议背诵顺序]
\begin{ReviewList}
  \item 先背三种不均匀性：温度 \(T\)、流速 \(u\)、数密度 \(n\)。
  \item 再背三种宏观规律：
  \[
    j_Q=-\kappa\frac{\dd T}{\dd z},
    \qquad
    \tau=-\eta\frac{\dd u}{\dd z},
    \qquad
    j_N=-D\frac{\dd n}{\dd z}.
  \]
  \item 再理解微观图像：分子热运动负责跨层转移，碰撞决定平均自由程。
  \item 再背三个输运系数：
  \[
    \kappa=\frac13\lambda\rho c_V\bar v,
    \qquad
    \eta=\frac13\lambda\rho\bar v,
    \qquad
    D=\frac13\lambda\bar v.
  \]
  \item 最后记压强依赖：普通稀薄气体中 \(\kappa,\eta\) 近似与压强无关，\(D\propto1/p\)。
\end{ReviewList}
\end{MethodBox}

\begin{KeyBox}[第五章一句话总结]
第五章的核心是“非平衡导致输运，输运消除不均匀”。热传导输运能量，黏滞输运动量，扩散输运分子数；宏观规律看梯度，微观解释看平均自由程、平均速率和分子碰撞。
\end{KeyBox}
\chapter{固、液体性质简介与相变}

\chaptergoal{
本章是热学B的收束章节，内容从气体扩展到凝聚态物质与相变现象。复习重点是掌握几个清晰的物理图像：晶体有长程有序结构，液体有近程有序结构；液体表面张力来自表面层分子受力不平衡；弯曲液面导致附加压强；相变伴随潜热和物性突变；克拉珀龙方程给出相平衡曲线斜率。
}

\exammap{
\begin{ReviewList}
  \item \textbf{会区分：}晶体、非晶体、单晶体、多晶体的基本特征。
  \item \textbf{会解释：}晶体各向异性、确定熔点、空间点阵结构。
  \item \textbf{会使用：}固体经典热容量 \(C_{V,m}\approx3R\)。
  \item \textbf{会理解：}液体近程有序、表面层、表面张力和表面能。
  \item \textbf{会计算：}液体表面张力力、表面积增大所需功、弯曲液面附加压强、毛细上升或下降高度。
  \item \textbf{会判断：}润湿与不润湿时毛细管液面上升或下降。
  \item \textbf{会掌握：}相、相变、一级相变、相平衡、相图、三相点、临界点等概念。
  \item \textbf{会使用：}克拉珀龙方程
  \[
    \frac{\dd p}{\dd T}=\frac{l}{T(v_2-v_1)}.
  \]
\end{ReviewList}
}

\section{固体的基本分类与宏观性质}

\subsection{固体的主要特征}

固体的主要特征是能够保持一定体积和一定形状。与气体相比，固体分子或原子之间平均距离小、相互作用强；与液体相比，固体在宏观上能维持稳定形状。

固体通常分为两类：
\begin{ReviewList}
  \item \concept{晶体}：内部粒子有规则、周期性排列，例如岩盐、云母、水晶、冰、金属、石墨等。
  \item \concept{非晶体}：内部粒子没有长程周期性排列，例如玻璃、松脂、沥青、橡胶、塑料等。
\end{ReviewList}

\begin{KeyBox}[晶体与非晶体的核心差别]
晶体具有长程有序结构，因此常有规则外形、各向异性和确定熔点；非晶体缺少长程有序结构，因此通常没有确定熔点，更像过冷液体。
\end{KeyBox}

\subsection{晶体的宏观性质}

晶体常见宏观性质包括：

\begin{FormulaBox}[晶体的三个典型宏观性质]
\begin{tabularx}{\textwidth}{@{}>{\bfseries}lX@{}}
规则几何外形 & 晶体常由若干平面围成凸多面体。属于同一品种的晶体，对应晶面夹角恒定。\\
各向异性 & 晶体在不同方向上原子排列不同，因此力学、热学、电学、光学性质可能随方向变化。\\
确定熔点 & 晶体熔化时温度保持在确定熔点附近，熔化曲线上出现平台。\\
\end{tabularx}
\end{FormulaBox}

\begin{WarnBox}[易错点：各向异性不是所有固体都有]
单晶体通常表现出明显各向异性；多晶体由于许多小晶粒取向杂乱，宏观上常表现为各向同性。非晶体通常也表现为各向同性。
\end{WarnBox}

\subsection{单晶体与多晶体}

\concept{单晶体}是指整个晶体内粒子排列方向基本一致的晶体。它的各向异性通常比较明显。

\concept{多晶体}由许多微小单晶粒组成，每个晶粒内部有规则排列，但不同晶粒取向杂乱。由于各方向平均效果，多晶体宏观上往往近似各向同性。

\begin{KeyBox}[单晶体与多晶体的对比]
\begin{tabularx}{\textwidth}{@{}>{\bfseries}lCC@{}}
\toprule
 & 单晶体 & 多晶体\\
\midrule
内部结构 & 整体取向一致 & 由许多取向不同的小晶粒组成\\
各向异性 & 明显 & 宏观上常近似各向同性\\
典型例子 & 水晶、单晶硅 & 普通金属材料\\
\bottomrule
\end{tabularx}
\end{KeyBox}

\section{晶体的微观结构与热学性质}

\subsection{空间点阵}

晶体中粒子有规则地、周期性地排列，形成\concept{空间点阵}结构。授课 PDF 中以 NaCl、Diamond、CsCl、Graphite 等为例说明晶体粒子的周期性排列。:contentReference[oaicite:1]{index=1}

空间点阵反映的是晶体结构的平移周期性。晶体的宏观规则外形和各向异性，根源都在于这种微观周期性。

\subsection{晶体中的结合力}

晶体粒子之间存在相互作用力。常见晶体结合类型包括：
\begin{itemize}
  \item 离子晶体：正负离子之间的库仑作用。
  \item 共价晶体：原子之间通过共价键结合。
  \item 金属晶体：金属离子与自由电子之间的相互作用。
  \item 分子晶体：分子之间通过较弱的分子间作用力结合。
\end{itemize}

结合能越大，晶体通常越稳定，熔点也往往较高。

\subsection{晶体中的热运动}

晶体中的粒子并非静止不动，而是在平衡位置附近做热振动。温度越高，振动越剧烈。

在经典近似下，每个原子有三个振动方向。每个方向包含动能和势能两个平方项，按能量均分定理，每个原子平均振动能量为
\[
  3k_{\mathrm B}T.
\]
因此 \(1\) mol 固体的摩尔内能近似为
\[
  u=3N_Ak_{\mathrm B}T=3RT.
\]
摩尔热容量为
\[
  \boxed{
  C_{V,m}=\frac{\dd u}{\dd T}=3R
  }
\]
这就是\concept{杜隆--珀替定律}。

\begin{FormulaBox}[固体经典热容量]
\[
  \boxed{
  C_{V,m}\approx3R\approx25\,\mathrm{J\cdot mol^{-1}\cdot K^{-1}}
  }.
\]
\end{FormulaBox}

\begin{WarnBox}[经典热容量的局限]
杜隆--珀替定律是高温经典近似。低温下，固体热容量明显低于 \(3R\)，这是量子效应的结果。热学B中重点掌握公式和适用范围，不需要深入量子理论。
\end{WarnBox}

\section{液体结构与基本性质}

\subsection{液体的宏观性质}

液体具有一定体积，但没有固定形状，能适应容器形状。液体分子之间距离较小，相互作用较强，因此液体比气体难压缩得多。

液体同时具有某种“介于气体与固体之间”的性质：
\begin{itemize}
  \item 与固体相似：分子间距小，具有较强相互作用，体积较稳定。
  \item 与气体相似：分子可以相对流动，液体没有固定形状。
\end{itemize}

\subsection{液体的微观结构}

液体分子排列不具有晶体那样的长程周期性，但在短距离范围内仍有一定有序性，称为\concept{近程有序}。

\begin{KeyBox}[液体结构图像]
液体不像气体那样分子彼此相距很远，也不像晶体那样有长程周期点阵。液体的结构特点是：短程有序，长程无序。
\end{KeyBox}

液体分子不断热运动，并可在相邻分子形成的临时“空隙”中移动，因此液体具有流动性。

\section{液体表面张力}

\subsection{表面层与表面张力}

液体表面附近有一个薄层，称为\concept{表面层}。液体内部的分子受到周围分子吸引，合力平均为零；而表面层分子上方液体分子较少，受力不对称，合力指向液体内部。

因此，液体表面有自动收缩的趋势，表现为\concept{表面张力}。

\begin{FormulaBox}[表面张力定义]
若液面边界线长度为 \(l\)，液体表面张力为
\[
  \boxed{
  f=\sigma l
  }.
\]
其中 \(\sigma\) 称为表面张力系数，单位为
\[
  \mathrm{N/m}.
\]
\end{FormulaBox}

表面张力系数也可从能量角度理解。若液体表面积增加 \(\Delta S\)，外力克服表面张力做功为
\[
  \boxed{
  \Delta A=\sigma\Delta S
  }.
\]
因此 \(\sigma\) 也等于单位面积表面能。

\begin{KeyBox}[表面张力的两个等价理解]
\begin{itemize}
  \item 力学理解：\(\sigma\) 是单位长度边界线两侧液面的相互拉力。
  \item 能量理解：\(\sigma\) 是增加单位液体表面积所需的功。
\end{itemize}
\end{KeyBox}

\subsection{液体表面积趋于最小}

由于增大表面积需要做功，液体在无外力或外力影响较小时，会趋向于使表面积最小。例如小液滴常近似呈球形，因为在给定体积下球形表面积最小。

\begin{WarnBox}[表面张力方向]
表面张力沿液面切向，并垂直于所取边界线。不要把它误认为垂直液面的力。垂直液面的压强差来自弯曲液面，而不是表面张力本身直接垂直作用。
\end{WarnBox}

\section{弯曲液面的附加压强}

\subsection{球形液面的附加压强}

弯曲液面两侧存在压强差，称为\concept{附加压强}。对半径为 \(R\) 的球形液面，
\[
  \boxed{
  \Delta p=\frac{2\sigma}{R}
  }.
\]
其中 \(\Delta p\) 表示凹侧压强相对于凸侧压强的增加。

\begin{KeyBox}[附加压强方向]
弯曲液面附加压强总是指向曲率中心一侧。也就是说，弯曲液面凹侧压强大于凸侧压强。
\end{KeyBox}

例如肥皂泡有内外两个表面，因此肥皂泡内外压强差为
\[
  \boxed{
  \Delta p=\frac{4\sigma}{R}
  }.
\]
普通液滴只有一个液面，内外压强差为
\[
  \Delta p=\frac{2\sigma}{R}.
\]

\subsection{柱形液面的附加压强}

对半径为 \(R\) 的柱形液面，附加压强为
\[
  \boxed{
  \Delta p=\frac{\sigma}{R}
  }.
\]

更一般地，对两个主曲率半径分别为 \(R_1,R_2\) 的弯曲液面，有拉普拉斯公式
\[
  \boxed{
  \Delta p=\sigma\left(\frac{1}{R_1}+\frac{1}{R_2}\right)
  }.
\]
球面 \(R_1=R_2=R\)，得到 \(2\sigma/R\)；柱面 \(R_1=R,R_2=\infty\)，得到 \(\sigma/R\)。

\begin{WarnBox}[易错点：液滴和肥皂泡差一倍]
液滴只有一个液气界面，\(\Delta p=2\sigma/R\)。肥皂泡薄膜有内外两个界面，\(\Delta p=4\sigma/R\)。这类题非常容易漏掉一个表面。
\end{WarnBox}

\section{毛细现象}

\subsection{润湿与不润湿}

液体与固体接触时，会形成一定接触角 \(\theta\)。若
\[
  \theta<90^\circ,
\]
称为\concept{润湿}，液面在毛细管中上升。若
\[
  \theta>90^\circ,
\]
称为\concept{不润湿}，液面在毛细管中下降。

\begin{KeyBox}[润湿判断]
\begin{itemize}
  \item 水润湿干净玻璃，毛细管中水面上升，液面为凹面。
  \item 水银不润湿玻璃，毛细管中水银面下降，液面为凸面。
\end{itemize}
\end{KeyBox}

\subsection{毛细上升或下降高度}

设毛细管半径为 \(r\)，液体密度为 \(\rho\)，表面张力系数为 \(\sigma\)，接触角为 \(\theta\)。液柱高度满足
\[
  \boxed{
  h=\frac{2\sigma\cos\theta}{\rho gr}
  }.
\]

若润湿管壁，
\[
  \theta<90^\circ,\qquad \cos\theta>0,
\]
则
\[
  h>0,
\]
液面上升。

若不润湿管壁，
\[
  \theta>90^\circ,\qquad \cos\theta<0,
\]
则
\[
  h<0,
\]
液面下降。

\begin{WarnBox}[毛细公式的符号]
毛细公式本身
\[
  h=\frac{2\sigma\cos\theta}{\rho gr}
\]
已经包含正负号。若题目只问下降高度的大小，可取绝对值；若问相对于外部液面的高度变化，要保留符号。
\end{WarnBox}

\section{相、相变与相平衡}

\subsection{相的概念}

\concept{相}是系统中物理性质和化学性质均匀、并与其它部分有明显界面的部分。常见相包括固相、液相和气相。

一个系统可以是单相，也可以是多相。例如密闭容器中的水和水蒸气共存时，液态水和水蒸气分别构成不同相。

\subsection{相变}

\concept{相变}是物质从一种相转变为另一种相的过程。常见相变包括：
\begin{itemize}
  \item 熔化与凝固：固相与液相之间转变。
  \item 汽化与液化：液相与气相之间转变。
  \item 升华与凝华：固相与气相之间转变。
\end{itemize}

\begin{FormulaBox}[常见相变]
\begin{tabularx}{\textwidth}{@{}>{\bfseries}lCC@{}}
\toprule
相变类型 & 正向过程 & 反向过程\\
\midrule
固液相变 & 熔化 & 凝固\\
液气相变 & 汽化 & 液化\\
固气相变 & 升华 & 凝华\\
\bottomrule
\end{tabularx}
\end{FormulaBox}

\subsection{一级相变的普遍特征}

\concept{一级相变}通常具有以下特征：
\begin{ReviewList}
  \item 相变过程中温度保持不变。
  \item 相变过程中吸收或放出潜热。
  \item 熵、体积等一阶热力学量发生突变。
  \item 两相可以在相变温度和相变压强下共存。
\end{ReviewList}

授课 PDF 将单元系一级相变、气液相变、克拉珀龙方程、固液和固气相变作为相变部分主线。:contentReference[oaicite:2]{index=2}

\begin{WarnBox}[易错点：相变时温度不变不代表不吸热]
例如水在 \(100^\circ\mathrm C\) 沸腾时，继续吸热但温度不升高，吸收的热量用于克服分子间作用并完成相变，表现为汽化潜热。
\end{WarnBox}

\subsection{相平衡与相图}

若两个相在一定温度和压强下共存，并且各相质量不随时间变化，则称为\concept{相平衡}。

相平衡状态可以在 \(p\)--\(T\) 图中表示。常见相图中有三条相平衡曲线：
\begin{itemize}
  \item 升华曲线：固相与气相平衡。
  \item 熔解曲线：固相与液相平衡。
  \item 汽化曲线：液相与气相平衡。
\end{itemize}

三条曲线交于\concept{三相点}，此时固、液、气三相共存。

\section{气液相变}

\subsection{汽化与液化}

\concept{汽化}是液体转变为气体的过程，包括蒸发和沸腾。蒸发可以在液体表面、任何温度下发生；沸腾是在液体内部和表面同时发生的剧烈汽化过程，通常在饱和蒸气压等于外界压强时发生。

\concept{液化}是气体转变为液体的过程。

\subsection{饱和蒸气压}

在密闭容器中，液体不断蒸发，同时蒸气也不断凝结。当蒸发和凝结达到动态平衡时，蒸气称为\concept{饱和蒸气}，其压强称为\concept{饱和蒸气压}。

饱和蒸气压只与温度有关：
\[
  p=p_{\mathrm s}(T).
\]
温度升高，饱和蒸气压增大。

\begin{KeyBox}[沸腾条件]
当液体的饱和蒸气压等于外界压强时，液体发生沸腾：
\[
  p_{\mathrm s}(T)=p_{\mathrm{ext}}.
\]
外界压强越低，沸点越低；外界压强越高，沸点越高。
\end{KeyBox}

\subsection{临界点}

气液相变曲线终止于\concept{临界点}。临界点对应临界温度 \(T_c\)、临界压强 \(p_c\)、临界体积 \(V_c\)。

当温度高于临界温度时，无论怎样压缩气体，都不能通过等温压缩使其液化。此时气相和液相的界限消失，物质处于超临界状态。

\begin{WarnBox}[临界温度的意义]
临界温度不是普通沸点。它表示能否通过单纯加压液化气体的最高温度。若 \(T>T_c\)，单纯加压不能液化。
\end{WarnBox}

\subsection{气液二相图}

在 \(p\)--\(V\) 图上，低于临界温度的等温线会出现气液共存区域。气液两相共存时，压强为该温度下的饱和蒸气压，随体积变化压强保持不变，直到某一相完全消失。

随着温度升高，气液共存区域缩小；到临界温度时，气相和液相差别消失。

\begin{KeyBox}[气液二相图的读法]
\begin{itemize}
  \item 饱和液体线与饱和蒸气线之间是气液共存区。
  \item 共存区内改变体积，主要改变两相比例，而不是改变压强。
  \item 临界点以上不再有清楚的气液分界。
\end{itemize}
\end{KeyBox}

\section{克拉珀龙方程}

\subsection{方程形式}

对单元系一级相变，两相平衡曲线斜率由\concept{克拉珀龙方程}给出：
\[
  \boxed{
  \frac{\dd p}{\dd T}
  =
  \frac{l}{T(v_2-v_1)}
  }.
\]

其中：
\begin{itemize}
  \item \(p\) 为相平衡压强，例如饱和蒸气压。
  \item \(T\) 为相变温度。
  \item \(l\) 为单位质量或单位摩尔相变潜热。
  \item \(v_1,v_2\) 分别为两相的比体积或摩尔体积。
\end{itemize}

若 \(l\) 用单位质量潜热，则 \(v_1,v_2\) 应使用单位质量体积；若 \(l\) 用摩尔潜热，则 \(v_1,v_2\) 应使用摩尔体积。

\begin{WarnBox}[单位必须配套]
克拉珀龙方程中 \(l\) 与 \(v\) 的单位必须一致：
\[
  l:\mathrm{J/kg}
  \quad\Rightarrow\quad
  v:\mathrm{m^3/kg};
\]
\[
  l:\mathrm{J/mol}
  \quad\Rightarrow\quad
  v:\mathrm{m^3/mol}.
\]
否则量纲会错。
\end{WarnBox}

\subsection{克拉珀龙方程的物理意义}

克拉珀龙方程说明：相平衡曲线斜率由相变潜热和相变体积变化共同决定。

\begin{KeyBox}[斜率判断]
\[
  \frac{\dd p}{\dd T}
  =
  \frac{l}{T(v_2-v_1)}.
\]
通常 \(l>0\)，\(T>0\)，因此斜率正负主要由 \(v_2-v_1\) 决定。
\end{KeyBox}

\section{克劳修斯--克拉珀龙近似}

\subsection{气液相变的近似}

对液体汽化过程，令相 1 为液相，相 2 为气相。通常
\[
  v_{\mathrm g}\gg v_{\mathrm l},
\]
因此
\[
  v_2-v_1\approx v_{\mathrm g}.
\]
若把蒸气近似看作理想气体，
\[
  v_{\mathrm g}=\frac{RT}{p}
\]
这里 \(v_{\mathrm g}\) 为摩尔体积，\(l\) 为摩尔汽化热。

代入克拉珀龙方程得
\[
  \frac{\dd p}{\dd T}
  =
  \frac{l}{T(RT/p)}
  =
  \boxed{
  \frac{lp}{RT^2}
  }.
\]

若 \(l\) 在温度范围内近似为常数，则
\[
  \frac{\dd p}{p}
  =
  \frac{l}{R}\frac{\dd T}{T^2}.
\]
积分得
\[
  \boxed{
  \ln p=-\frac{l}{RT}+C
  }.
\]

\begin{KeyBox}[饱和蒸气压随温度变化]
由
\[
  \ln p=-\frac{l}{RT}+C
\]
可知温度升高时，饱和蒸气压迅速增大。
\end{KeyBox}

\subsection{两温度下饱和蒸气压关系}

若在 \(T_1,T_2\) 下饱和蒸气压分别为 \(p_1,p_2\)，且汽化潜热 \(l\) 近似不变，则
\[
  \boxed{
  \ln\frac{p_2}{p_1}
  =
  -\frac{l}{R}
  \left(
  \frac{1}{T_2}-\frac{1}{T_1}
  \right)
  }.
\]
也可写作
\[
  \boxed{
  \ln\frac{p_2}{p_1}
  =
  \frac{l}{R}
  \left(
  \frac{1}{T_1}-\frac{1}{T_2}
  \right)
  }.
\]

\begin{WarnBox}[指数符号检查]
若 \(T_2>T_1\)，则 \(p_2>p_1\)，所以
\[
  \ln\frac{p_2}{p_1}>0.
\]
写完公式后可以用这个物理判断检查符号。
\end{WarnBox}

\section{固液相变与固气相变}

\subsection{固液相变}

固液相变包括熔化和凝固。熔化时吸热，凝固时放热。对多数物质，熔化后体积增大：
\[
  v_{\mathrm l}>v_{\mathrm s}.
\]
因此
\[
  \frac{\dd p}{\dd T}>0.
\]
这意味着压强增大时，熔点升高。

但水是重要例外。冰熔化成水时体积减小：
\[
  v_{\mathrm l}<v_{\mathrm s}.
\]
因此
\[
  \frac{\dd p}{\dd T}<0.
\]
压强增大时，冰的熔点降低。

\begin{WarnBox}[冰的熔解曲线斜率]
大多数物质熔解曲线斜率为正；冰的熔解曲线斜率为负。根源是冰的比体积大于水的比体积。
\end{WarnBox}

\subsection{固气相变}

固气相变包括升华和凝华。升华吸热，凝华放热。固气相变也伴随潜热吸放和体积突变。

在 \(p\)--\(T\) 相图上，升华曲线是固相和气相的分界线。授课 PDF 中指出，升华曲线、熔解曲线、汽化曲线的斜率均由克拉珀龙方程决定。:contentReference[oaicite:3]{index=3}

升华热与熔解热、汽化热之间通常满足近似关系：
\[
  \boxed{
  l_{\mathrm{sub}}
  \approx
  l_{\mathrm{fus}}+l_{\mathrm{vap}}
  }.
\]
即升华可看作先熔化再汽化的组合过程。

\section{第六章题型模板}

\subsection{题型一：表面张力做功}

\begin{MethodBox}[表面积变化题]
若题目给出液膜或液面面积变化 \(\Delta S\)，表面张力做功常用：
\[
  \Delta A=\sigma\Delta S.
\]
若是液膜，例如肥皂膜，通常有两个表面，面积变化要乘以 \(2\)。关键是先数清楚有几个液气界面。
\end{MethodBox}

\subsection{题型二：弯曲液面压强差}

\begin{MethodBox}[液滴、气泡、肥皂泡]
\begin{ReviewList}
  \item 普通液滴：
  \[
    \Delta p=\frac{2\sigma}{R}.
  \]
  \item 液体内气泡，只有一个界面：
  \[
    \Delta p=\frac{2\sigma}{R}.
  \]
  \item 肥皂泡，有内外两个界面：
  \[
    \Delta p=\frac{4\sigma}{R}.
  \]
  \item 柱形液面：
  \[
    \Delta p=\frac{\sigma}{R}.
  \]
\end{ReviewList}
\end{MethodBox}

\subsection{题型三：毛细现象}

\begin{MethodBox}[毛细高度计算]
直接使用
\[
  h=\frac{2\sigma\cos\theta}{\rho gr}.
\]
\begin{itemize}
  \item \(\theta<90^\circ\)，\(\cos\theta>0\)，液面上升。
  \item \(\theta>90^\circ\)，\(\cos\theta<0\)，液面下降。
  \item 管半径 \(r\) 越小，毛细高度绝对值越大。
\end{itemize}
\end{MethodBox}

\subsection{题型四：克拉珀龙方程}

\begin{MethodBox}[相平衡曲线斜率]
\begin{ReviewList}
  \item 先确定相 1 和相 2。
  \item 写出相变潜热 \(l\)。
  \item 判断体积变化 \(v_2-v_1\)。
  \item 使用
  \[
    \frac{\dd p}{\dd T}=\frac{l}{T(v_2-v_1)}.
  \]
  \item 若是汽化或升华，并且气相可近似为理想气体，可进一步使用克劳修斯--克拉珀龙近似。
\end{ReviewList}
\end{MethodBox}

\section{第六章公式速查}

\formulatable{
\begin{tabularx}{\textwidth}{@{}>{\bfseries}lXl@{}}
\toprule
内容 & 公式 & 条件\\
\midrule

固体经典摩尔热容量 &
\(\displaystyle C_{V,m}\approx3R\)
& 高温经典近似\\[0.8em]

表面张力 &
\(\displaystyle f=\sigma l\)
& 边界线长度 \(l\)\\[0.8em]

表面能 &
\(\displaystyle \Delta A=\sigma\Delta S\)
& 增加表面积 \(\Delta S\)\\[0.8em]

球形液面附加压强 &
\(\displaystyle \Delta p=\frac{2\sigma}{R}\)
& 单个球形界面\\[0.8em]

肥皂泡压强差 &
\(\displaystyle \Delta p=\frac{4\sigma}{R}\)
& 两个界面\\[0.8em]

柱形液面附加压强 &
\(\displaystyle \Delta p=\frac{\sigma}{R}\)
& 柱形界面\\[0.8em]

一般拉普拉斯公式 &
\(\displaystyle \Delta p=\sigma\left(\frac1{R_1}+\frac1{R_2}\right)\)
& 两主曲率半径\\[0.8em]

毛细高度 &
\(\displaystyle h=\frac{2\sigma\cos\theta}{\rho gr}\)
& 半径 \(r\) 的毛细管\\[0.8em]

一级相变熵变 &
\(\displaystyle \Delta S=\frac{l}{T}\)
& 可逆等温相变\\[0.8em]

克拉珀龙方程 &
\(\displaystyle \frac{\dd p}{\dd T}=\frac{l}{T(v_2-v_1)}\)
& 单元系一级相变\\[1em]

克劳修斯--克拉珀龙近似 &
\(\displaystyle \frac{\dd p}{\dd T}=\frac{lp}{RT^2}\)
& 气相近似理想气体\\[1em]

饱和蒸气压积分式 &
\(\displaystyle \ln p=-\frac{l}{RT}+C\)
& \(l\) 近似常数\\[1em]

两温度饱和蒸气压关系 &
\(\displaystyle \ln\frac{p_2}{p_1}=\frac{l}{R}\left(\frac1{T_1}-\frac1{T_2}\right)\)
& \(l\) 近似常数\\

\bottomrule
\end{tabularx}
}

\section{第六章易错点总表}

\begin{WarnBox}[本章高频错误]
\begin{PitfallList}
  \item 把晶体和非晶体都当作有确定熔点。晶体有确定熔点，非晶体通常没有确定熔点。
  \item 认为所有固体都各向异性。单晶体明显各向异性，多晶体宏观上常近似各向同性。
  \item 忘记杜隆--珀替定律 \(C_{V,m}\approx3R\) 是高温经典近似。
  \item 把表面张力方向画成垂直液面。表面张力沿液面切向。
  \item 液滴和肥皂泡压强差混淆。肥皂泡有两个表面，压强差是 \(4\sigma/R\)。
  \item 毛细公式中忘记 \(\cos\theta\)，导致无法判断上升还是下降。
  \item 把沸腾条件误写成温度达到 \(100^\circ\mathrm C\)。正确条件是饱和蒸气压等于外界压强。
  \item 认为相变时温度不变所以没有热量交换。一级相变有潜热。
  \item 克拉珀龙方程中 \(l\) 和 \(v\) 的单位不匹配。
  \item 判断熔解曲线斜率时忘记看 \(v_{\mathrm l}-v_{\mathrm s}\) 的符号。
  \item 认为所有气体都能靠加压液化。高于临界温度时，单纯加压不能液化。
\end{PitfallList}
\end{WarnBox}

\section{第六章复习路线}

\begin{MethodBox}[建议背诵顺序]
\begin{ReviewList}
  \item 先背固体分类：晶体、非晶体；单晶体、多晶体。
  \item 再背晶体性质：规则外形、各向异性、确定熔点、空间点阵。
  \item 再背固体热容量：
  \[
    C_{V,m}\approx3R.
  \]
  \item 再理解液体：近程有序、长程无序、具有流动性。
  \item 再背表面张力：
  \[
    f=\sigma l,\qquad \Delta A=\sigma\Delta S.
  \]
  \item 再背弯曲液面与毛细现象：
  \[
    \Delta p=\frac{2\sigma}{R},\qquad
    h=\frac{2\sigma\cos\theta}{\rho gr}.
  \]
  \item 再背相变概念：相、相平衡、一级相变、潜热、三相点、临界点。
  \item 最后背克拉珀龙方程：
  \[
    \frac{\dd p}{\dd T}=\frac{l}{T(v_2-v_1)}.
  \]
\end{ReviewList}
\end{MethodBox}

\begin{KeyBox}[第六章一句话总结]
第六章的核心是把凝聚态和相变现象用热学语言描述：晶体的性质来自周期结构，液体表面现象来自表面张力，相变曲线由潜热和体积突变决定，而克拉珀龙方程把这些量统一到相平衡曲线的斜率中。
\end{KeyBox}
\chapter{总复习框架与考前速查}

\chaptergoal{
本章不是第七章，而是对前六章的压缩整理。目标是在考前快速建立全局框架：热学B主要由三条线组成。第一条是宏观热力学线：状态、过程、第一定律、第二定律、熵；第二条是微观统计线：分子动理论、麦克斯韦分布、玻尔兹曼分布、能量均分；第三条是物质性质线：输运、固液体性质、表面现象与相变。
}

\section{全课程知识地图}

\begin{KeyBox}[热学B的三条主线]
\begin{ReviewList}
  \item \textbf{宏观热力学：}
  \[
    \text{系统与状态}
    \rightarrow
    \text{第一定律}
    \rightarrow
    \text{第二定律}
    \rightarrow
    \text{熵与过程方向}
  \]
  这一部分回答：状态如何描述？过程如何计算？能量如何守恒？过程能否自发发生？

  \item \textbf{微观统计物理：}
  \[
    \text{分子热运动}
    \rightarrow
    \text{速度分布}
    \rightarrow
    \text{玻尔兹曼分布}
    \rightarrow
    \text{能量均分与热容量}
  \]
  这一部分回答：压强和温度的微观来源是什么？大量分子速度如何分布？热容量为什么与自由度有关？

  \item \textbf{物质性质与相变：}
  \[
    \text{输运过程}
    \rightarrow
    \text{固液体性质}
    \rightarrow
    \text{表面张力}
    \rightarrow
    \text{相变与相图}
  \]
  这一部分回答：非平衡态如何趋于平衡？固体和液体有什么典型性质？相变曲线如何由热力学量决定？
\end{ReviewList}
\end{KeyBox}

\section{六章核心内容一句话}

\begin{FormulaBox}[章节主旨]
\begin{tabularx}{\textwidth}{@{}>{\bfseries}lX@{}}
第一章 & 建立热学语言：系统、状态、平衡态、准静态过程、温度、温标与状态方程。\\
第二章 & 用能量守恒处理过程计算：\(\dd U=\dbar Q+\dbar W\)，重点是 \(Q,W,\Delta U,\Delta H\)。\\
第三章 & 用第二定律和熵判断过程方向：自然过程满足孤立系统总熵不减。\\
第四章 & 从微观统计解释宏观量：压强来自分子碰壁，温度来自平均平动动能。\\
第五章 & 描述非平衡态下的输运：热量、动量、分子数都沿负梯度方向输运。\\
第六章 & 介绍凝聚态与相变：晶体、液体、表面张力、毛细现象、相图与克拉珀龙方程。\\
\end{tabularx}
\end{FormulaBox}

\section{最重要的概念对比}

\subsection{态函数与过程量}

\begin{FormulaBox}[态函数与过程量]
\begin{tabularx}{\textwidth}{@{}>{\bfseries}lXX@{}}
 & 态函数 & 过程量\\
\midrule
典型量 & \(p,V,T,U,H,S\) & \(Q,W\)\\
特点 & 只由状态决定 & 与具体过程路径有关\\
数学表现 & \(\oint\dd X=0\) & 一般 \(\oint\dbar Q\ne0,\ \oint\dbar W\ne0\)\\
典型错误 & 把熵当作过程量 & 把热量、功当作状态中“含有”的量\\
\end{tabularx}
\end{FormulaBox}

\begin{WarnBox}[必须记住]
\[
  U,H,S
\]
是态函数；
\[
  Q,W
\]
是过程量。热学B的大量错误都来自这一区分不清。
\end{WarnBox}

\subsection{准静态、可逆、绝热、等熵}

\begin{FormulaBox}[四个概念不要混]
\begin{tabularx}{\textwidth}{@{}>{\bfseries}lX@{}}
准静态过程 & 系统经过一系列无限接近平衡态的状态；可以画热力学过程曲线。\\
可逆过程 & 系统和外界都可以完全恢复原状；必须准静态且无耗散。\\
绝热过程 & 与外界无热交换，\(Q=0\)。\\
等熵过程 & \(\Delta S=0\)。可逆绝热过程才是等熵过程。\\
\end{tabularx}
\end{FormulaBox}

\begin{WarnBox}[高频陷阱]
\begin{itemize}
  \item 准静态不一定可逆。
  \item 绝热不一定等熵。
  \item 可逆绝热才等熵。
  \item 不可逆绝热可以有 \(\Delta S>0\)，例如理想气体自由膨胀。
\end{itemize}
\end{WarnBox}

\section{热力学过程计算总路线}

\subsection{通用过程计算}

\begin{MethodBox}[算 \(Q,W,\Delta U,\Delta H,\Delta S\) 的统一流程]
\begin{ReviewList}
  \item \textbf{第一步：判断过程类型。}
  看过程是等容、等压、等温、绝热、多方、节流、循环，还是不可逆过程。

  \item \textbf{第二步：写状态方程。}
  对理想气体：
  \[
    pV=\nu RT.
  \]

  \item \textbf{第三步：算功。}
  若是准静态体积功：
  \[
    W=-\int p\dd V.
  \]

  \item \textbf{第四步：算内能变化。}
  对理想气体：
  \[
    \Delta U=\nu C_{V,m}(T_2-T_1).
  \]

  \item \textbf{第五步：用第一定律补热量。}
  \[
    \Delta U=Q+W,
    \qquad
    Q=\Delta U-W.
  \]
  这里 \(W\) 是外界对系统做功。

  \item \textbf{第六步：若涉及方向性，算总熵变。}
  \[
    \Delta S_{\mathrm{total}}
    =
    \Delta S_{\mathrm{sys}}+\Delta S_{\mathrm{env}}.
  \]
\end{ReviewList}
\end{MethodBox}

\subsection{理想气体五类过程总表}

\begin{FormulaBox}[理想气体过程速查]
\begin{tabularx}{\textwidth}{@{}>{\bfseries}lXXX@{}}
\toprule
过程 & \(W\)：外界对系统做功 & \(\Delta U\) & \(Q\)\\
\midrule

等容 &
\(0\) &
\(\nu C_{V,m}\Delta T\) &
\(\nu C_{V,m}\Delta T\)\\[0.5em]

等压 &
\(-\nu R\Delta T\) &
\(\nu C_{V,m}\Delta T\) &
\(\nu C_{p,m}\Delta T\)\\[0.5em]

等温 &
\(-\nu RT\ln(V_2/V_1)\) &
\(0\) &
\(\nu RT\ln(V_2/V_1)\)\\[0.5em]

绝热可逆 &
\(\nu C_{V,m}\Delta T\) &
\(\nu C_{V,m}\Delta T\) &
\(0\)\\[0.5em]

多方 \(n\ne1\) &
\(\dfrac{\nu R\Delta T}{n-1}\) &
\(\nu C_{V,m}\Delta T\) &
\(\nu C_{n,m}\Delta T\)\\

\bottomrule
\end{tabularx}
\end{FormulaBox}

\begin{WarnBox}[符号检查]
若气体膨胀，通常 \(V_2>V_1\)，系统对外做正功，因此外界对系统做功 \(W<0\)。做题时先判断膨胀还是压缩，可以避免符号错。
\end{WarnBox}

\section{最核心公式总表}

\subsection{第一部分：状态方程与基本热力学量}

\formulatable{
\begin{tabularx}{\textwidth}{@{}>{\bfseries}lX@{}}
理想气体状态方程 &
\(\displaystyle pV=\nu RT=Nk_{\mathrm B}T\)\\[0.7em]

范德瓦耳斯方程 &
\(\displaystyle \left(p+\frac{\nu^2a}{V^2}\right)(V-\nu b)=\nu RT\)\\[0.9em]

固液体微分状态方程 &
\(\displaystyle \dd V=\alpha V\dd T-\beta V\dd p\)\\[0.7em]

第一定律 &
\(\displaystyle \dd U=\dbar Q+\dbar W\)\\[0.7em]

准静态体积功 &
\(\displaystyle \dbar W=-p\dd V\)\\[0.7em]

焓 &
\(\displaystyle H=U+pV\)\\[0.7em]

理想气体内能 &
\(\displaystyle \dd U=\nu C_{V,m}\dd T\)\\[0.7em]

理想气体焓 &
\(\displaystyle \dd H=\nu C_{p,m}\dd T\)\\[0.7em]

迈耶公式 &
\(\displaystyle C_{p,m}-C_{V,m}=R\)\\
\end{tabularx}
}

\subsection{第二部分：第二定律与熵}

\formulatable{
\begin{tabularx}{\textwidth}{@{}>{\bfseries}lX@{}}
卡诺热机效率 &
\(\displaystyle \eta_{\mathrm C}=1-\frac{T_2}{T_1}\)\\[0.8em]

卡诺制冷系数 &
\(\displaystyle \varepsilon_{\mathrm C}=\frac{T_2}{T_1-T_2}\)\\[0.8em]

克劳修斯等式 &
\(\displaystyle \oint_{\mathrm{rev}}\frac{\dbar Q}{T}=0\)\\[0.8em]

克劳修斯不等式 &
\(\displaystyle \oint\frac{\dbar Q}{T}\le0\)\\[0.8em]

熵定义 &
\(\displaystyle \dd S=\frac{\dbar Q_{\mathrm{rev}}}{T}\)\\[0.8em]

孤立系统熵增加 &
\(\displaystyle \Delta S_{\mathrm{isolated}}\ge0\)\\[0.8em]

理想气体熵变 &
\(\displaystyle \Delta S=\nu C_{V,m}\ln\frac{T_2}{T_1}+\nu R\ln\frac{V_2}{V_1}\)\\[1em]

理想气体熵变另一式 &
\(\displaystyle \Delta S=\nu C_{p,m}\ln\frac{T_2}{T_1}-\nu R\ln\frac{p_2}{p_1}\)\\
\end{tabularx}
}

\subsection{第三部分：气体动理论与统计分布}

\formulatable{
\begin{tabularx}{\textwidth}{@{}>{\bfseries}lX@{}}
压强公式 &
\(\displaystyle p=\frac13nm\overline{v^2}\)\\[0.8em]

温度统计解释 &
\(\displaystyle \overline{\varepsilon_k}=\frac32k_{\mathrm B}T\)\\[0.8em]

麦克斯韦速率分布 &
\(\displaystyle f(v)=4\pi\left(\frac{m}{2\pi k_{\mathrm B}T}\right)^{3/2}v^2
e^{-mv^2/(2k_{\mathrm B}T)}\)\\[1em]

最概然速率 &
\(\displaystyle v_{\mathrm p}=\sqrt{\frac{2k_{\mathrm B}T}{m}}=\sqrt{\frac{2RT}{\mu}}\)\\[0.8em]

平均速率 &
\(\displaystyle \bar v=\sqrt{\frac{8k_{\mathrm B}T}{\pi m}}=\sqrt{\frac{8RT}{\pi\mu}}\)\\[0.8em]

方均根速率 &
\(\displaystyle v_{\mathrm{rms}}=\sqrt{\frac{3k_{\mathrm B}T}{m}}=\sqrt{\frac{3RT}{\mu}}\)\\[0.8em]

玻尔兹曼分布 &
\(\displaystyle n(\bm r)=n_0\exp\left[-\frac{\varepsilon_p(\bm r)}{k_{\mathrm B}T}\right]\)\\[1em]

能量均分 &
\(\displaystyle \overline{\varepsilon_i}=\frac12k_{\mathrm B}T\)\\[0.8em]

理想气体热容量 &
\(\displaystyle C_{V,m}=\frac{i}{2}R,\qquad C_{p,m}=\frac{i+2}{2}R\)\\
\end{tabularx}
}

\subsection{第四部分：输运、表面现象与相变}

\formulatable{
\begin{tabularx}{\textwidth}{@{}>{\bfseries}lX@{}}
傅里叶定律 &
\(\displaystyle j_Q=-\kappa\frac{\dd T}{\dd z}\)\\[0.8em]

牛顿黏性定律 &
\(\displaystyle \tau=-\eta\frac{\dd u}{\dd z}\)\\[0.8em]

菲克定律 &
\(\displaystyle j_N=-D\frac{\dd n}{\dd z}\)\\[0.8em]

热传导系数 &
\(\displaystyle \kappa=\frac13\lambda\rho c_V\bar v\)\\[0.8em]

黏滞系数 &
\(\displaystyle \eta=\frac13\lambda\rho\bar v\)\\[0.8em]

扩散系数 &
\(\displaystyle D=\frac13\lambda\bar v\)\\[0.8em]

平均自由程 &
\(\displaystyle \lambda=\frac{1}{\sqrt2\pi d^2n}\)\\[0.8em]

表面张力 &
\(\displaystyle f=\sigma l,\qquad \Delta A=\sigma\Delta S\)\\[0.8em]

球形液面附加压强 &
\(\displaystyle \Delta p=\frac{2\sigma}{R}\)\\[0.8em]

毛细高度 &
\(\displaystyle h=\frac{2\sigma\cos\theta}{\rho gr}\)\\[0.8em]

克拉珀龙方程 &
\(\displaystyle \frac{\dd p}{\dd T}=\frac{l}{T(v_2-v_1)}\)\\
\end{tabularx}
}

\section{题型总路线}

\subsection{题型一：热力学过程计算}

\begin{MethodBox}[核心路线]
这类题通常给出理想气体的初末状态和过程条件，要求 \(Q,W,\Delta U\) 或效率。
\begin{ReviewList}
  \item 先判断过程。
  \item 用 \(pV=\nu RT\) 找缺失状态量。
  \item 用 \(W=-\int p\dd V\) 算功。
  \item 用 \(\Delta U=\nu C_{V,m}\Delta T\) 算内能变化。
  \item 用 \(Q=\Delta U-W\) 算热量。
  \item 循环题用 \(\Delta U=0\)，热机题用 \(\eta=1-Q_2/Q_1\)。
\end{ReviewList}
\end{MethodBox}

\subsection{题型二：熵变与自发方向}

\begin{MethodBox}[核心路线]
这类题常考自由膨胀、有限温差传热、热机可逆性判断。
\begin{ReviewList}
  \item 熵变只看初末态。
  \item 可逆过程直接算：
  \[
    \Delta S=\int\frac{\dbar Q_{\mathrm{rev}}}{T}.
  \]
  \item 不可逆过程要构造同初末态的可逆路径。
  \item 判断自发方向看总熵：
  \[
    \Delta S_{\mathrm{total}}\ge0.
  \]
\end{ReviewList}
\end{MethodBox}

\subsection{题型三：麦克斯韦分布与速率}

\begin{MethodBox}[核心路线]
\begin{ReviewList}
  \item 若问速率，先判断是 \(v_{\mathrm p}\)、\(\bar v\)、还是 \(v_{\mathrm{rms}}\)。
  \item 若问曲线变化，看温度 \(T\) 与摩尔质量 \(\mu\)。
  \item 若问压强或温度的微观解释，用
  \[
    p=\frac13nm\overline{v^2},
    \qquad
    \overline{\varepsilon_k}=\frac32k_{\mathrm B}T.
  \]
  \item 若问外力场分布，用
  \[
    n\propto e^{-\varepsilon_p/k_{\mathrm B}T}.
  \]
\end{ReviewList}
\end{MethodBox}

\subsection{题型四：输运过程}

\begin{MethodBox}[核心路线]
\begin{ReviewList}
  \item 温度梯度：热传导。
  \item 流速梯度：黏滞。
  \item 数密度梯度：扩散。
  \item 三个宏观规律都按
  \[
    \text{流密度}=-\text{系数}\times\text{梯度}
  \]
  来记。
  \item 微观输运系数看 \(\lambda,\bar v,\rho,c_V\)。
\end{ReviewList}
\end{MethodBox}

\subsection{题型五：表面张力与相变}

\begin{MethodBox}[核心路线]
\begin{ReviewList}
  \item 表面积变化题：用 \(\Delta A=\sigma\Delta S\)，注意肥皂膜两个表面。
  \item 弯曲液面题：液滴 \(2\sigma/R\)，肥皂泡 \(4\sigma/R\)。
  \item 毛细题：用
  \[
    h=\frac{2\sigma\cos\theta}{\rho gr}.
  \]
  \item 相变题：看潜热、体积变化、相平衡曲线。
  \item 相图斜率题：用
  \[
    \frac{\dd p}{\dd T}=\frac{l}{T(v_2-v_1)}.
  \]
\end{ReviewList}
\end{MethodBox}

\section{考前最后一页}

\begin{WarnBox}[全书最高频错误]
\begin{PitfallList}
  \item 不看符号约定，直接套功的公式。本资料中 \(W\) 是外界对系统做功。
  \item 把 \(Q,W\) 当作态函数。
  \item 非准静态过程乱用 \(W=-\int p\dd V\)。
  \item 把绝热过程直接等同于等熵过程。
  \item 计算不可逆过程熵变时，沿不可逆路径积分 \(\dbar Q/T\)。
  \item 卡诺公式中使用摄氏温度。
  \item 混淆 \(v_{\mathrm p},\bar v,v_{\mathrm{rms}}\)。
  \item 用 \(\bar v^2\) 代替 \(\overline{v^2}\)。
  \item 玻尔兹曼分布指数符号写反。
  \item 输运公式负号方向理解错误。
  \item 忘记肥皂泡有两个表面。
  \item 克拉珀龙方程中潜热和体积单位不匹配。
\end{PitfallList}
\end{WarnBox}

\begin{KeyBox}[考前检查顺序]
\begin{ReviewList}
  \item 先把第一、第二、第三章的热力学主线背熟。这是分数最高、题型最稳定的部分。
  \item 再背第四章的三种速率、麦克斯韦分布和能量均分。
  \item 再背第五章三个输运定律和三个输运系数。
  \item 最后看第六章表面张力、毛细现象和克拉珀龙方程。
\end{ReviewList}
\end{KeyBox}

\finalnote{
热学B的复习不要追求推导越多越好。最实用的目标是：看到题目能立即判断系统、过程、适用公式和符号约定。只要这四件事稳定，绝大多数计算题和判断题都能顺着做下去。
}

\appendix

\chapter{符号与约定}

\begin{FormulaBox}[常用符号]
\begin{tabularx}{\textwidth}{@{}>{\bfseries}lX@{}}
\(p,V,T\) & 压强、体积、热力学温度。\\
\(U,H,S\) & 内能、焓、熵。\\
\(C_V,C_p\) & 定容热容、定压热容；摩尔量常写作 \(C_{V,m},C_{p,m}\)。\\
\(k,R,N_A\) & 玻尔兹曼常数、气体常数、阿伏伽德罗常数，满足 \(R=N_A k\)。\\
\(\alpha,\beta\) & 体膨胀系数、等温压缩系数。\\
\(\gamma\) & 热容比，\(\gamma=C_p/C_V\)。\\
\end{tabularx}
\end{FormulaBox}

\backmatter

\chapter*{版本记录}
\addcontentsline{toc}{chapter}{版本记录}

\begin{SourceBox}[版本说明]
本资料为中国科学技术大学《热学B》总复习资料，版本号为 \textbf{2026.1}，由 \textbf{\DocAuthor} 整理。资料内容以 2026 春季学期授课 PDF 与《热学B》知识要点总结为基础，定位为考前系统复习材料。整理重点包括：知识框架、核心公式、适用条件、典型题型、易错点与考前速查。本资料按《热学B》课程范围编写。
\end{SourceBox}

\begin{longtable}{@{}p{0.22\textwidth}p{0.68\textwidth}@{}}
\toprule
\textbf{版本} & \textbf{内容说明}\\
\midrule
\endhead

\textbf{Ver. 2026.1} &
完成《热学B》总复习资料初版。资料采用 \texttt{main.tex}、\texttt{preamble.tex}、\texttt{macros.tex} 与 \texttt{chapters/} 文件夹结构编写，正文覆盖六章课程内容，并附有总复习框架。\\[0.6em]

&
第一章“热力学基础知识与温度”：整理热力学系统、状态参量、平衡态、准静态过程、热力学第零定律、温度、温标、理想气体状态方程、范德瓦耳斯方程以及简单固液体状态方程。\\[0.6em]

&
第二章“热力学第一定律”：整理功和热量、内能、焓、热容、热力学第一定律、理想气体典型过程、多方过程、绝热节流、焦耳--汤姆孙效应、热机效率与制冷系数。\\[0.6em]

&
第三章“热力学第二定律与熵”：整理可逆与不可逆过程、第二定律的开尔文表述和克劳修斯表述、卡诺定理、热力学温标、克劳修斯等式与不等式、熵、熵增加原理、理想气体熵变与 \(T\)--\(S\) 图。\\[0.6em]

&
第四章“麦克斯韦--玻尔兹曼分布律”：整理理想气体压强公式、温度的统计解释、麦克斯韦速度分布与速率分布、最概然速率、平均速率、方均根速率、泻流、玻尔兹曼分布、能量均分定理与理想气体热容量。\\[0.6em]

&
第五章“气体中的输运过程”：整理热传导、黏滞、扩散三种输运现象，傅里叶定律、牛顿黏性定律、菲克定律，平均自由程、碰撞频率、输运系数的微观表达，以及输运系数与压强的关系。\\[0.6em]

&
第六章“固、液体性质简介与相变”：整理晶体与非晶体、单晶体与多晶体、晶体热学性质、液体结构、表面张力、弯曲液面附加压强、毛细现象、相变、相图、临界点、克拉珀龙方程与克劳修斯--克拉珀龙近似。\\[0.6em]

&
总结部分：整理全课程知识地图、核心概念对比、热力学过程计算路线、核心公式总表、题型总路线与考前高频易错点。\\

\bottomrule
\end{longtable}

\begin{KeyBox}[后续修订方向]
后续版本可根据课程作业、习题课、TA 重点和往年题继续补充典型例题、题型索引和公式适用条件索引；若最终页数超过目标，可优先压缩概念性说明，保留公式、题型模板和易错点。
\end{KeyBox}

\end{document}